\documentclass[10pt,letterpaper,oneside]{scrreprt}
\usepackage[letterpaper,top=19mm,bottom=19mm,left=23mm,right=23mm,headsep=5.5mm,footskip=9mm]{geometry}
\usepackage{fontspec}
\usepackage{xeCJK}
\setmainfont{Palatino}
\setsansfont{Avenir Next}
\setmonofont[Scale=0.82]{Menlo}
\setCJKmainfont{Songti SC}
\setCJKsansfont{Hiragino Sans GB}
\setCJKmonofont{Songti SC}
\newfontfamily\symbolfont{Arial Unicode MS}
\usepackage{graphicx}
\usepackage[table]{xcolor}
\usepackage{booktabs,longtable,array,tabularx,colortbl}
\usepackage{calc}
\usepackage{enumitem}
\usepackage{fancyhdr}
\usepackage{microtype}
\usepackage{hyperref}
\usepackage{bookmark}
\usepackage{ragged2e}
\usepackage{setspace}
\usepackage{needspace}
\usepackage{etoolbox}
\usepackage[most]{tcolorbox}

\definecolor{GewuBlue}{HTML}{235EA7}
\definecolor{GewuInk}{HTML}{172033}
\definecolor{GewuMuted}{HTML}{68758B}
\definecolor{GewuPale}{HTML}{EDF4FC}
\definecolor{GewuRule}{HTML}{D7E1EC}
\definecolor{GewuTableHeader}{HTML}{E3EEF9}
\definecolor{GewuTableStripe}{HTML}{F5F8FC}
\definecolor{GewuSealPanel}{HTML}{526985}

\hypersetup{colorlinks=true,linkcolor=GewuBlue,urlcolor=GewuBlue,citecolor=GewuBlue,pdfauthor={Lingyuan Kong and Zhengran Wu},pdftitle={Gewu Insight User Guide}}
\urlstyle{same}
\Urlmuskip=0mu plus 1mu
\setlength{\parindent}{0pt}
\setlength{\parskip}{3pt plus .8pt minus .5pt}
\setstretch{1.01}
\setlist{topsep=2.5pt,itemsep=1.25pt,parsep=.5pt,leftmargin=5.5mm}
% Number the reader-facing hierarchy, but stop before implementation-detail
% subsections.  Each part owns its chapter sequence, so the visible chapter
% labels are 1-1, 1-2, 2-1, 3-1, ... instead of one global counter.
\setcounter{secnumdepth}{1}
\setcounter{tocdepth}{1}
\makeatletter
\@addtoreset{chapter}{part}
\makeatother
\renewcommand{\thechapter}{\arabic{part}-\arabic{chapter}}
\renewcommand{\thesection}{\thechapter.\arabic{section}}
\renewcommand{\thesubsection}{\thesection.\arabic{subsection}}
\renewcommand{\arraystretch}{1.16}
\setlength{\tabcolsep}{3.5pt}
\setlength{\LTpre}{4pt}
\setlength{\LTpost}{5pt}
\setlength{\heavyrulewidth}{0.55pt}
\setlength{\lightrulewidth}{0.3pt}
\arrayrulecolor{GewuRule}
\setlength{\emergencystretch}{3em}
\hyphenpenalty=10000
\exhyphenpenalty=10000
\widowpenalty=10000
\clubpenalty=10000
\setlength{\headheight}{16pt}
\providecommand{\tightlist}{\setlength{\itemsep}{1pt}\setlength{\parskip}{0pt}}
\providecommand{\pandocbounded}[1]{#1}

\setkomafont{part}{\sffamily\bfseries\color{GewuInk}\fontsize{25}{29}\selectfont}
\setkomafont{partnumber}{\sffamily\bfseries\color{GewuBlue}\fontsize{11}{13}\selectfont}
\setkomafont{chapter}{\sffamily\bfseries\color{GewuInk}\fontsize{21}{24}\selectfont}
\setkomafont{section}{\sffamily\bfseries\color{GewuInk}\fontsize{15.5}{18.5}\selectfont}
\setkomafont{subsection}{\sffamily\bfseries\color{GewuInk}\fontsize{11.8}{14}\selectfont}
\setkomafont{subsubsection}{\sffamily\bfseries\color{GewuInk}\fontsize{10.5}{12.5}\selectfont}
\renewcommand*{\partformat}{\partname~\thepart\autodot\enskip}
\RedeclareSectionCommand[beforeskip=2.1ex plus .4ex minus .2ex,afterskip=.85ex,tocnumwidth=3.8em]{chapter}
\RedeclareSectionCommand[beforeskip=1.8ex plus .3ex minus .2ex,afterskip=.48ex,tocnumwidth=5.2em]{section}
\RedeclareSectionCommand[beforeskip=1.35ex plus .25ex minus .15ex,afterskip=.3ex]{subsection}
\renewcommand*{\chapterformat}{\color{GewuBlue}\rule[-2pt]{1.7mm}{23pt}\hspace{3mm}\thechapter\enskip}
\renewcommand*{\sectionformat}{\thesection\enskip}
\pretocmd{\section}{\Needspace{9\baselineskip}}{}{}
\pretocmd{\subsection}{\Needspace{6\baselineskip}}{}{}
\pretocmd{\subsubsection}{\Needspace{3\baselineskip}}{}{}
% Product Identity is one continuous reading sequence. Suppress only the
% chapter's forced break locally; normal page overflow, numbering, TOC entries
% and labels are still handled by KOMA-Script. Other parts remain unchanged.
\newcommand{\GewuIdentityChapter}[1]{%
  \par\Needspace{9\baselineskip}%
  \let\GewuSavedClearPage\clearpage
  \let\GewuSavedClearDoublePage\cleardoublepage
  \let\clearpage\relax
  \let\cleardoublepage\relax
  \chapter{#1}%
  \let\clearpage\GewuSavedClearPage
  \let\cleardoublepage\GewuSavedClearDoublePage
}
\raggedbottom
\newenvironment{GewuRelated}{\par\nobreak\vspace{3pt}\nobreak\begingroup\small\sffamily\begin{samepage}\noindent}{\par\end{samepage}\endgroup}

\pagestyle{fancy}
\fancyhf{}
\fancyhead[L]{\footnotesize\sffamily\bfseries\color{GewuInk}Gewu Insight User Guide}
\fancyhead[R]{\footnotesize\sffamily\color{GewuMuted}\leftmark}
\fancyfoot[L]{\footnotesize\sffamily\color{GewuMuted}Gewu Insight}
\fancyfoot[R]{\footnotesize\sffamily\color{GewuMuted}\thepage}
\renewcommand{\headrulewidth}{0.35pt}
\renewcommand{\headrule}{\hbox to\headwidth{\color{GewuRule}\leaders\hrule height \headrulewidth\hfill}}

\newenvironment{GewuLead}
  {\begin{list}{}{\leftmargin=0pt\rightmargin=0pt\topsep=0pt\partopsep=0pt\parsep=0pt}\item\relax
    \sffamily\fontsize{10.5}{13.2}\selectfont\color{GewuMuted}}
  {\end{list}}

\newenvironment{GewuResult}
  {\begin{tcolorbox}[enhanced,breakable,colback=GewuPale,colframe=GewuPale,boxrule=0pt,
    borderline west={1.6pt}{0pt}{GewuBlue},arc=1.2mm,left=2.5mm,right=2.5mm,top=1.7mm,bottom=1.7mm,
    before skip=4pt,after skip=5pt,fontupper=\color{GewuInk}\small\sffamily]}
  {\end{tcolorbox}}

\newenvironment{GewuCaption}
  {\begin{list}{}{\leftmargin=2.5mm\rightmargin=2.5mm\topsep=0pt}\item\relax
    \footnotesize\itshape\color{GewuMuted}}
  {\end{list}}

\renewenvironment{quote}
  {\begin{tcolorbox}[enhanced,breakable,colback=GewuPale,colframe=GewuPale,boxrule=0pt,
    borderline west={1.6pt}{0pt}{GewuBlue},arc=1.2mm,left=2.5mm,right=2.5mm,top=1.7mm,bottom=1.7mm,
    before skip=4pt,after skip=5pt,fontupper=\color{GewuInk}\small]}
  {\end{tcolorbox}}

\newcommand{\GewuTableStart}{\begingroup\fontsize{8.6}{10.6}\selectfont\rowcolors{2}{GewuTableStripe}{white}}
\newcommand{\GewuTableEnd}{\endgroup}

\begin{document}
\begin{titlepage}
  \centering
  \vspace*{25mm}
  \includegraphics[width=36mm]{/Users/lingyuan/PaperMonitor-20260516/assets/gewu-insight-app-icon.png}\\[8mm]
  {\sffamily\bfseries\fontsize{36}{40}\selectfont\color{GewuInk} Gewu Insight}\\[4mm]
  \includegraphics[width=57mm]{/Users/lingyuan/PaperMonitor-20260516/assets/branding/GewuAboutCalligraphy.png}\\[5mm]
  {\sffamily\color{GewuMuted}\fontsize{15}{18}\selectfont User Guide}\\[2.5mm]
  {\sffamily\small\color{GewuMuted} Version 0.1 (Build 673606679)}\\[19mm]
  \begin{minipage}{0.72\textwidth}
    \centering\sffamily\fontsize{13}{18}\selectfont\color{GewuInk}
    Your personal research workspace, keeping papers, evidence, and files connected over time.
  \end{minipage}
  \vfill
  {\color{GewuRule}\rule{68mm}{0.45pt}}\\[7mm]
  {\small\sffamily\color{GewuMuted}Designed and developed by Lingyuan Kong and Zhengran Wu,\\with coding assistance from OpenAI Codex.}\\[4mm]
  {\footnotesize\sffamily\color{GewuMuted}Copyright \textcopyright\ 2026 Lingyuan Kong and Zhengran Wu. All Rights Reserved.}\\[10mm]
\end{titlepage}

\chapter*{About this manual}
\addcontentsline{toc}{chapter}{About this manual}
This PDF is the fixed-layout edition of the in-app Gewu Insight User Guide. It contains the same product guidance as the searchable Help window, arranged for comfortable reading, printing, and archival use. Interface names remain in English so that the manual matches the application.

\section*{Three complementary reading paths}
\begin{description}[style=nextline,leftmargin=0pt,labelindent=0pt,itemsep=3pt]
  \item[Operate the workspace] Parts I and III explain the interface and everyday workflows: navigation, papers and files, importing, monitors, settings, recovery, and transfer.
  \item[Understand the product] Parts II and IV explain the local-first philosophy, the visible-workspace model, SQLite's role as the source of truth, and the product's visual identity.
  \item[Inspect the technical contract] Part V documents storage, database relationships, network behavior, migration, troubleshooting, and other details needed for maintenance or development.
\end{description}

These paths intentionally share one manual. The user-facing behavior, the product model, and the implementation constraints can therefore be read independently without drifting into separate and contradictory documents.

\clearpage
\section*{Contents}
% Keep the complete navigation list compact without changing body typography.
\begingroup
\setstretch{0.96}
\setlength{\parskip}{0pt}
\makeatletter
\@starttoc{toc}
\makeatother
\endgroup
\clearpage
\RaggedRight
\part{Orientation}\label{part-1}

\chapter{Welcome}\label{welcome}

\begin{GewuLead}
Gewu Insight is a local workspace for collecting papers, preserving their files and versions, and following journals, subjects, authors, and your own publications.
\end{GewuLead}

Guide revision · September 9, 2026 · Illustrated windows and settings

\section{Why this workspace exists}\label{why-this-workspace-exists}

Research rarely arrives as one clean collection. A paper may first appear in a feed, later acquire an arXiv version and a DOI, be downloaded under several filenames, reappear through an author\textquotesingle s profile, and eventually become part of a project or your own publication record. Publisher pages, alerts, PDFs, folders, notes, and citation services each preserve only part of that history.

Gewu Insight is designed to keep those fragments connected. It treats discovery, identity, files, versions, organization, reading state, and provenance as parts of one continuing research process. The aim is not to accumulate the largest possible list. The aim is to preserve enough context that a useful paper remains understandable and usable months or years after it was first found.

\section{Why the name 格物洞见}\label{why-the-name-ux683cux7269ux6d1eux89c1}

The name recalls the Chinese idea of \textbf{格物致知}: approaching things closely enough to understand their underlying order. In this practical research setting, \emph{gewu}---the investigation of things---means continuously examining the research object, its sources, authors, versions, files, relationships, and changing state.

{\def\LTcaptype{none} % do not increment counter
\GewuTableStart
\begin{longtable}[]{@{}>{\RaggedRight\arraybackslash}p{0.27\textwidth}>{\RaggedRight\arraybackslash}p{0.63\textwidth}@{}}
\toprule\noalign{}
\rowcolor{GewuTableHeader}
Element & Role \\
\midrule\noalign{}
\endhead
\bottomrule\noalign{}
\endlastfoot
{\textbf{格物 · Gewu}} & {Stay close to the research object. Keep the paper, identifiers, versions, sources, files, and relationships inspectable instead of reducing them to an untraceable card or summary.} \\
{\textbf{洞见 · Insight}} & {Let understanding emerge from connected evidence, repeated reading, comparison, and the researcher\textquotesingle s judgment. Insight is accumulated; it is not a conclusion the app claims to manufacture.} \\
\end{longtable}
\GewuTableEnd
}

\begin{GewuCaption}
The Chinese name is part of the product idea, not a decorative translation of the English name.
\end{GewuCaption}

\begin{quote}
Keep research objects visible, evidence traceable, and insight cumulative.{让研究对象保持可见，让证据保持可追，让洞见能够积累。}
\end{quote}

\section{Five design commitments}\label{five-design-commitments}

{\def\LTcaptype{none} % do not increment counter
\GewuTableStart
\begin{longtable}[]{@{}>{\RaggedRight\arraybackslash}p{0.27\textwidth}>{\RaggedRight\arraybackslash}p{0.63\textwidth}@{}}
\toprule\noalign{}
\rowcolor{GewuTableHeader}
Principle & What it means \\
\midrule\noalign{}
\endhead
\bottomrule\noalign{}
\endlastfoot
\textbf{Local ownership} & {Your live database and managed research files remain in the local project. Online services contribute evidence; they do not own the workspace.} \\
\textbf{Visible evidence} & {PDFs and attachments remain real files that Finder, Preview, and other applications can inspect.} \\
\textbf{Continuity over time} & {A preprint, later DOI, publication version, annotations, and personal contexts can accumulate around one persistent paper object.} \\
\textbf{Explicit state and provenance} & {Completed, empty, partial, failed, downloaded, imported, and retried are different states. Sources and operational outcomes should remain visible.} \\
\textbf{Researcher control} & {Automation assists discovery and maintenance, but classification, interpretation, organization, and scientific judgment remain yours.} \\
\end{longtable}
\GewuTableEnd
}

\begin{quote}
\textbf{Local-first does not mean offline-only}

The durable workspace is local, while supported discovery, metadata, profile, and download features may require network access and remain subject to source, publisher, and institutional availability. Verify uncertain matches and versions; the app does not bypass access controls.
\end{quote}

\section{Cross-platform core, platform-specific integration}\label{cross-platform-core-platform-specific-integration}

Gewu Insight separates its portable Python, SQLite, and web-interface core from the native desktop layer. Paper identity, database migrations, monitor state, resource relationships, and most interface logic are designed without making macOS the definition of the product.

{\def\LTcaptype{none} % do not increment counter
\GewuTableStart
\begin{longtable}[]{@{}>{\RaggedRight\arraybackslash}p{0.27\textwidth}>{\RaggedRight\arraybackslash}p{0.63\textwidth}@{}}
\toprule\noalign{}
\rowcolor{GewuTableHeader}
Platform layer & Current status \\
\midrule\noalign{}
\endhead
\bottomrule\noalign{}
\endlastfoot
Portable core & Python, SQLite, the local application service, and the web-based workspace form the shared architectural foundation. \\
macOS native release & This is the current primary release and the platform on which the complete desktop workflow is developed and tested, including Finder, Preview, Quick Look, Dock, Trash, and native drag and drop. \\
Planned Windows release & Windows support is planned to begin as an experimental build. Explorer, the default PDF viewer, Recycle Bin, native drag and drop, packaging, and updates require their own implementation and validation. \\
\end{longtable}
\GewuTableEnd
}

\begin{quote}
\textbf{Current distribution status}

This guide does not claim that a supported Windows package is available today. Instructions that name Finder, Preview, Quick Look, the Dock, or macOS Trash describe the currently tested macOS integration. Equivalent Windows instructions should be added when the experimental Windows package exists and has been verified.
\end{quote}

\section{What this guide covers}\label{what-this-guide-covers}

This guide follows the work you do in the app: put papers into the Library, organize them, manage publication versions and attachments, review new literature, and keep the database and files safe. It describes the current interface rather than the history of the project.

\section{What Gewu Insight is---and is not}\label{what-gewu-insight-isand-is-not}

{\def\LTcaptype{none} % do not increment counter
\GewuTableStart
\begin{longtable}[]{@{}>{\RaggedRight\arraybackslash}p{0.27\textwidth}>{\RaggedRight\arraybackslash}p{0.63\textwidth}@{}}
\toprule\noalign{}
\rowcolor{GewuTableHeader}
It is & It is not only \\
\midrule\noalign{}
\endhead
\bottomrule\noalign{}
\endlastfoot
A local research workspace that connects discovery, identity, files, versions, organization, and durable state. & An RSS reader that forgets a result after it leaves the feed. \\
A visible file and resource system built around persistent paper identities. & A PDF folder whose filenames must carry every piece of meaning. \\
A personal research memory shared by Library, Spaces, publications, and monitors. & Five unrelated tools that happen to occupy the same window. \\
A system that records evidence and assists repeatable work. & A black-box ranking or summary engine that replaces scientific judgment. \\
\end{longtable}
\GewuTableEnd
}

\section{Six main tabs}\label{six-main-tabs}

\begin{itemize}
\tightlist
\item
  \textbf{Library} --- your managed papers, folders, spaces, versions, and related files.
\item
  \textbf{My Publications} --- your publication list, citation metrics, author roles, and export tools.
\item
  \textbf{Journal Monitor} --- papers published on a selected date, organized by journal and field.
\item
  \textbf{Subject Monitor} --- topic-based discovery from supported online sources and feeds.
\item
  \textbf{Author Monitor} --- monitored researchers, profile-linked discovery, and newly found work.
\item
  \textbf{Browser} --- Google search, websites, online PDFs, hierarchical bookmarks, and local browsing history.
\end{itemize}

\section{One continuous research loop}\label{one-continuous-research-loop}

{\def\LTcaptype{none} % do not increment counter
\GewuTableStart
\begin{longtable}[]{@{}
  >{\raggedright\arraybackslash}p{(\linewidth - 2\tabcolsep) * \real{0.5000}}
  >{\raggedright\arraybackslash}p{(\linewidth - 2\tabcolsep) * \real{0.5000}}@{}}
\toprule\noalign{}
\rowcolor{GewuTableHeader}
\begin{minipage}[b]{\linewidth}\raggedright
Element
\end{minipage} & \begin{minipage}[b]{\linewidth}\raggedright
Role
\end{minipage} \\
\midrule\noalign{}
\endhead
\bottomrule\noalign{}
\endlastfoot
\begin{minipage}[t]{\linewidth}\raggedright
{{Acquire}\\
\textbf{Find}}\strut
\end{minipage} & {Journal, Subject, and Author Monitor introduce relevant candidates.} \\
\begin{minipage}[t]{\linewidth}\raggedright
{{Identify}\\
\textbf{Recognize}}\strut
\end{minipage} & {DOI, arXiv, title, authors, and source evidence are connected to a shared record.} \\
\begin{minipage}[t]{\linewidth}\raggedright
{{Resolve}\\
\textbf{Preserve}}\strut
\end{minipage} & {Metadata, files, versions, roles, paths, and provenance are made explicit.} \\
\begin{minipage}[t]{\linewidth}\raggedright
{{Remember}\\
\textbf{Understand}}\strut
\end{minipage} & {Library, Spaces, My Publications, reading state, and personal organization retain meaning.} \\
\end{longtable}
\GewuTableEnd
}

\begin{GewuCaption}
What you save, read, classify, and follow becomes context for the next cycle. The five research views share this loop; Browser supplies a separate place to explore websites.
\end{GewuCaption}

\begin{quote}
\textbf{Important}

SQLite is the source of truth for Gewu Insight. Do not replace, rebuild, or edit the database by hand. Use the app to move, classify, or delete managed items.
\end{quote}

\section{Choose a path through this guide}\label{choose-a-path-through-this-guide}

This is one comprehensive manual with three complementary reading paths. The \textbf{operational path} explains how to use the app; the \textbf{concept and architecture path} explains why the workspace is organized this way; and the \textbf{technical reference path} documents storage, database, network, migration, and troubleshooting details. They remain together so visible behavior, product principles, and implementation constraints do not drift apart.

{\def\LTcaptype{none} % do not increment counter
\GewuTableStart
\begin{longtable}[]{@{}>{\RaggedRight\arraybackslash}p{0.27\textwidth}>{\RaggedRight\arraybackslash}p{0.63\textwidth}@{}}
\toprule\noalign{}
\rowcolor{GewuTableHeader}
If you want to & Continue with \\
\midrule\noalign{}
\endhead
\bottomrule\noalign{}
\endlastfoot
Start using the app now & Read \textbf{Part I · Orientation}, then use \textbf{Part III · User Manual} as the operational guide. \\
Understand the product idea and architecture & Read \textbf{Part II · Product Concepts and Architecture}, followed by \textbf{Part IV · Product Identity}. \\
Maintain, migrate, or inspect the system & Use \textbf{Part V · Technical and Operational Reference}, with Part II for the underlying model. \\
Look up one operation & Use the topic list or press \texttt{{\symbolfont ⌘}}\texttt{F} to search this guide. \\
Prepare an upgrade or move to another Mac & Read \hyperref[backup-migrate]{Back Up or Move to Another Mac} before changing the installation. \\
\end{longtable}
\GewuTableEnd
}

\section{Open this guide again}\label{open-this-guide-again}

Choose \textbf{Help \textgreater{} User Guide}. The guide opens in its own window and works without an internet connection. Search filters the topic list; choose a result to read it, or use Previous and Next to follow the manual. Select \textbf{Print Guide} to open the complete, fixed-layout PDF in macOS Preview, where you can print or save a copy. In a separately hosted web UI, that button opens the PDF in a browser tab. It opens the whole guide, not just the topic currently on screen.

\begin{GewuRelated}
\textbf{See also:} \hyperref[guided-tour]{Guided Tour} · \hyperref[first-session]{Your First Session} · \hyperref[system-overview]{System Architecture at a Glance} · \hyperref[storage-model]{How Gewu Insight Stores Files}
\end{GewuRelated}

\chapter{Guided Tour}\label{guided-tour}

\begin{GewuLead}
The top chooses a work mode; the panes move from context on the left to the active object and its evidence; the bottom preserves navigation and progress without reshaping the workspace.
\end{GewuLead}

\section{Read the interface as a sentence}\label{read-the-interface-as-a-sentence}

Gewu Insight uses a consistent left-to-right grammar. The exact number of panes changes with the task, but their meaning remains stable: choose a research context, narrow it to an object, inspect its evidence, then act. This lets a dense workspace remain predictable without hiding the underlying files and relationships.

{\def\LTcaptype{none} % do not increment counter
\GewuTableStart
\begin{longtable}[]{@{}
  >{\raggedright\arraybackslash}p{(\linewidth - 2\tabcolsep) * \real{0.5000}}
  >{\raggedright\arraybackslash}p{(\linewidth - 2\tabcolsep) * \real{0.5000}}@{}}
\toprule\noalign{}
\rowcolor{GewuTableHeader}
\begin{minipage}[b]{\linewidth}\raggedright
Element
\end{minipage} & \begin{minipage}[b]{\linewidth}\raggedright
Role
\end{minipage} \\
\midrule\noalign{}
\endhead
\bottomrule\noalign{}
\endlastfoot
\begin{minipage}[t]{\linewidth}\raggedright
{{Top}\\
\textbf{Choose a mode}}\strut
\end{minipage} & {Library, My Publications, one of the three discovery monitors, or Browser.} \\
\begin{minipage}[t]{\linewidth}\raggedright
{{Left}\\
\textbf{Set the context}}\strut
\end{minipage} & {A root, folder, date, journal, subject, author, or filter defines the working scope.} \\
\begin{minipage}[t]{\linewidth}\raggedright
{{Center}\\
\textbf{Select the object}}\strut
\end{minipage} & {Browse folders, paper lists, search results, or the resources belonging to one paper.} \\
\begin{minipage}[t]{\linewidth}\raggedright
{{Right}\\
\textbf{Inspect and act}}\strut
\end{minipage} & {Read metadata and evidence; open, download, classify, refresh, export, or reorganize.} \\
\end{longtable}
\GewuTableEnd
}

\begin{GewuCaption}
The bottom edge is reserved for progress, breadcrumbs, and view controls so background work does not repeatedly resize the main content.
\end{GewuCaption}

\section{The Library window}\label{the-library-window}

{\def\LTcaptype{none} % do not increment counter
\GewuTableStart
\begin{longtable}[]{@{}>{\RaggedRight\arraybackslash}p{0.27\textwidth}>{\RaggedRight\arraybackslash}p{0.63\textwidth}@{}}
\toprule\noalign{}
\rowcolor{GewuTableHeader}
Area & Purpose \\
\midrule\noalign{}
\endhead
\bottomrule\noalign{}
\endlastfoot
\textbf{1. Source tabs} & {Library, shortcuts, and independent Spaces.} \\
\textbf{2. Folder columns} & {Move from broad folders to the selected item.} \\
\textbf{3. Item column} & {Papers, files, and child folders in the current location.} \\
\textbf{4. Paper Files} & {One flat list of resources with independent type, annotation, source, and Default status.} \\
\textbf{5. Details} & {Metadata, identifiers, actions, citations, and folder membership.} \\
\end{longtable}
\GewuTableEnd
}

\section{Follow selection from broad to specific}\label{follow-selection-from-broad-to-specific}

A selected item does not always mean ``open this file.'' In a folder column, selection establishes the next level of context. In the item column, it identifies the active paper or standalone file. In Paper Files, it identifies one resource inside the active paper bundle. Paper Details then describes the shared work rather than merely repeating the selected filename.

{\def\LTcaptype{none} % do not increment counter
\GewuTableStart
\begin{longtable}[]{@{}>{\RaggedRight\arraybackslash}p{0.19\textwidth}>{\RaggedRight\arraybackslash}p{0.29\textwidth}>{\RaggedRight\arraybackslash}p{0.36\textwidth}@{}}
\toprule\noalign{}
\rowcolor{GewuTableHeader}
Selection & What becomes active & What remains shared \\
\midrule\noalign{}
\endhead
\bottomrule\noalign{}
\endlastfoot
A Library folder & Its child folders and items. & The papers keep their identities, resources, and other placements. \\
A paper in the item column & Its Paper Files bundle and Paper Details. & The same work can still belong to other folders, monitors, or My Publications. \\
A resource in Paper Files & That version or attachment for preview, classification, or file action. & The parent paper remains the active scholarly object. \\
A Monitor result & Its discovery evidence, report, and available actions. & Adding it to Library reuses the cataloged paper record when identity has been resolved. \\
\end{longtable}
\GewuTableEnd
}

\section{Tabs across the top}\label{tabs-across-the-top}

Click a top tab to change workspaces and focus that tab. Left/Right then switch tabs; Down enters the page. Monitors show calendars, filters and papers. Browser has a Bookmarks/History sidebar beside a native webpage. Progress messages stay in reserved status areas.

The compact, flat search control at the top right expands on hover or focus. Each research workspace retains its own query; clear it with ×. On macOS, double-click only blank title-bar space to zoom or restore the window. Tabs, search and other controls do not resize it.

{\def\LTcaptype{none} % do not increment counter
\GewuTableStart
\begin{longtable}[]{@{}>{\RaggedRight\arraybackslash}p{0.19\textwidth}>{\RaggedRight\arraybackslash}p{0.29\textwidth}>{\RaggedRight\arraybackslash}p{0.36\textwidth}@{}}
\toprule\noalign{}
\rowcolor{GewuTableHeader}
Area & Primary question & Typical next action \\
\midrule\noalign{}
\endhead
\bottomrule\noalign{}
\endlastfoot
\textbf{Library} & What have I chosen to keep, and how is it organized? & Read, classify files, move, copy, attach, or recover. \\
\textbf{My Publications} & Which works are mine, and what authorship or citation context do they carry? & Verify roles and metrics, refresh, order, or export. \\
\textbf{Journal Monitor} & What appeared in selected journals on a given date? & Filter, inspect, download, or add to Library. \\
\textbf{Subject Monitor} & What matches a defined scientific topic across supported sources? & Review source evidence and preserve useful results. \\
\textbf{Author Monitor} & What is new in the work of a researcher I follow? & Inspect the author relationship, open the work, or add a Library placement. \\
\textbf{Browser} & What can I find on the web without leaving this workspace? & Search Google, read a website or online PDF, save a bookmark, or revisit history. \\
\end{longtable}
\GewuTableEnd
}

\section{Watch one paper move through the workspace}\label{watch-one-paper-move-through-the-workspace}

{\def\LTcaptype{none} % do not increment counter
\GewuTableStart
\begin{longtable}[]{@{}>{\RaggedRight\arraybackslash}p{0.27\textwidth}>{\RaggedRight\arraybackslash}p{0.63\textwidth}@{}}
\toprule\noalign{}
\rowcolor{GewuTableHeader}
Element & Role \\
\midrule\noalign{}
\endhead
\bottomrule\noalign{}
\endlastfoot
{\textbf{Discovered}} & {A Monitor records where and when the work was observed.} \\
{\textbf{Cataloged}} & {Identity evidence connects the observation to a persistent paper record.} \\
{\textbf{Organized}} & {A Library folder, Space, or My Publications adds personal context.} \\
{\textbf{File-backed}} & {Downloaded or attached resources become available in Paper Files and Finder.} \\
\end{longtable}
\GewuTableEnd
}

\begin{GewuCaption}
These states can overlap. A cataloged Monitor result may have no local file, and a downloaded file does not have to belong to a Library folder immediately.
\end{GewuCaption}

\section{Resize a pane}\label{resize-a-pane}

\begin{enumerate}
\tightlist
\item
  \textbf{Point to a divider.} Move the pointer to the vertical line between two panes.
\item
  \textbf{Drag the divider.} Move it left or right until the useful content fits.
\item
  \textbf{Release.} The app uses the new working width without changing the underlying files.
\end{enumerate}

\begin{quote}
\textbf{Note}

The right-side Paper Details area scrolls independently. A long author list or citation section does not move the Library columns.
\end{quote}

\section{Recognize stable and temporary state}\label{recognize-stable-and-temporary-state}

Selection, keyboard focus, menus, and scroll position are temporary interface state. Paper identity, placements, resource properties, read state, provenance, and completed operations persist when saved. A breathing progress message means work is active, not that a paper or folder was created.

\begin{quote}
\textbf{Keep the active context visible}

Before moving, deleting, or attaching, check the focused pane and selected Paper Details. Similar names can identify different targets: a paper placement, one Paper Files resource, or a standalone file.
\end{quote}

\par\noindent\begin{minipage}{\linewidth}

\begin{center}\includegraphics[width=0.96\linewidth,height=0.72\textheight,keepaspectratio]{/Users/lingyuan/PaperMonitor-20260516/assets/user-guide/screenshots/user-guide-window.png}\end{center}

\begin{GewuCaption}
The User Guide window. Help opens a separate searchable manual. The sidebar selects topics; Print Guide opens the complete fixed-layout PDF.
\end{GewuCaption}

\end{minipage}\par\medskip

\begin{GewuRelated}
\textbf{See also:} \hyperref[library-layout]{Library Layout} · \hyperref[navigation]{Navigate with Mouse and Keyboard} · \hyperref[keyboard]{Keyboard Shortcuts}
\end{GewuRelated}

\chapter{Your First Session}\label{first-session}

\begin{GewuLead}
Start with one paper, confirm its identity and context, then inspect or add a local file when the work is file-backed.
\end{GewuLead}

\section{Choose your starting route}\label{choose-your-starting-route}

You do not need to prepare a perfect record before beginning. Start from the strongest evidence you already have, then let the paper object become richer over time.

{\def\LTcaptype{none} % do not increment counter
\GewuTableStart
\begin{longtable}[]{@{}>{\RaggedRight\arraybackslash}p{0.19\textwidth}>{\RaggedRight\arraybackslash}p{0.29\textwidth}>{\RaggedRight\arraybackslash}p{0.36\textwidth}@{}}
\toprule\noalign{}
\rowcolor{GewuTableHeader}
You already have & Start here & What Gewu Insight does first \\
\midrule\noalign{}
\endhead
\bottomrule\noalign{}
\endlastfoot
A local PDF or another research file & Drag it into the target Library folder, or use the folder\textquotesingle s add command. & Preserves a visible file, creates a catalog record, and attempts to extract or resolve identity evidence. \\
A DOI, arXiv identifier, supported URL, or title & Choose \textbf{Add Paper} in the target folder and enter the evidence. & Creates or reuses a paper record, resolves metadata, and offers a download when a supported file is available. \\
A useful result in Journal, Subject, or Author Monitor & Select the result and choose \textbf{Add to Library}. & Adds a Library placement for the cataloged paper while retaining its discovery provenance. \\
\end{longtable}
\GewuTableEnd
}

\begin{quote}
\textbf{External originals stay separate}

Folder Loader and ordinary file import are designed to preserve the external source rather than silently reorganize it. The managed copy and the source path can both remain part of the provenance.
\end{quote}

\section{Add and inspect a paper}\label{add-and-inspect-a-paper}

\begin{enumerate}
\tightlist
\item
  \textbf{Open \textbf{Library}.} Select an existing folder, or create one from the context menu.
\item
  \textbf{Add a paper.} Drag a PDF into the folder, or open the target folder, Control-click blank space in its contents, choose \textbf{Add Paper}, and enter a DOI, arXiv identifier, supported URL, or title.
\item
  \textbf{Select the new item.} Paper Details appears at the right. Gewu Insight attempts to identify the work. A local file is recorded only when you imported one or a supported download succeeded.
\item
  \textbf{Check \textbf{DOI} and \textbf{arXiv}.} Enter a corrected DOI in its field and refresh it, or edit arXiv inline and press \texttt{Return}.
\item
  \textbf{Inspect the available evidence.} If a local PDF exists, choose \textbf{Open in Preview}, or enter Paper Files and press \texttt{Space} for Quick Look. If no file exists yet, use \textbf{Open Website} or download a supported version when one is available.
\end{enumerate}

\begin{GewuResult}
\textbf{Result.} The paper now has a cataloged paper record, a stable Gewu Key, and a Library placement. It may have no local file; that cataloged-only state is valid. A file resource exists only after a file is imported, downloaded, or attached.
\end{GewuResult}

\section{Understand what may now exist}\label{understand-what-may-now-exist}

{\def\LTcaptype{none} % do not increment counter
\GewuTableStart
\begin{longtable}[]{@{}
  >{\raggedright\arraybackslash}p{(\linewidth - 2\tabcolsep) * \real{0.5000}}
  >{\raggedright\arraybackslash}p{(\linewidth - 2\tabcolsep) * \real{0.5000}}@{}}
\toprule\noalign{}
\rowcolor{GewuTableHeader}
\begin{minipage}[b]{\linewidth}\raggedright
Element
\end{minipage} & \begin{minipage}[b]{\linewidth}\raggedright
Role
\end{minipage} \\
\midrule\noalign{}
\endhead
\bottomrule\noalign{}
\endlastfoot
\begin{minipage}[t]{\linewidth}\raggedright
{{Always}\\
\textbf{Catalog record}}\strut
\end{minipage} & {Title, authors, identifiers, source evidence, and a stable internal paper identity.} \\
\begin{minipage}[t]{\linewidth}\raggedright
{{Always}\\
\textbf{Library placement}}\strut
\end{minipage} & {A relationship that makes the paper visible in the folder you selected.} \\
\begin{minipage}[t]{\linewidth}\raggedright
{{When file-backed}\\
\textbf{File record}}\strut
\end{minipage} & {Type, annotation state, source, stored path, fingerprints, and Default status when known.} \\
\begin{minipage}[t]{\linewidth}\raggedright
{{When file-backed}\\
\textbf{Visible local file}}\strut
\end{minipage} & {The actual PDF or attachment that Finder and Preview can open.} \\
\end{longtable}
\GewuTableEnd
}

\begin{GewuCaption}
A cataloged-only paper legitimately has the first two objects and no local file. A file-backed paper also has resource metadata and real bytes. These objects cooperate, but they are not interchangeable.
\end{GewuCaption}

\section{Verify the identity before enriching it}\label{verify-the-identity-before-enriching-it}

\begin{enumerate}
\tightlist
\item
  \textbf{Compare the title and authors with the PDF.} A similar title is useful evidence but not sufficient proof.
\item
  \textbf{Check the DOI and arXiv fields.} A work can have both. Correct an inaccurate identifier before refreshing more metadata or downloading another version.
\item
  \textbf{Read the source and journal information.} Confirm that a publication record belongs to this work rather than a cited paper or a similarly named article.
\item
  \textbf{Keep uncertainty visible.} It is safer to leave a field empty than to confirm a plausible but wrong identifier.
\end{enumerate}

\section{Complete the first Paper Files bundle}\label{complete-the-first-paper-files-bundle}

A cataloged-only paper may have no local file, and nothing is missing merely because Paper Files is empty. Continue with these steps after you import, download, or attach a file and the paper becomes file-backed.

\begin{enumerate}
\tightlist
\item
  \textbf{Select the paper to show \textbf{Paper Files}.} Cards initially show File type and size. Click a card\textquotesingle s size footer to show or hide Annotation and Source.
\item
  \textbf{Classify the available file.} Use its small type, annotation, and source controls. Do not infer these facts from the filename alone.
\item
  \textbf{Add related material if you have it.} Add a supplement, annotated copy, note, presentation, data file, code file, or other attachment. Import first, then classify the resource.
\item
  \textbf{Check the default file.} \textbf{Open in Preview} and Library preview use the resolved default when applicable. Selecting \textbf{Set as Default} changes preference without renaming the resource.
\item
  \textbf{Reveal one file in Finder.} Choose \textbf{Show in Finder} and confirm the real managed file exists where the database expects it.
\end{enumerate}

{\def\LTcaptype{none} % do not increment counter
\GewuTableStart
\begin{longtable}[]{@{}>{\RaggedRight\arraybackslash}p{0.27\textwidth}>{\RaggedRight\arraybackslash}p{0.63\textwidth}@{}}
\toprule\noalign{}
\rowcolor{GewuTableHeader}
Principle & What it means \\
\midrule\noalign{}
\endhead
\bottomrule\noalign{}
\endlastfoot
\textbf{Cataloged and organized} & {The shared paper record and Library placement are durable even when no local file has been added.} \\
\textbf{File-backed} & {After a file is added, its independent file properties and physical bytes can be inspected separately from paper identity and folder context.} \\
\end{longtable}
\GewuTableEnd
}

\section{Optional: verify shared identity}\label{optional-verify-shared-identity}

Copy the paper to a second Library folder, then select it there. The title, Gewu Key, Paper Files, and metadata should remain the same because the second folder contains another placement for the shared paper---not an unrelated duplicate. Delete only the extra placement when you are finished with the check.

\section{Next useful actions}\label{next-useful-actions}

\begin{itemize}
\tightlist
\item
  Use \textbf{Download} when a better publication version is available.
\item
  Use \textbf{Paper Files} to classify imported PDFs or add notes, data, slides, and supplements.
\item
  Add the paper to another folder without creating a second independent work record.
\item
  Back up the project before a large import or migration.
\end{itemize}

\section{Finish with a safe baseline}\label{finish-with-a-safe-baseline}

Your first session is complete when you can answer two questions for every paper: What is this work, and why is it in this folder? If the paper is file-backed, answer two more: Which file is the default, and where are the real bytes? Paper Details, Paper Files, and Finder should answer the applicable questions without contradicting one another.

{\def\LTcaptype{none} % do not increment counter
\GewuTableStart
\begin{longtable}[]{@{}>{\RaggedRight\arraybackslash}p{0.27\textwidth}>{\RaggedRight\arraybackslash}p{0.63\textwidth}@{}}
\toprule\noalign{}
\rowcolor{GewuTableHeader}
Principle & What it means \\
\midrule\noalign{}
\endhead
\bottomrule\noalign{}
\endlastfoot
\textbf{Identity checked} & {The title, authors, DOI, and arXiv evidence do not point to different works.} \\
\textbf{Placement understood} & {The Library folder expresses personal organization, not a second scholarly identity.} \\
\textbf{Resource classified when present} & {If the paper is file-backed, its files can have independent type, annotation, and source descriptions.} \\
\textbf{File visible when present} & {If the paper is file-backed, the expected managed file opens from the app and can be revealed in Finder.} \\
\end{longtable}
\GewuTableEnd
}

\begin{quote}
\textbf{Create the backup habit early}

Before a major import, upgrade, or migration, quit Gewu Insight and back up both the SQLite database and the managed file roots. A database-only copy lacks some bytes; a files-only copy lacks the relationships and state that make the workspace intelligible.
\end{quote}

\begin{GewuRelated}
\textbf{See also:} \hyperref[import-files]{Add Papers and Files} · \hyperref[paper-details]{Use Paper Details} · \hyperref[paper-files]{Understand Paper Files} · \hyperref[backup-migrate]{Back Up or Move to Another Mac}
\end{GewuRelated}

\part{Product Concepts and Architecture}\label{part-2}

\chapter{System Architecture at a Glance}\label{system-overview}

\begin{GewuLead}
Gewu Insight is a local research system that coordinates structured state, visible files, and a personal workspace. It is not simply a web page placed in front of a folder.
\end{GewuLead}

\section{How the parts cooperate}\label{how-the-parts-cooperate}

{\def\LTcaptype{none} % do not increment counter
\GewuTableStart
\begin{longtable}[]{@{}>{\RaggedRight\arraybackslash}p{0.27\textwidth}>{\RaggedRight\arraybackslash}p{0.63\textwidth}@{}}
\toprule\noalign{}
\rowcolor{GewuTableHeader}
Element & Role \\
\midrule\noalign{}
\endhead
\bottomrule\noalign{}
\endlastfoot
{\textbf{Native desktop shell}} & {Starts the app, owns Help and Browser, and connects Preview, Finder, and native file actions to the portable core.} \\
{\textbf{Personal Research Workspace}} & {Library, Spaces, My Publications, and the three Monitor views you see and operate.} \\
{\textbf{External research sources}} & {Publisher pages, feeds, metadata services, scholarly profiles, and preprint services.} \\
{\textbf{Local application service}} & {Interprets your actions, resolves identities, schedules network work, and coordinates database and file changes.} \\
{\textbf{SQLite state core}} & {Owns identity, relationships, resource descriptions, preferences, provenance, and durable workflow state.} \\
{\textbf{Visible Research Files}} & {PDFs, versions, supplements, notes, data, and other files that remain inspectable in the native file manager and document applications.} \\
\end{longtable}
\GewuTableEnd
}

\begin{GewuCaption}
The interface sends an intention to the local service. The service validates it and coordinates SQLite with the relevant files; the UI does not directly treat a Finder folder as the database.
\end{GewuCaption}

The desktop window presents a local interface served from the project on the current computer. In today\textquotesingle s primary release, a macOS shell supplies Finder, Preview, Quick Look, Dock, and native drag-and-drop integration. The Python, SQLite, and web-interface core is designed to support other desktop shells without changing the meaning of the research data. Network sources contribute evidence and downloadable material, but the durable workspace stays local.

\section{What each persistent layer answers}\label{what-each-persistent-layer-answers}

\textbf{Browser is a separate web surface.} Its native WebKit views load websites directly, not through the Python service. The local service stores bookmarks and history in SQLite; WebKit manages a separate persistent website profile. External pages do not receive the Library\textquotesingle s local API bridge. Visiting a page, saving a bookmark, or downloading through Browser does not create a paper or assign a Gewu Key; import the paper explicitly.

{\def\LTcaptype{none} % do not increment counter
\GewuTableStart
\begin{longtable}[]{@{}>{\RaggedRight\arraybackslash}p{0.19\textwidth}>{\RaggedRight\arraybackslash}p{0.29\textwidth}>{\RaggedRight\arraybackslash}p{0.36\textwidth}@{}}
\toprule\noalign{}
\rowcolor{GewuTableHeader}
Layer & Questions it answers & What it does not decide alone \\
\midrule\noalign{}
\endhead
\bottomrule\noalign{}
\endlastfoot
SQLite & What is this paper? Where is it referenced? Which resources belong to it? What happened in the last operation? & It does not contain the bytes of your PDFs and attachments. \\
Filesystem & Which real files can the native file manager, PDF reader, or another application open? & A filename or folder position alone does not define paper identity, version priority, or monitor membership. \\
Personal Research Workspace & How should the shared state be presented for discovery, reading, organization, and action? & It is a projection of saved state, not a separate copy of the collection. \\
\end{longtable}
\GewuTableEnd
}

\begin{quote}
System relationship

Persistent research workspace = structured SQLite state + visible research files, coordinated by the local application.
\end{quote}

\section{Why this separation matters}\label{why-this-separation-matters}

A paper can appear in a Monitor, the Library, a Space, and My Publications without becoming four unrelated papers. It can also acquire a DOI after first arriving as an arXiv preprint, or acquire a publication PDF while retaining earlier versions. The database preserves continuity across those changes; visible files preserve user ownership and inspectability.

\begin{quote}
\textbf{Terminology}

This guide uses \textbf{Personal Research Workspace} for the on-screen projection and \textbf{Visible Research Files} for assets visible to the native filesystem. These are explanatory terms, not menu names. Together they form the visible workspace, but they serve different roles.
\end{quote}

\begin{GewuRelated}
\textbf{See also:} \hyperref[persistent-paper-object]{The Persistent Paper Object} · \hyperref[sqlite-core]{Why SQLite Is the Core} · \hyperref[visible-research-files]{Visible Research Files}
\end{GewuRelated}

\chapter{The Persistent Paper Object}\label{persistent-paper-object}

\begin{GewuLead}
Gewu Insight treats a scholarly work as an object that persists through discovery, publication, file changes, and personal organization---not as a temporary search result or a single PDF.
\end{GewuLead}

\section{One identity can accumulate evidence}\label{one-identity-can-accumulate-evidence}

{\def\LTcaptype{none} % do not increment counter
\GewuTableStart
\begin{longtable}[]{@{}>{\RaggedRight\arraybackslash}p{0.27\textwidth}>{\RaggedRight\arraybackslash}p{0.63\textwidth}@{}}
\toprule\noalign{}
\rowcolor{GewuTableHeader}
Element & Role \\
\midrule\noalign{}
\endhead
\bottomrule\noalign{}
\endlastfoot
{\textbf{Identity evidence}} & {DOI, arXiv identifier, title, authors, journal, dates, and source identifiers.} \\
{\textbf{Resources}} & {Publication PDFs, preprints, accepted manuscripts, supplements, annotations, and attachments.} \\
{\textbf{Contexts}} & {Library folders, Spaces, My Publications, and Monitor relationships.} \\
{\textbf{Durable state}} & {Read status, metadata quality, provenance, download state, retries, and timestamps.} \\
{\textbf{One shared paper identity · one stable Gewu Key}} & {The internal identity stays stable when an external identifier, filename, folder location, or preferred PDF changes.} \\
\end{longtable}
\GewuTableEnd
}

\begin{GewuCaption}
Identifiers help resolve and enrich a paper, while the Gewu Key provides continuity inside Gewu Insight.
\end{GewuCaption}

The central record uses an internal paper identity. A DOI or arXiv identifier is important evidence attached to that identity, but neither is a safe universal primary key: an early preprint may have no DOI, an imported PDF may initially have no reliable identifier, and an incorrect identifier may later be corrected.

\section{Do not collapse four different things}\label{do-not-collapse-four-different-things}

{\def\LTcaptype{none} % do not increment counter
\GewuTableStart
\begin{longtable}[]{@{}>{\RaggedRight\arraybackslash}p{0.19\textwidth}>{\RaggedRight\arraybackslash}p{0.29\textwidth}>{\RaggedRight\arraybackslash}p{0.36\textwidth}@{}}
\toprule\noalign{}
\rowcolor{GewuTableHeader}
Object & Meaning & Example \\
\midrule\noalign{}
\endhead
\bottomrule\noalign{}
\endlastfoot
Paper record & The shared scholarly identity and normalized metadata. & The work that later connects arXiv and a journal DOI. \\
Placement or relation & Why and where the paper appears in your workspace. & A Library folder entry, a shortcut, a My Publications role, or an Author Monitor association. \\
Resource record & A described file or retrievable version belonging to the paper. & Preprint v2, published version, supplementary material, peer review file, or data file. \\
Physical file & The bytes stored at a real path on disk. & A PDF under \nolinkurl{data/Library} or an associated attachment. \\
\end{longtable}
\GewuTableEnd
}

\begin{quote}
System relationship

Paper record {\symbolfont ≠} workspace placement {\symbolfont ≠} resource record {\symbolfont ≠} physical file.
\end{quote}

\section{What happens when the same paper is found again}\label{what-happens-when-the-same-paper-is-found-again}

Journal, Subject, and Author Monitor can encounter the same work from different directions. Identity resolution is designed to reuse one shared paper record and add source or context relationships. Copying a paper to another Library folder likewise creates another placement rather than a second independent identity. Its metadata, Gewu Key, versions, and attachments remain shared.

When the available evidence is insufficient or contradictory, separate candidates can remain until their identifiers or metadata are corrected. A major DOI, title, author, or venue contradiction is quarantined as an identity conflict instead of overwriting the saved paper or forcing a merge. The shared-object model is the resolution target; it does not justify merging uncertain works merely because their titles look similar.

A network observation can be transient while acquisition is running. Once a result is cataloged as a shared paper record, it has a Gewu Key even if a particular Monitor view does not display that key. A Monitor result may still have no local Paper Files because no file has been downloaded or attached; adding it to Library creates a visible placement rather than inventing a new scholarly identity.

Gewu Insight also uses the same catalog and resource infrastructure for an imported standalone file whose scholarly identity is not yet known. In that case, the catalog record preserves the file and its relationships while later identification can enrich it.

\section{Identity corrections should preserve continuity}\label{identity-corrections-should-preserve-continuity}

Editing a DOI or arXiv identifier changes evidence on the existing paper. It should not silently make an unrelated second paper, rename every file, or discard your folder placements and read history. When evidence conflicts, inspect the title, authors, venue, and files before confirming the correction.

\begin{quote}
\textbf{Important}

Similar titles are not proof of identity, and matching filenames are not proof that two resources are the same. Use identifiers, metadata, file fingerprints, and provenance together.
\end{quote}

\begin{GewuRelated}
\textbf{See also:} \hyperref[identifiers]{Edit DOI and arXiv Identifiers} · \hyperref[paper-files]{Understand Paper Files} · \hyperref[personal-research-workspace]{The Personal Research Workspace}
\end{GewuRelated}

\chapter{Why SQLite Is the Core}\label{sqlite-core}

\begin{GewuLead}
SQLite is the source of truth for Gewu Insight because it preserves the identity, relationships, resource registry, and operational memory that neither a filename nor an exported report can express.
\end{GewuLead}

\section{Four families of structured state}\label{four-families-of-structured-state}

{\def\LTcaptype{none} % do not increment counter
\GewuTableStart
\begin{longtable}[]{@{}>{\RaggedRight\arraybackslash}p{0.27\textwidth}>{\RaggedRight\arraybackslash}p{0.63\textwidth}@{}}
\toprule\noalign{}
\rowcolor{GewuTableHeader}
Element & Role \\
\midrule\noalign{}
\endhead
\bottomrule\noalign{}
\endlastfoot
{\textbf{SQLite state authority}} & {A local transactional database stored with the project.} \\
{\textbf{Identity and metadata}} & {Shared papers, external identifiers, normalized title and authors, dates, fields, and Gewu Key.} \\
{\textbf{Relationships and context}} & {Library placements, folder tree, Spaces, My Publications, and Monitor associations.} \\
{\textbf{Resource registry}} & {File types, paths, source evidence, hashes, annotation state, and default preferences.} \\
{\textbf{Operational memory}} & {Fetch runs, import batches, source tasks, reads, settings, failures, and retries; also Browser bookmarks and history.} \\
\end{longtable}
\GewuTableEnd
}

\begin{GewuCaption}
Table names and schema versions can evolve; these four responsibilities are the stable user-level model.
\end{GewuCaption}

The database does not store the PDF bytes themselves. It stores enough structured information to understand the files: which paper they belong to, what file type they represent, where they are, where they came from, whether they are annotated, and which one should be used by default.

\section{Why a folder scan is not an equivalent database}\label{why-a-folder-scan-is-not-an-equivalent-database}

Finder can show that \texttt{paper.pdf} exists. It cannot reliably say whether it is a published version or a preprint, whether it is annotated, where it came from, whether the same work is also in My Publications, or which file is the manual default. Those facts live in SQLite.

\section{Transactions, relationships, and WAL}\label{transactions-relationships-and-wal}

{\def\LTcaptype{none} % do not increment counter
\GewuTableStart
\begin{longtable}[]{@{}>{\RaggedRight\arraybackslash}p{0.27\textwidth}>{\RaggedRight\arraybackslash}p{0.63\textwidth}@{}}
\toprule\noalign{}
\rowcolor{GewuTableHeader}
Database mechanism & User-facing purpose \\
\midrule\noalign{}
\endhead
\bottomrule\noalign{}
\endlastfoot
Transactions & Keep a multi-step state change together so a partly applied edit is less likely to become the visible truth. \\
Foreign-key relationships & Connect a resource or placement to the correct shared paper instead of relying on a coincidentally similar filename. \\
WAL (write-ahead logging) & Lets foreground reading and background updates coexist more safely than repeated full-file rewrites. \\
Schema migration & Adds new structures and backfills required state during an upgrade rather than asking the user to rebuild the database. \\
\end{longtable}
\GewuTableEnd
}

WAL improves normal concurrency; it is not a substitute for a coherent backup. Close Gewu Insight before copying its live database, and back up the managed files with it.

\section{Durable preferences and shipped defaults}\label{durable-preferences-and-shipped-defaults}

Journal subscriptions are durable SQLite state as well. The database preserves which journals are enabled, their identifiers and source settings, and their display order. The journal catalog shipped with a newer app version is a set of defaults, not an instruction to overwrite your personal list.

On the first upgrade that introduces durable subscriptions, Gewu Insight seeds the database from the existing journal configuration and history. After that one-time seed, adding, disabling, restoring, or reordering a journal changes the database record. Replacing source code or rebuilding the app does not silently re-enable a journal you removed or discard journals you added.

\begin{quote}
\textbf{Cross-layer changes}

SQLite transactions protect database state. A file move and its database update cannot be one native SQLite transaction, so the local service coordinates the steps and uses validation or compensating recovery when a filesystem operation fails.
\end{quote}

\section{Exports and provenance}\label{exports-and-provenance}

JSON, HTML, Markdown, Word, LaTeX, and PDF outputs are derived artifacts. They can be regenerated, shared, or inspected, but editing them does not alter the source records. Raw source evidence may also be retained inside SQLite for provenance; this does not turn an external JSON export into a second source of truth.

\section{Paper-level Monitor provenance}\label{paper-level-monitor-provenance}

Normalized paper-source observations are the authority for how a canonical paper was encountered. Each observation identifies a source kind, provider name, and stable source key, and keeps first-seen and last-seen times. The same paper and effective source update one observation instead of increasing the provenance count indefinitely.

Journal Monitor also connects each observed paper to the relevant fetch run. The run supplies the journal, requested date window, status, and timing; the paper-level observation records that the reporting source saw this particular work. Subject and Author Monitor relationships can point to the same canonical paper, so three discovery routes add context without manufacturing three scholarly identities.

A bibliographic metadata refresh can add a newly proven source, but it cannot erase an earlier Journal, Subject, or Author observation. Display counts and the older raw-source list are compatibility projections of the normalized observations rather than an independent truth. When two paper identities are safely merged, their source observations move to the surviving canonical paper and duplicate observations converge.

\begin{quote}
\textbf{Discovery is not storage}

A provenance relationship does not by itself create a Library placement or a local PDF. It records how the work entered view. Organization and file acquisition remain separate, explicit actions.
\end{quote}

\section{What the app does during an upgrade}\label{what-the-app-does-during-an-upgrade}

On startup, the storage layer reads the database schema version, checks for required tables and fields, applies additive migrations in order, ensures required roots exist, and backfills stable keys or resource records where older data needs them. A normal migration preserves paper identities, folder placements, monitored-journal choices, fetch provenance, and file paths. This is why a project upgrade should preserve the existing database and files rather than replace them with an empty database.

\begin{quote}
\textbf{Do not edit the live database manually}

A seemingly small manual change can violate a relationship that is not visible in one table. Use the app for normal work and make a closed, complete backup before any diagnostic intervention.
\end{quote}

\begin{GewuRelated}
\textbf{See also:} \hyperref[backup-migrate]{Back Up or Move to Another Mac} · \hyperref[privacy]{Local Data and Privacy} · \hyperref[storage-reference]{Storage Reference}
\end{GewuRelated}

\chapter{Visible Research Files}\label{visible-research-files}

\begin{GewuLead}
Your PDFs and attachments remain real local files. Gewu Insight adds a resource model around them so several versions and related materials can belong to one paper without hiding them inside an opaque cloud library.
\end{GewuLead}

\section{One paper, one file bundle}\label{one-paper-one-file-bundle}

{\def\LTcaptype{none} % do not increment counter
\GewuTableStart
\begin{longtable}[]{@{}>{\RaggedRight\arraybackslash}p{0.27\textwidth}>{\RaggedRight\arraybackslash}p{0.63\textwidth}@{}}
\toprule\noalign{}
\rowcolor{GewuTableHeader}
Element & Role \\
\midrule\noalign{}
\endhead
\bottomrule\noalign{}
\endlastfoot
{\textbf{Persistent Paper Object}} & {Owns the resource relationships and decides which usable file is the default.} \\
{\textbf{Default file}} & {The automatically ranked or explicitly preferred valid file used by Open in Preview and Library preview.} \\
{\textbf{Associated Files}} & {Earlier versions, supplements, annotated copies, presentations, notes, data, code, images, and other attachments.} \\
{\textbf{Paper Files}} & {Flat peer records with independent type, annotation state, and source.} \\
{\textbf{Supplements}} & {Publisher or author-provided supporting files.} \\
{\textbf{Files and figures}} & {User material such as notes, slides, data, code, and extracted or attached images.} \\
\end{longtable}
\GewuTableEnd
}

\begin{GewuCaption}
The Paper Files pane presents this as a compact bundle rather than another Finder hierarchy.
\end{GewuCaption}

\section{Managed and linked resources}\label{managed-and-linked-resources}

{\def\LTcaptype{none} % do not increment counter
\GewuTableStart
\begin{longtable}[]{@{}>{\RaggedRight\arraybackslash}p{0.19\textwidth}>{\RaggedRight\arraybackslash}p{0.29\textwidth}>{\RaggedRight\arraybackslash}p{0.36\textwidth}@{}}
\toprule\noalign{}
\rowcolor{GewuTableHeader}
Storage mode & Meaning & Safe expectation \\
\midrule\noalign{}
\endhead
\bottomrule\noalign{}
\endlastfoot
Managed & The active stored path is inside the Gewu Insight project data area. & Move, classify, replace, or delete it through the app so paths and relationships change together. \\
Linked & The active stored path remains outside the managed project data area. & Availability depends on that external path, and the app should not silently reorganize it. \\
\end{longtable}
\GewuTableEnd
}

A resource can be managed while retaining a separate external source path as provenance. The existence of that source path does not by itself make the managed copy a linked resource, and the external original remains protected.

Folder Loader is a migration-oriented workflow. It imports a visible collection as a collection instead of asking you to classify every incoming file as an attachment before identity resolution. The original source path and the managed result remain conceptually distinct.

\section{How the default file is resolved}\label{how-the-default-file-is-resolved}

An explicit manual default takes precedence while it still exists and is not classified as Invalid file. Otherwise, usable files are ranked as Corrected published version, Published version, Accepted manuscript, then Preprint. If no Default or main PDF exists, Browser Save establishes its first article PDF as a fallback Default even when its type is Unknown. This keeps the version classification intact.

\textbf{Supplementary material}, \textbf{Peer review file}, \textbf{Other}, and \textbf{Invalid file} do not become the automatic default. Other Unknown attachments remain peers until selected manually. Any valid resource can still be selected manually; \textbf{Invalid file} is never eligible and remains visible only for audit and repair.

\begin{quote}
\textbf{Published files}

The user-facing type is \textbf{Published version}. A later formal correction can be classified as \textbf{Corrected published version}; the correction notice itself remains a separate peer.
\end{quote}

\section{Fingerprints and annotations}\label{fingerprints-and-annotations}

When a PDF is imported, attached, or downloaded, its exact SHA-256, content-oriented fingerprint, annotation count, and annotated state are computed as the resource is registered. They do not depend on a later refresh. Exact SHA-256 answers ``are these bytes identical?'' Adding a highlight or note normally changes those bytes. The content fingerprint is designed to recognize the same underlying document while discounting annotation-only changes, allowing clean and annotated copies to be related without treating them as blindly interchangeable.

Gewu Insight counts user-created comments and markup---such as highlights, notes, drawings, stamps, and form edits---as annotations. Structural PDF elements such as publisher hyperlinks, internal navigation links, and ordinary page actions do not make a freshly downloaded publisher or arXiv PDF \textbf{Annotated}.

An exact duplicate can be reused instead of stored twice. A clean PDF and an annotated copy may have different exact hashes but the same content fingerprint; both can be retained when the annotation state is meaningful.

\section{Why Finder and the UI do not look identical}\label{why-finder-and-the-ui-do-not-look-identical}

The filesystem must expose practical paths and folders; the UI must expose the meaning of one paper bundle. A Library paper can therefore retain a compatibility primary file in the visible topic folder while related files live under a mirrored Associated Files path. Paper Files deliberately presents all resource records as flat peers instead of mirroring those physical directories.

\begin{quote}
\textbf{Important}

Do not manually move or rename managed resources in Finder. The bytes may still exist, but the saved path and live workspace can disagree. Use an app-supported move, repair, or restore; do not assume a restart will rediscover an arbitrary manual move.
\end{quote}

\begin{GewuRelated}
\textbf{See also:} \hyperref[paper-files]{Understand Paper Files} · \hyperref[versions-current]{Choose File Types and the Default File} · \hyperref[storage-model]{How Gewu Insight Stores Files}
\end{GewuRelated}

\chapter{The Personal Research Workspace}\label{personal-research-workspace}

\begin{GewuLead}
The five research views---Library, My Publications, Journal Monitor, Subject Monitor, and Author Monitor---share research objects. Browser is the sixth main tab, with its own website session. Moving between research views does not reset a paper\textquotesingle s identity.
\end{GewuLead}

{\def\LTcaptype{none} % do not increment counter
\GewuTableStart
\begin{longtable}[]{@{}>{\RaggedRight\arraybackslash}p{0.27\textwidth}>{\RaggedRight\arraybackslash}p{0.63\textwidth}@{}}
\toprule\noalign{}
\rowcolor{GewuTableHeader}
Element & Role \\
\midrule\noalign{}
\endhead
\bottomrule\noalign{}
\endlastfoot
{\textbf{Journal Monitor}} & {Discovers work by publication venue and date.} \\
{\textbf{Subject Monitor}} & {Discovers work through topic definitions and supported feeds or searches.} \\
{\textbf{Author Monitor}} & {Discovers work through a researcher relationship.} \\
{\textbf{One shared paper identity}} & {Metadata, identifiers, resources, provenance, and state accumulate here.} \\
{\textbf{Library and shortcuts}} & {A durable managed collection with logical placements and fast entry points.} \\
{\textbf{Spaces}} & {Independent organizational roots for projects, courses, or collaborations.} \\
{\textbf{My Publications}} & {Your authorship, role, metrics, ordering, and export context on the same paper.} \\
\end{longtable}
\GewuTableEnd
}

\begin{GewuCaption}
The same paper can participate in several contexts. Each context adds meaning; it does not require an unrelated duplicate record.
\end{GewuCaption}

\section{Views are relational projections}\label{views-are-relational-projections}

A Library folder answers ``where do I organize this?'' My Publications answers ``how does this relate to my authorship?'' Author Monitor answers ``how does this relate to this researcher?'' These are different relationships to the same possible work. The interface queries those relationships and projects the relevant controls, status, and files for the current task.

\section{Library placements and shortcuts}\label{library-placements-and-shortcuts}

A paper can have more than one Library placement. Copying it to another folder reuses the shared identity and resources. A shortcut tab is another entry point to a Library location, not another paper collection. A Space is an independent organizational root, but papers placed there can still reuse the same identity model.

\section{Discovery does not require immediate file management}\label{discovery-does-not-require-immediate-file-management}

A Monitor result may remain a discovery relationship with metadata and status but no local Paper Files. A cataloged work already has a Gewu Key, which Author Monitor also shows. Downloading or attaching a file creates a resource relationship; adding the paper to Library creates a managed placement. Neither action should manufacture a second work merely because the user entered from another tab.

\section{Why the same action can have different scope}\label{why-the-same-action-can-have-different-scope}

Removing an author from Author Monitor removes that monitored-author relationship. Removing a shortcut removes an entry point. Moving a Library placement changes organization. Deleting a managed resource changes the paper\textquotesingle s file bundle. The selected view supplies scope, while SQLite preserves what remains shared.

\begin{quote}
\textbf{A useful mental model}

The UI shows a personal map of a shared object graph. It is ``visible'' because you can navigate and act on it, but SQLite---not pixels on the screen---preserves the graph.
\end{quote}

\begin{GewuRelated}
\textbf{See also:} \hyperref[organize]{Folders, Shortcuts, and Spaces} · \hyperref[my-publications]{My Publications} · \hyperref[state-and-scope]{State, Scope, and Safe Change}
\end{GewuRelated}

\chapter{From Discovery to Durable Memory}\label{research-loop}

\begin{GewuLead}
Gewu Insight connects monitoring, identity resolution, version handling, and personal organization into one repeatable research loop.
\end{GewuLead}

{\def\LTcaptype{none} % do not increment counter
\GewuTableStart
\begin{longtable}[]{@{}>{\RaggedRight\arraybackslash}p{0.27\textwidth}>{\RaggedRight\arraybackslash}p{0.63\textwidth}@{}}
\toprule\noalign{}
\rowcolor{GewuTableHeader}
Element & Role \\
\midrule\noalign{}
\endhead
\bottomrule\noalign{}
\endlastfoot
{\textbf{1. Acquire}} & {Journal, Subject, Author, Folder Loader, or a direct identifier introduces a candidate.} \\
{\textbf{2. Identify}} & {DOI, arXiv, title, authors, source evidence, and existing records are reconciled.} \\
{\textbf{3. Resolve}} & {Metadata, file types, annotation states, sources, paths, hashes, and a default preference are established.} \\
{\textbf{4. Remember}} & {Placements, authorship, reading state, provenance, completion, and failures persist.} \\
\end{longtable}
\GewuTableEnd
}

\begin{GewuCaption}
The loop is continuous because a work can acquire better evidence, a new publication version, or a new personal context later.
\end{GewuCaption}

\section{A paper can change without losing its history}\label{a-paper-can-change-without-losing-its-history}

\begin{enumerate}
\tightlist
\item
  \textbf{Acquire an early work.} Subject Monitor finds an arXiv preprint. Its source and discovery context are stored.
\item
  \textbf{Identify the work.} The title, authors, and arXiv identifier connect it to one shared paper identity.
\item
  \textbf{Resolve its first resource.} A preprint PDF becomes the automatic default because no higher-priority file is available.
\item
  \textbf{Remember your context.} You add it to a Library folder, read it, and attach notes or data.
\item
  \textbf{Acquire new evidence later.} Journal Monitor or a metadata refresh finds the published DOI and a published file.
\item
  \textbf{Resolve without erasing.} The published file becomes the automatic default unless a manual choice is sticky; the preprint and its provenance remain peer records.
\end{enumerate}

\section{Operational state is part of scientific memory}\label{operational-state-is-part-of-scientific-memory}

A result count alone is ambiguous. The app distinguishes work that was never attempted from work that completed with no results, partially completed, failed, or entered a cooldown before retry. It also remembers imports, downloads, metadata refreshes, source tasks, and read state so the next action can be informed rather than blindly repeated.

{\def\LTcaptype{none} % do not increment counter
\GewuTableStart
\begin{longtable}[]{@{}>{\RaggedRight\arraybackslash}p{0.19\textwidth}>{\RaggedRight\arraybackslash}p{0.29\textwidth}>{\RaggedRight\arraybackslash}p{0.36\textwidth}@{}}
\toprule\noalign{}
\rowcolor{GewuTableHeader}
State & Meaning & Reason to preserve it \\
\midrule\noalign{}
\endhead
\bottomrule\noalign{}
\endlastfoot
Empty and complete & The requested source or date was checked and returned no matching records. & Prevents ``zero'' from being mistaken for ``not run.'' \\
Partial & Some sources completed while another did not. & Allows a targeted retry instead of discarding useful completed work. \\
Failed & An operation produced a specific error. & Makes the failure visible and avoids pretending the result is authoritative. \\
Cooldown or retry scheduled & A remote service asked the app to wait or a transient error needs backoff. & Reduces repeated traffic and protects normal foreground use. \\
\end{longtable}
\GewuTableEnd
}

\section{Background work should not seize the foreground}\label{background-work-should-not-seize-the-foreground}

Fetches, metadata refreshes, and imports report progress in reserved status areas and save their state. You should still be able to browse already stored material. Where callers share an online service, service-specific pacing, locks, cooldowns, and retry policies coordinate their requests. These controls reduce unnecessary traffic; they cannot guarantee that a publisher will never rate-limit or block a network address.

\begin{quote}
\textbf{The loop is user-directed}

Gewu Insight preserves the consequences of what you monitor, save, read, ignore, classify, and publish. It does not silently replace your research judgment with an invisible global keyword filter.
\end{quote}

\begin{GewuRelated}
\textbf{See also:} \hyperref[journal-monitor]{Journal Monitor} · \hyperref[subject-monitor]{Subject Monitor} · \hyperref[author-monitor]{Author Monitor}
\end{GewuRelated}

\chapter{State, Scope, and Safe Change}\label{state-and-scope}

\begin{GewuLead}
A safe local research workspace changes only the object and relationship you intended while preserving shared identity, resources, provenance, and recoverable state.
\end{GewuLead}

\section{Eight invariants behind normal operations}\label{eight-invariants-behind-normal-operations}

{\def\LTcaptype{none} % do not increment counter
\GewuTableStart
\begin{longtable}[]{@{}>{\RaggedRight\arraybackslash}p{0.27\textwidth}>{\RaggedRight\arraybackslash}p{0.63\textwidth}@{}}
\toprule\noalign{}
\rowcolor{GewuTableHeader}
Principle & What it means \\
\midrule\noalign{}
\endhead
\bottomrule\noalign{}
\endlastfoot
\textbf{One shared identity} & {The same scholarly work should not fragment merely because several sources or folders found it.} \\
\textbf{Many contexts, not blind copies} & {Folders, Spaces, monitors, and My Publications add relationships around the shared paper.} \\
\textbf{Deletion has scope} & {Removing one relationship should not destroy unrelated relationships or shared resources.} \\
\textbf{State survives restart} & {Saved reading and workflow state, retries, failures, versions, and provenance are durable.} \\
\textbf{Repeated operations converge} & {Refreshing, importing, or downloading the same evidence should not create endless duplicates.} \\
\textbf{Good metadata is protected} & {A weak or empty response should not silently erase stronger saved evidence.} \\
\textbf{Database and files agree} & {Saved resource paths should point to the intended real files after an app-managed move.} \\
\textbf{Failure remains inspectable} & {Deletion, import, move, and network failure should preserve enough state to recover or retry.} \\
\end{longtable}
\GewuTableEnd
}

\section{Read an action by its scope}\label{read-an-action-by-its-scope}

{\def\LTcaptype{none} % do not increment counter
\GewuTableStart
\begin{longtable}[]{@{}>{\RaggedRight\arraybackslash}p{0.19\textwidth}>{\RaggedRight\arraybackslash}p{0.29\textwidth}>{\RaggedRight\arraybackslash}p{0.36\textwidth}@{}}
\toprule\noalign{}
\rowcolor{GewuTableHeader}
Action & Intended scope & What should remain \\
\midrule\noalign{}
\endhead
\bottomrule\noalign{}
\endlastfoot
Remove an Author Monitor profile & The monitored-author relation and its profile context. & Shared papers, Library placements, My Publications, and managed files. \\
Remove a shortcut & That fast entry point. & The underlying Library folder and its papers. \\
Copy a paper to another folder & Add another placement. & One shared paper identity and resource bundle. \\
Move a paper in the app & Change its managed placement and corresponding stored paths. & Gewu Key, metadata, versions, and associations. \\
Make another file default & Change the reading preference without physical promotion. & The file\textquotesingle s original name, location, and all other saved resources. \\
Delete an associated file & That resource relationship and its managed file, subject to confirmation. & The paper identity and unrelated resources. \\
\end{longtable}
\GewuTableEnd
}

\section{Recovery is part of the design}\label{recovery-is-part-of-the-design}

Recycle keeps ordinary removals recoverable. Import batches and operation status provide evidence about what ran. Resource paths, hashes, and provenance help diagnose missing or duplicate files. The app can migrate older storage structure forward, but migration is not a reason to discard the previous database or original files.

\section{Back up one coherent research asset}\label{back-up-one-coherent-research-asset}

A reliable backup contains the SQLite database and the managed file roots from the same closed application state. Copying only the database preserves descriptions without every artifact; copying only folders preserves bytes without all identities, relationships, classifications, and operation history. A portable desktop copy is useful for sharing visible folder material, but it is not automatically a complete live-workspace backup.

\begin{quote}
\textbf{Before a major update or migration}

Quit Gewu Insight, preserve the entire project data set, verify the backup can be read, then update the software in place. Do not create a fresh empty database and attempt to reconstruct the workspace from export files.
\end{quote}

\begin{GewuRelated}
\textbf{See also:} \hyperref[recycle]{Delete and Recover Items} · \hyperref[backup-migrate]{Back Up or Move to Another Mac} · \hyperref[troubleshooting]{Troubleshooting}
\end{GewuRelated}

\part{User Manual}\label{part-3}

\chapter{Workspace Navigation}\label{chapter-11}

\section{Library Layout}\label{library-layout}

\begin{GewuLead}
The Library uses a Finder-style column view, but the database---not the visible folder alone---defines the paper and its relationships.
\end{GewuLead}

\par\noindent\begin{minipage}{\linewidth}

\begin{center}\includegraphics[width=0.96\linewidth,height=0.72\textheight,keepaspectratio]{/Users/lingyuan/PaperMonitor-20260516/assets/user-guide/screenshots/library-paper-files.png}\end{center}

\begin{GewuCaption}
Library and Paper Files. This real paper has two already-classified files: Preprint and Published version. The published file is marked Default. Its folder placement, related files, and Paper Details remain connected.
\end{GewuCaption}

\end{minipage}\par\medskip

\subsection{Read the columns from left to right}\label{read-the-columns-from-left-to-right}

\begin{enumerate}
\tightlist
\item
  \textbf{Source:} the main Library, a shortcut, or a Space.
\item
  \textbf{Folders:} each selection opens its children in the next column.
\item
  \textbf{Items:} the papers and other files in the selected folder.
\item
  \textbf{Paper Files:} shown only when a paper or managed file can expose related files.
\item
  \textbf{Paper Details:} metadata and actions for the selected work.
\end{enumerate}

\subsection{Use the breadcrumb}\label{use-the-breadcrumb}

The breadcrumb at the bottom shows the active path. Select a segment to return to that location. The \textbf{List View / Table View} control at the lower right changes the representation of items without moving them.

\subsection{Folder Details}\label{folder-details}

When you select a folder, the right side shows its name, type, and the number of child folders and items. This is informational; it is not another folder level.

\begin{GewuRelated}
\textbf{See also:} \hyperref[organize]{Folders, Shortcuts, and Spaces} · \hyperref[search-views]{Search, List View, and Table View} · \hyperref[paper-files]{Understand Paper Files}
\end{GewuRelated}

\section{Navigate with Mouse and Keyboard}\label{navigation}

\begin{GewuLead}
Library and Space columns form one keyboard path with Paper Files. Focus moves into Paper Files with the Right Arrow and returns to the selected paper with the Left Arrow.
\end{GewuLead}

\subsection{Move through folders and papers}\label{move-through-folders-and-papers}

\begin{itemize}
\tightlist
\item
  Press \texttt{{\symbolfont ↑}} or \texttt{{\symbolfont ↓}} to change the selection in the focused column.
\item
  Press \texttt{{\symbolfont →}} on a folder to enter it.
\item
  Press \texttt{{\symbolfont ←}} to return to its parent.
\item
  Double-click a Library paper row to open its Default file: PDFs open in macOS Preview; other formats use their macOS application. An unclassified legacy import falls back to its available current file.
\end{itemize}

\subsection{Move between the source rail and Library columns}\label{move-between-the-source-rail-and-library-columns}

\begin{enumerate}
\tightlist
\item
  \textbf{Return to the leftmost Library column.} Press \texttt{{\symbolfont ←}} once more to move focus to the vertical source rail.
\item
  \textbf{Press \texttt{{\symbolfont ↑}} or \texttt{{\symbolfont ↓}} on the rail.} Move through Library, shortcuts, and Spaces. The highlighted source opens immediately; you do not need to press Return.
\item
  \textbf{Press \texttt{{\symbolfont →}}.} Focus enters the first visible Library column for that source. Its active row changes from the quieter background selection to the stronger blue content selection.
\end{enumerate}

While the source rail owns focus, the current row in the adjacent column remains visible in a quieter blue so you do not lose your place. The rail keeps focus while you move between sources and releases it only after you press \texttt{{\symbolfont →}}. Once you are navigating Library content, pressing \texttt{{\symbolfont ↑}} at the first row stays in the Library instead of unexpectedly moving to the app tabs across the top.

\subsection{Enter and leave Paper Files}\label{enter-and-leave-paper-files}

\begin{enumerate}
\tightlist
\item
  \textbf{Select a paper in the main item column.} Paper Files appears to its right. No Paper Files card is highlighted yet.
\item
  \textbf{Press \texttt{{\symbolfont →}}.} Focus enters Paper Files and the active card gains a light-blue selection. The main paper remains visible with a quieter selection color.
\item
  \textbf{Press \texttt{{\symbolfont ↑}} or \texttt{{\symbolfont ↓}}.} Move through the flat Paper Files list.
\item
  \textbf{Press \texttt{{\symbolfont ←}}.} Focus returns to the main paper. Paper Files no longer shows an active card.
\end{enumerate}

The same focus rules apply in the main Library, a shortcut, and every Space. Paper Files is a database-backed view of the selected work rather than a physical child folder, but the keyboard transition is intentionally continuous: \texttt{{\symbolfont →}} enters the bundle, \texttt{{\symbolfont ←}} returns to the same paper, and the next \texttt{{\symbolfont ↑}} or \texttt{{\symbolfont ↓}} moves in the main item column.

\subsection{Preview while navigating}\label{preview-while-navigating}

In Paper Files, press \texttt{Space} to open Quick Look. While Quick Look is open, use \texttt{{\symbolfont ↑}} and \texttt{{\symbolfont ↓}} to follow adjacent files. Press \texttt{Space} again or close Quick Look to dismiss it.

\subsection{Use the Quick Look window}\label{use-the-quick-look-window}

\begin{itemize}
\tightlist
\item
  For a PDF, use the thumbnail strip to jump to a page and scroll the page area independently.
\item
  Drag the title bar to reposition the preview. Drag its resize corner when you need more reading space.
\item
  Choose \textbf{Open in Preview} when you want the full macOS Preview application.
\item
  When Quick Look is already open, changing the Paper Files selection updates the preview without requiring another open action.
\end{itemize}

\begin{quote}
\textbf{Note}

Single-click selects a Paper Files item. Double-click opens it. This prevents an accidental external app launch while you are classifying files.
\end{quote}

\begin{GewuRelated}
\textbf{See also:} \hyperref[paper-files]{Understand Paper Files} · \hyperref[quick-look]{Preview PDFs with Quick Look} · \hyperref[keyboard]{Keyboard Shortcuts}
\end{GewuRelated}

\section{Folders, Shortcuts, and Spaces}\label{organize}

\begin{GewuLead}
Folders organize a root. Shortcuts provide another entrance to the Library. Spaces keep an independent root for a project or research context.
\end{GewuLead}

\par\noindent\begin{minipage}{\linewidth}

\begin{center}\includegraphics[width=0.96\linewidth,height=0.72\textheight,keepaspectratio]{/Users/lingyuan/PaperMonitor-20260516/assets/user-guide/screenshots/space-workspace.png}\end{center}

\begin{GewuCaption}
A Space. Spaces use the same paper identity and Paper Files model as the Library while keeping an independent organizational root.
\end{GewuCaption}

\end{minipage}\par\medskip

\subsection{Choose the right structure}\label{choose-the-right-structure}

{\def\LTcaptype{none} % do not increment counter
\GewuTableStart
\begin{longtable}[]{@{}>{\RaggedRight\arraybackslash}p{0.19\textwidth}>{\RaggedRight\arraybackslash}p{0.29\textwidth}>{\RaggedRight\arraybackslash}p{0.36\textwidth}@{}}
\toprule\noalign{}
\rowcolor{GewuTableHeader}
Structure & Use it for & What it changes \\
\midrule\noalign{}
\endhead
\bottomrule\noalign{}
\endlastfoot
Library folder & A durable topic hierarchy inside the main collection. & Creates or moves a managed path. \\
Shortcut tab & Fast access to a Library location you use often. & Adds another entry point; it does not duplicate the paper record. \\
Space & An independent project, course, collaboration, or temporary collection. & Uses its own managed root under \nolinkurl{data/Spaces}. \\
\end{longtable}
\GewuTableEnd
}

\subsection{Create, rename, or move}\label{create-rename-or-move}

\begin{enumerate}
\tightlist
\item
  \textbf{Control-click the relevant folder area.} Choose \textbf{New Folder}, \textbf{Rename}, \textbf{Copy}, or \textbf{Paste}.
\item
  \textbf{For a move, drag the item to its destination.} Wait for the destination to indicate it can receive the item, then release.
\item
  \textbf{Check the breadcrumb and Paper Details.} The visible path and folder membership update immediately.
\end{enumerate}

Shortcut and Space tabs can be dragged to change their order. Their inactive colors distinguish the tab groups; the active tab uses the common selected style. With the keyboard, press \texttt{{\symbolfont ←}} from the leftmost content column to focus this rail, use \texttt{{\symbolfont ↑}} and \texttt{{\symbolfont ↓}} to open another source, then press \texttt{{\symbolfont →}} to enter its content.

\par\noindent\begin{minipage}{\linewidth}

\begin{center}\includegraphics[width=0.96\linewidth,height=0.72\textheight,keepaspectratio]{/Users/lingyuan/PaperMonitor-20260516/assets/user-guide/screenshots/add-shortcut.png}\end{center}

\begin{GewuCaption}
Add Shortcut. Choose an existing Library folder through the column picker. Add creates a navigation shortcut to that folder.
\end{GewuCaption}

\end{minipage}\par\medskip

\par\noindent\begin{minipage}{\linewidth}

\begin{center}\includegraphics[width=0.96\linewidth,height=0.72\textheight,keepaspectratio]{/Users/lingyuan/PaperMonitor-20260516/assets/user-guide/screenshots/add-space.png}\end{center}

\begin{GewuCaption}
Add Space. Give the independent workspace a name, then choose Add. Cancel closes this form without creating a Space.
\end{GewuCaption}

\end{minipage}\par\medskip

\subsection{Color folders and remove navigation entries}\label{color-folders-and-remove-navigation-entries}

Control-click an ordinary folder and choose \textbf{Color} to change its folder color. Protected roots do not offer this command. On the source rail, \textbf{Remove Shortcut} removes only the shortcut, leaving its folder and papers intact. \textbf{Remove Empty Space} works only after both its managed contents and remaining files are empty; move the Space into Library if you want to keep its hierarchy.

\subsection{Work with several entries}\label{work-with-several-entries}

Command-click to add or remove individual entries from a selection. Shift-click, or hold Shift while using the arrow keys, to extend a range in the current folder. Copy the selected entries with \texttt{{\symbolfont ⌘}}\texttt{C}, paste into the active destination with \texttt{{\symbolfont ⌘}}\texttt{V}, or press \texttt{Delete} to apply the normal reference-aware deletion rules.

To include an existing Library paper in your own publication list, select one or more paper rows, Control-click, and choose \textbf{Add to My Publications}. This adds the publication relationship to the same paper and Gewu Key; it does not remove the Library placement. Review authorship roles in My Publications afterward.

\subsection{Copy a paper to another folder}\label{copy-a-paper-to-another-folder}

A Library copy creates another reference to the same paper record rather than a second unrelated identity. The local filesystem may use a secondary symbolic link where appropriate. Paper versions and attachments remain associated with the work.

\subsection{Move an entire Space into Library}\label{move-an-entire-space-into-library}

Use this when a temporary or project Space should become part of the permanent Library hierarchy.

\begin{enumerate}
\tightlist
\item
  \textbf{Select the Space on the left rail.} Open the Space that you want to keep as a normal Library folder.
\item
  \textbf{Choose \textbf{Move Space to Library}.} The move sheet explains that the Space will become a folder with the same name.
\item
  \textbf{Choose \textbf{Destination in My Library}.} Select the Library folder that should contain the moved Space.
\item
  \textbf{Confirm the move.} The complete Space hierarchy, including its files and subfolders, moves together.
\end{enumerate}

When the move finishes, the Space tab is removed and Library opens the moved folder. Canonical paper identities, folder placements, and Paper Files resource associations remain connected; the operation does not create duplicate paper records.

\begin{quote}
\textbf{Important}

Move managed papers in Gewu Insight, not directly in Finder. The app updates the visible file, associated bundle, and live database paths together.
\end{quote}

\begin{GewuRelated}
\textbf{See also:} \hyperref[storage-model]{How Gewu Insight Stores Files} · \hyperref[recycle]{Delete and Recover Items}
\end{GewuRelated}

\section{Search, List View, and Table View}\label{search-views}

\begin{GewuLead}
Search locates items across the managed collection. List View emphasizes browsing; Table View emphasizes metadata comparison.
\end{GewuLead}

\par\noindent\begin{minipage}{\linewidth}

\begin{center}\includegraphics[width=0.96\linewidth,height=0.72\textheight,keepaspectratio]{/Users/lingyuan/PaperMonitor-20260516/assets/user-guide/screenshots/library-search.png}\end{center}

\begin{GewuCaption}
Search Library. The search field filters the current workspace. Select a result or use Show in Library to reveal its folder location.
\end{GewuCaption}

\end{minipage}\par\medskip

\subsection{Search the collection}\label{search-the-collection}

\begin{enumerate}
\tightlist
\item
  \textbf{Select the search control in the top-right corner.}
\item
  \textbf{Enter a title, author, journal, identifier, or filename.} Results may come from the Library and Spaces.
\item
  \textbf{Select a result.} Its Paper Details and managed path become available.
\item
  \textbf{Clear the search.} Return to the previous Library location.
\end{enumerate}

\subsection{Switch views}\label{switch-views}

Choose \textbf{View \textgreater{} List View} or \textbf{View \textgreater{} Table View}. You can also use \texttt{{\symbolfont ⇧}}\texttt{{\symbolfont ⌘}}\texttt{V}. Table columns can be sorted or resized where available; changing the view does not rename, move, or rewrite files. Active Library search uses its result list and disables Table View. Clear the search to use Table View for the current folder again.

\subsection{Sort a folder}\label{sort-a-folder}

Control-click the folder contents and choose \textbf{Sort by Name} or \textbf{Sort by Modified Time}. Sorting changes presentation, not the paper identity.

\begin{GewuRelated}
\textbf{See also:} \hyperref[library-layout]{Library Layout} · \hyperref[keyboard]{Keyboard Shortcuts}
\end{GewuRelated}

\chapter{Papers and Files}\label{chapter-12}

\section{Add Papers and Files}\label{import-files}

\begin{GewuLead}
A normal Library import creates a managed local copy and a database record. The original source file remains where it was.
\end{GewuLead}

\subsection{Drag a paper into a folder}\label{drag-a-paper-into-a-folder}

\begin{enumerate}
\tightlist
\item
  \textbf{Open the destination folder in \textbf{Library}.}
\item
  \textbf{Drag a PDF from Finder into the item area.} Drop it on the folder or in its contents.
\item
  \textbf{Wait for the item to appear.} The app copies the file into managed storage and begins local identification.
\item
  \textbf{Review Paper Details.} Confirm the title, DOI, arXiv identifier, and file classification.
\end{enumerate}

\subsection{Add a paper by DOI, arXiv ID, URL, or title}\label{add-a-paper-by-doi-arxiv-id-url-or-title}

\begin{enumerate}
\tightlist
\item
  \textbf{Open the destination folder.} The resolved paper will be placed in the active Library folder.
\item
  \textbf{Choose \textbf{Add Paper}.} Control-click blank space in the folder contents and choose the command to open the inline lookup field.
\item
  \textbf{Enter a DOI, arXiv identifier, supported article URL, or title.} A DOI or arXiv identifier is the strongest identity clue. A title search may return alternatives that need review.
\item
  \textbf{Review and confirm the resolved match.} Check the title, authors, journal, and year, then choose \textbf{Add Paper} again.
\end{enumerate}

If the identity already exists, Gewu Insight reuses the canonical paper and adds the selected folder placement rather than creating an unrelated duplicate.

\begin{quote}
\textbf{Note}

A paper record does not require a local PDF. \textbf{Add Paper} creates or reuses the metadata record; it does not automatically download a PDF. Use \textbf{Download} or \textbf{Add File} separately when you want a managed local file.
\end{quote}

\begin{quote}
\textbf{Note}

A file added to an existing paper's \textbf{Paper Files} bundle is an attachment or version of that paper. A normal folder drop creates or reuses a paper placement in that folder. When identity matching finds an existing paper, the incoming file joins its Paper Files; different PDFs are retained even when they share a DOI.
\end{quote}

Folder drops and Folder Loader use the same import rules. An identity mismatch keeps the file with the matched paper, preserves its existing metadata and Default, and records the filename and identifiers in \textbf{Details}. It does not pause for a review decision. Import status stays in the bottom bar until dismissed; failed or stopped imports also refresh Library so completed items remain visible. A temporary progress connection failure triggers reconnection rather than reporting that the import stopped.

\begin{GewuRelated}
\textbf{See also:} \hyperref[attachments]{Add and Classify Attachments} · \hyperref[folder-loader]{Use Folder Loader} · \hyperref[drag-attach]{Move a Standalone File into Paper Files}
\end{GewuRelated}

\section{Use Paper Details}\label{paper-details}

\begin{GewuLead}
Paper Details combines actions, identifiers, metadata, and citation counts. Library, My Publications, and Author Monitor start with the paper title, without a separate heading bar. Action labels stay on one line; whole buttons wrap when space is narrow.
\end{GewuLead}

\subsection{Top actions}\label{top-actions}

{\def\LTcaptype{none} % do not increment counter
\GewuTableStart
\begin{longtable}[]{@{}>{\RaggedRight\arraybackslash}p{0.27\textwidth}>{\RaggedRight\arraybackslash}p{0.63\textwidth}@{}}
\toprule\noalign{}
\rowcolor{GewuTableHeader}
Action & What it does \\
\midrule\noalign{}
\endhead
\bottomrule\noalign{}
\endlastfoot
\textbf{Open in Preview} & Opens the preferred local PDF in macOS Preview. \\
\textbf{Download} & Retrieves the best available paper file without overwriting a different preserved version. \\
\textbf{Open Website} & Opens the publisher, DOI, arXiv, or best known work page using \hyperref[website-link-settings]{Website Links}: Gewu Insight Browser by default, or the system default browser. \\
\textbf{Google Scholar} & Searches for the work in Google Scholar. \\
\textbf{Get Identifier from Files} & Examines local paper files for DOI or arXiv evidence. \\
\end{longtable}
\GewuTableEnd
}

\subsection{File management}\label{file-management}

Library uses the adjacent \textbf{Paper Files} column; Paper Details does not repeat a Files box. My Publications keeps its collapsible \textbf{Files · n} management panel. Author Monitor shows a compact local filename with the full path on hover; \textbf{Get Identifier from Files} appears when a local PDF is available.

\subsection{Metadata rows}\label{metadata-rows}

The title appears once at the top. The rows below include authors, journal, date, DOI, arXiv, Gewu Key, citation, fields, and folder membership when available. \textbf{Gewu Key} is the stable, human-visible storage handle for a cataloged work record. It is not a DOI, arXiv identifier, or the database\textquotesingle s relational primary key.

\subsection{Compare citation sources}\label{compare-citation-sources}

\textbf{Merged Citations} brings together the counts returned by OpenAlex, Semantic Scholar, Crossref, NASA ADS, and saved Google Scholar data when those services have a match. Select a source label to show its saved count in the headline; this does not start a refresh. The merged estimate compares published, preprint, and already-merged records rather than simply adding every provider\textquotesingle s total. Separate published/preprint estimates can be combined, so this is not a verified count of unique citing papers. The yearly histogram uses OpenAlex data and does not switch provider when you select another count. PubMed Details supplies identifiers and indexing information, not another citation-count source.

\subsection{Inspect influential citations}\label{inspect-influential-citations}

When Semantic Scholar reports an \textbf{Influential} count, select it to inspect citing papers that Semantic Scholar classifies as especially meaningful to the cited work. This is Semantic Scholar's classification, not a score calculated by Gewu Insight. The list can include citation context, intent, citing authors, venue, year, and citation counts when the service supplies them.

Influential Citations is fetched on demand and cached locally for several hours. A saved API key gives Gewu Insight an individual Semantic Scholar allowance, but it is still limited to about one request per second and can be temporarily throttled. If the service returns a rate-limit message, the app pauses repeat requests for a few minutes; the paper record, existing citation totals, and any previously displayed result remain unaffected.

\subsection{Read and reuse the record}\label{read-and-reuse-the-record}

\begin{itemize}
\tightlist
\item
  A representative first figure, when available, provides a visual checkpoint above the actions.
\item
  \textbf{Abstract} keeps the paper summary in the independently scrolling details pane.
\item
  \textbf{BibTeX} and \textbf{Copy Code} provide a citation record for reference software or LaTeX.
\item
  \textbf{Nature Reference} and \textbf{Copy Reference} provide a formatted reference string.
\item
  \textbf{Folders} lists every Library location for the shared work. Remove a location there only when you intend to remove that reference.
\end{itemize}

\begin{quote}
\textbf{Note}

Gewu Key is assigned when a work is first cataloged in SQLite, before Add to Library. Its row remains visible in Author Monitor, including while details load; a missing value shows \texttt{(undefined)}. Adding to Library reuses the existing key and updates the view without creating a second identity.
\end{quote}

\begin{GewuRelated}
\textbf{See also:} \hyperref[identifiers]{Edit DOI and arXiv Identifiers} · \hyperref[paper-files]{Understand Paper Files} · \hyperref[storage-model]{How Gewu Insight Stores Files}
\end{GewuRelated}

\section{Edit DOI and arXiv Identifiers}\label{identifiers}

\begin{GewuLead}
DOI and arXiv can identify the same work across publication stages. Correct either value in Paper Details before downloading or synchronizing versions.
\end{GewuLead}

\subsection{Refresh or correct a DOI}\label{refresh-or-correct-a-doi}

\begin{enumerate}
\tightlist
\item
  \textbf{Select the paper.} Find the \textbf{DOI} row in Paper Details.
\item
  \textbf{Select the DOI field.} Enter or paste the correct DOI. You can omit a leading \texttt{https://doi.org/}.
\item
  \textbf{Choose \textbf{Refresh} or press \texttt{Return}.} Gewu Insight saves the identifier, checks supported metadata services, and updates matching bibliographic fields.
\item
  \textbf{Inspect the result.} Verify the title, journal, year, and authors before downloading a new file.
\end{enumerate}

\subsection{Edit an arXiv identifier}\label{edit-an-arxiv-identifier}

\begin{enumerate}
\tightlist
\item
  \textbf{Select the arXiv value.} The value changes to an inline text field.
\item
  \textbf{Enter an identifier.} Use forms such as \texttt{arXiv:2106.14849} or a valid legacy identifier.
\item
  \textbf{Press \texttt{Return}.} The database saves the new value. Enter an empty value and press \texttt{Return} to clear it.
\end{enumerate}

The \textbf{arXiv} row remains visible when empty so you can add a corrected identifier later. Use the arrow icon to open arXiv and the download icon to retrieve the latest available arXiv PDF.

\begin{quote}
\textbf{Important}

Confirm an identifier before running a version sync. A wrong identifier can attach files from a different work.
\end{quote}

\begin{GewuRelated}
\textbf{See also:} \hyperref[download-sync]{Download and Sync Versions} · \hyperref[citation-settings]{Open Source Settings}
\end{GewuRelated}

\section{Understand Paper Files}\label{paper-files}

\begin{GewuLead}
A paper is one work record that can own several PDFs and related files. Paper Files shows that bundle without pretending it is another Library folder level.
\end{GewuLead}

\subsection{Compact cards for one flat bundle}\label{compact-cards-for-one-flat-bundle}

The Paper Files column is a flat peer list, not a hierarchy of version or attachment folders. Each card shows its filename, \textbf{File type}, and size by default. Click the size footer or its chevron to expand \textbf{Annotation} and \textbf{Source} into separate rows; click again to collapse them. The chevron points down when collapsed and up when expanded. Each card expands independently, and the choice lasts for the current session. File type remains visible because it identifies a preprint, accepted manuscript, published version, supplement, or other resource. The \textbf{Default} marker identifies the preferred file.

\subsection{Select, open, or act on a file}\label{select-open-or-act-on-a-file}

\begin{itemize}
\tightlist
\item
  Single-click a card to select it.
\item
  Drag a card\textquotesingle s file icon or filename to Finder or the Desktop to copy it, or onto a compatible application icon in the Dock or on the Desktop to open the original file with that application. Library item file icons support the same native drag. The target application decides which file types it accepts.
\item
  Double-click a card to open it.
\item
  Double-click the paper row to the left to open the paper\textquotesingle s Default file, independently of the highlighted resource card. PDFs open in macOS Preview. If an old unclassified import has no Default, its available current file still opens.
\item
  Press \texttt{Space} to use Quick Look.
\item
  Control-click a card for \textbf{Open}, \textbf{Quick Look}, \textbf{Show in Finder}, default-file controls, \textbf{Change File Type}, or \textbf{Delete from Paper Files} where allowed. Control-click blank space in Paper Files for \textbf{Add File} and, when the paper has an arXiv identifier, \textbf{Sync Versions}. Those bundle-wide commands do not appear on an individual file card.
\item
  Press \texttt{Delete} or \texttt{Backspace} while Paper Files has keyboard focus to remove only the selected removable resource after confirmation. This does not delete the parent paper or its other files.
\end{itemize}

The physical primary file used by older compatibility paths remains protected from direct deletion. Choosing \textbf{Set as Default} makes that resource the file used by \textbf{Open in Preview} and paper preview without moving, renaming, or rewriting it. If you set it manually, \textbf{Use Automatic Selection} returns control to the file-type priority rule.

Library Paper Details no longer repeats this bundle. In My Publications, which has no adjacent Paper Files column, the collapsible \textbf{Files} panel keeps its file-management controls. \textbf{Show in Finder} reveals a physical file.

Deleting a private managed attachment moves its bytes to macOS Trash before removing its resource record. Shared files and external originals are preserved. A failed database deletion restores the moved file when macOS returns its Trash path. The interface shows progress and updates only after success; path indexes keep the shared-file safety check fast even in large libraries.

\subsection{Correct a file\textquotesingle s paper assignment}\label{correct-a-files-paper-assignment}

If an imported PDF was matched to the wrong paper, Control-click its card and choose \textbf{Move to Another Paper\ldots{}}. Search the destination in Library by title, DOI, or Gewu Key, select it, and choose \textbf{Move File}. To make it independent, choose \textbf{Create Paper from File\ldots{}}, enter a title, and choose \textbf{Create Paper}. The new paper appears in the current Library folder with a new Gewu Key and no guessed DOI.

Only the selected file changes ownership. Its bytes, filename, file ID, classification and source information remain intact; the existing paper\textquotesingle s title and DOI stay unchanged. If you move its Default or primary file, the source falls back to a remaining eligible file. The new independent paper uses the selected file as its Default. A repeat import of the same source path and size respects the explicit reassignment. Historical analysis text is not rewritten, so check any notes that refer to the old grouping.

These correction commands are user interface actions. AI Access cannot move, edit, overwrite or delete Paper Files; it can only read authorized files and append new Markdown analyses.

\subsection{Understand a Missing File row}\label{understand-a-missing-file-row}

\textbf{Missing File} means that SQLite still has a resource record, but the file is not present at the recorded local path. Gewu Insight keeps genuine missing resources visible so you can identify and repair them; it hides only an older untyped migration alias when it can prove that the alias has already been replaced by the same paper\textquotesingle s live canonical file.

A missing attachment is not another copy of the default PDF merely because both rows have a similar name. Check \textbf{Show in Finder}, the paper\textquotesingle s Associated Files bundle, and Recycle before deleting or importing anything. The default resource and explicitly typed attachments remain visible even when their paths need repair.

\begin{GewuRelated}
\textbf{See also:} \hyperref[quick-look]{Preview Files with Quick Look} · \hyperref[versions-current]{Choose File Types and the Default File} · \hyperref[attachments]{Add and Classify Attachments} · \hyperref[navigation]{Navigate with Mouse and Keyboard}
\end{GewuRelated}

\section{Preview Files with Quick Look}\label{quick-look}

\begin{GewuLead}
Press Space to inspect the selected resource without changing its paper record. Quick Look can stay inside Gewu Insight, use the macOS system panel, or combine both approaches.
\end{GewuLead}

\subsection{Open and follow a preview}\label{open-and-follow-a-preview}

\begin{enumerate}
\tightlist
\item
  \textbf{Select a file in Library, a Space, or Paper Files.} Press \texttt{Space} to open the configured Quick Look.
\item
  \textbf{Move with \texttt{{\symbolfont ↑}} or \texttt{{\symbolfont ↓}}.} When Quick Look is open, the preview follows the current Gewu Insight selection.
\item
  \textbf{Press \texttt{Space} again.} Close Quick Look and return keyboard focus to the same file.
\end{enumerate}

With macOS Quick Look, Gewu Insight keeps keyboard focus so you can continue moving through Library or Paper Files while the system panel follows the selection. Rapid navigation is serialized so an older preview request cannot replace a newer selection.

\par\noindent\begin{minipage}{\linewidth}

\begin{center}\includegraphics[width=0.96\linewidth,height=0.72\textheight,keepaspectratio]{/Users/lingyuan/PaperMonitor-20260516/assets/user-guide/screenshots/pdf-quick-look.png}\end{center}

\begin{GewuCaption}
Built-in PDF Quick Look. The selected Library PDF appears with page thumbnails and Open in Preview. The document and thumbnail column scroll independently.
\end{GewuCaption}

\end{minipage}\par\medskip

\subsection{Read the built-in PDF loading sequence}\label{read-the-built-in-pdf-loading-sequence}

The document viewer starts first and requests the bytes it needs from the local file instead of waiting for the whole PDF to be copied into memory. Page count and other metadata load in parallel. Page thumbnails are scheduled after the embedded document reports loading, so the document and thumbnail strip may appear at different times. That loading signal is not a guarantee that every PDF page has finished painting.

On the first open, the app may still need to inspect the PDF and render thumbnail pages. Later opens can reuse stable local-file validators, remembered page information, and a bounded thumbnail cache. The cache is automatically trimmed and can always be rebuilt; it is not part of the paper record or a substitute for the PDF.

\subsection{Read and navigate the built-in viewer}\label{read-and-navigate-the-built-in-viewer}

The PDF opens at page width. Scroll the document, click a page thumbnail to jump, or use a trackpad pinch to zoom the PDF. The document and thumbnail strip scroll independently. Drag the preview title bar to reposition it, or its resize corner to change the reading area. The strip provides page buttons for the first 200 pages; thumbnail images are limited to the first 160. The document remains scrollable beyond that range; use \textbf{Open in Preview} for the full reader\textquotesingle s navigation and tools.

Markdown files (\texttt{.md}, \texttt{.markdown}, \texttt{.mdown}, and \texttt{.mkd}) use a formatted reading view in built-in Quick Look, including headings, emphasis, nested lists, quotations, code blocks, and tables. The view follows keyboard selection, including Markdown attachments in Paper Files. Nearby images in the same folder or its subfolders can appear inline; remote images are links that you open explicitly. Website links follow Website Links settings. Raw HTML is shown as text. Large files are limited to a 256 KiB preview and marked {[}Preview truncated{]}; the original file is unchanged. Complex notes may use a shorter preview to keep navigation responsive. Native Quick Look modes keep the system\textquotesingle s own renderer.

Supported images, audio, video, and text files can also be inspected inside the workspace. A long text preview may end with {[}Preview truncated{]}; use Open File to read the complete original. Other formats show a file card with \textbf{Open File}; that opens the file in its macOS application. Choose a native mode in \hyperref[quick-look-settings]{Quick Look Settings} for formats better handled by system Quick Look.

Switching between PDFs in the built-in viewer can briefly flash while the next document loads. It does not guarantee the same switching speed or continuity as macOS Quick Look. For rapid comparison, try \textbf{macOS Quick Look for all files}.

\par\noindent\begin{minipage}{\linewidth}

\begin{center}\includegraphics[width=0.96\linewidth,height=0.72\textheight,keepaspectratio]{/Users/lingyuan/PaperMonitor-20260516/assets/user-guide/screenshots/markdown-quick-look.png}\end{center}

\begin{GewuCaption}
Built-in Markdown Quick Look. This existing README.md renders headings and links. Remote images remain explicit links; nearby local images can appear inline.
\end{GewuCaption}

\end{minipage}\par\medskip

\subsection{Choose Quick Look or Preview}\label{choose-quick-look-or-preview}

{\def\LTcaptype{none} % do not increment counter
\GewuTableStart
\begin{longtable}[]{@{}>{\RaggedRight\arraybackslash}p{0.27\textwidth}>{\RaggedRight\arraybackslash}p{0.63\textwidth}@{}}
\toprule\noalign{}
\rowcolor{GewuTableHeader}
Use & Best for \\
\midrule\noalign{}
\endhead
\bottomrule\noalign{}
\endlastfoot
\textbf{Built-in Quick Look} & PDF inspection and keyboard comparison among versions while remaining inside the workspace. \\
\textbf{macOS Quick Look} & System-native preview support for presentations, media, images, PDFs, and other formats supported by macOS. \\
\textbf{Open in Preview} & Full macOS Preview tools, annotation, export, printing, and a separate reading window. \\
\end{longtable}
\GewuTableEnd
}

\begin{quote}
\textbf{Note}

Previewing is read-only from Gewu Insight\textquotesingle s perspective. Opening Quick Look, loading thumbnails, or moving between pages does not change file type, annotation state, source, Default status, or the SQLite paper record.
\end{quote}

\begin{GewuRelated}
\textbf{See also:} \hyperref[quick-look-settings]{Quick Look Settings} · \hyperref[navigation]{Navigate with Mouse and Keyboard} · \hyperref[paper-files]{Understand Paper Files} · \hyperref[troubleshooting]{Troubleshooting}
\end{GewuRelated}

\section{Choose File Types and the Default File}\label{versions-current}

\begin{GewuLead}
The default file is the preferred local representation of the paper. Choose it manually or let Gewu Insight select it. Default appears first in the file list; Up and Down follow the visible order.
\end{GewuLead}

\subsection{Automatic selection}\label{automatic-selection}

The automatic order is \textbf{Corrected published version}, \textbf{Published version}, \textbf{Accepted manuscript}, then \textbf{Preprint}. Within a type, a later recognized version normally wins. Correction notices, supplementary material, peer review files, Other, and Invalid file never become the automatic default. When a paper has no Default or main PDF, its first article PDF saved through Browser becomes a fallback Default even if File type is Unknown. It enters the main file folder immediately; later Unknown copies do not replace it.

\subsection{Automatic and manual eligibility}\label{automatic-and-manual-eligibility}

The four publication types above take priority in automatic selection, ahead of the first Browser PDF fallback. You may manually choose any existing file except one classified as \textbf{Invalid file}. A manual choice is sticky across rescans, metadata refreshes, and later downloads until you choose \textbf{Use Automatic Selection}. Newly arriving files remain peer rows and do not erase or silently reclassify older files.

\subsection{Choose a file manually}\label{choose-a-file-manually}

\begin{enumerate}
\tightlist
\item
  \textbf{Enter Paper Files, or expand \textbf{Files} in My Publications.}
\item
  \textbf{Select the file you want.}
\item
  \textbf{Choose \textbf{Set as Default}.} This file will be used by \textbf{Open in Preview} and paper preview. You can also pin the automatically selected Default this way.
\end{enumerate}

\begin{GewuResult}
\textbf{Result.} The selected file becomes the manual default. When it is a PDF, Gewu Insight places it in the main managed file folder and keeps the previous primary PDF in Associated Files. The article and file identities stay unchanged. External originals and files shared with another article are copied safely; existing filenames are retained unless a name conflict requires a unique name. Non-PDF defaults keep their existing location.
\end{GewuResult}

\subsection{Return to automatic selection}\label{return-to-automatic-selection}

Choose \textbf{Use Automatic Selection} for a manually selected default. Gewu Insight evaluates the eligible file types again.

\subsection{Classify files without collapsing annotation state}\label{classify-files-without-collapsing-annotation-state}

In Library, File type stays visible on each card. Click the size footer to reveal Annotation and Source, then change each control independently. In My Publications, expand Files to see all three controls directly. File type uses exactly these labels:

\begin{itemize}
\tightlist
\item
  \textbf{Preprint}, \textbf{Accepted manuscript}, \textbf{Published version}, and \textbf{Corrected published version}
\item
  \textbf{Correction notice}, \textbf{Supplementary material}, and \textbf{Peer review file}
\item
  \textbf{Other}, \textbf{Unknown}, and \textbf{Invalid file}
\end{itemize}

Annotation state is selected independently as \textbf{Annotated}, \textbf{Unannotated}, or \textbf{Unknown}. Source is independently recorded as \textbf{Publisher}, \textbf{Preprint server}, \textbf{Institutional repository}, \textbf{User import}, \textbf{Other}, or \textbf{Unknown}. Manual type and annotation choices survive later fingerprint rescans.

\textbf{Publisher} is used only when a recognized publisher host or a publisher-specific acquisition path provides that evidence. arXiv, bioRxiv, and medRxiv are shown as \textbf{Preprint server}; repository hosts remain \textbf{Institutional repository}; ordinary local additions are \textbf{User import}; and a generic web download with no stronger evidence remains \textbf{Other}. A real host can appear after the source label when useful. Internal migration or ingestion markers are never user-facing source details, and choosing a different source clears a stale provider detail rather than appending it to the new choice.

\begin{quote}
\textbf{Note}

Changing File type can change the automatic Default. When the selected Default is a PDF, its main file location is synchronized too. Annotation and Source remain independent properties; changing a label does not apply a filename template or rewrite the PDF.
\end{quote}

\begin{GewuRelated}
\textbf{See also:} \hyperref[download-sync]{Download and Sync Versions} · \hyperref[attachments]{Add and Classify Attachments}
\end{GewuRelated}

\section{Add and Classify Attachments}\label{attachments}

\begin{GewuLead}
Import first, classify afterward. This keeps the add operation simple and lets you set type, annotation state, and source after seeing the file in context.
\end{GewuLead}

\subsection{Add a file to the paper}\label{add-a-file-to-the-paper}

\begin{enumerate}
\tightlist
\item
  \textbf{Control-click blank space in the Paper Files area and choose \textbf{Add File}.} In My Publications, expand Files and use Add File.
\item
  \textbf{Select one or more local files.} Confirm the file chooser.
\item
  \textbf{Wait for the new row.} A new item begins as \textbf{Unknown} unless its source proves a more specific type.
\item
  \textbf{Select its compact controls.} Choose File type; in Library, click the size footer to expand Annotation and Source. My Publications shows those controls directly in its Files panel.
\end{enumerate}

\subsection{Attachment types}\label{attachment-types}

The same ten file types apply to every row. File extension and MIME type remain separate technical facts, so an image, note, dataset, or code archive can use \textbf{Other} without being mistaken for a PDF version.

\subsection{Add by drag and drop}\label{add-by-drag-and-drop}

Drag a Finder file over the expanded Paper Files area, wait for the drop target to respond, then release. The file is copied into the selected paper's associated bundle and appears in Paper Files.

\begin{quote}
\textbf{Important}

Invalid or unreadable PDFs are rejected at file ingress. Existing migrated records that can be proven invalid remain visible as \textbf{Invalid file} for audit and repair, but cannot become a default.
\end{quote}

\begin{GewuRelated}
\textbf{See also:} \hyperref[drag-attach]{Move a Standalone File into Paper Files} · \hyperref[storage-model]{How Gewu Insight Stores Files}
\end{GewuRelated}

\section{Download and Sync Versions}\label{download-sync}

\begin{GewuLead}
Downloads add a version to the paper bundle. They do not silently overwrite a different preserved PDF.
\end{GewuLead}

\subsection{Download the best available paper}\label{download-the-best-available-paper}

\begin{enumerate}
\tightlist
\item
  \textbf{Verify DOI and arXiv.} Make sure the identifiers belong to the selected work.
\item
  \textbf{Choose \textbf{Download} at the top of Paper Details.} Gewu Insight first tries the best available published file, then supported alternatives.
\item
  \textbf{Wait for the status message.} Existing files are preserved if the incoming file is a different version; a manual default stays selected.
\item
  \textbf{Inspect Paper Files.} Every preserved file appears in the same flat list. Automatic ranking changes only the Default marker; a manual default remains unchanged.
\end{enumerate}

\subsection{Download arXiv files}\label{download-arxiv-files}

\begin{itemize}
\tightlist
\item
  Use the download icon in the \textbf{arXiv} row to retrieve the latest arXiv PDF.
\item
  Control-click blank space in Paper Files and choose \textbf{Sync Versions} to retrieve all available arXiv revisions such as v1, v2, and v3. In My Publications, use Sync Versions inside the expanded Files panel. The command is available when the paper has an arXiv identifier.
\item
  Each revision is labeled with its arXiv version when known.
\end{itemize}

\subsection{Duplicate protection}\label{duplicate-protection}

Gewu Insight computes exact SHA-256 and PDF content fingerprints as each new resource is registered. An exact duplicate is reused instead of stored twice. A clean PDF and a genuinely annotated copy can share a content fingerprint while retaining different exact hashes, annotation counts, and labels.

\subsection{Validation before replacement}\label{validation-before-replacement}

A response is not accepted merely because it begins with \texttt{\%PDF-}. Structural PDF validation rejects empty files, HTML responses, truncated downloads, missing valid endings, and content that cannot be parsed as a PDF. A download is written to a temporary file and replaces the destination only after validation succeeds. If the request or stream fails, an existing valid file remains untouched and a partial file is not registered as a resource.

\subsection{Classification follows the file that actually arrived}\label{classification-follows-the-file-that-actually-arrived}

A publisher download is classified from the file and source actually returned. A corrected published version outranks a published version automatically; a publisher reference, First Release, or accepted preview remains an accepted manuscript. A correction notice remains a separate peer and never becomes automatic merely because it came from a publisher.

A PDF URL stored in bibliographic metadata is only a retrieval lead. It does not create a Paper Files row until a download actually succeeds, the returned bytes pass validation, and a local resource is registered.

If a download is skipped because the local bytes already exist, that skipped operation does not relabel the existing file from metadata or a newly discovered URL. Correct the file type explicitly if the saved classification itself is wrong.

\begin{quote}
\textbf{Note}

A Monitor result is not forced into a Library placement. Where download is available, a local resource can belong to the cataloged paper even before you add that paper to Library or My Publications.
\end{quote}

\begin{GewuRelated}
\textbf{See also:} \hyperref[versions-current]{Choose File Types and the Default File} · \hyperref[network-access]{Publisher and Network Access} · \hyperref[identifiers]{Edit DOI and arXiv Identifiers}
\end{GewuRelated}

\section{Move a Standalone File into Paper Files}\label{drag-attach}

\begin{GewuLead}
If a data, image, presentation, or other file was imported as its own Library item, drag it onto the correct paper's Paper Files area to make it an attachment.
\end{GewuLead}

\subsection{Convert an independent item}\label{convert-an-independent-item}

\begin{enumerate}
\tightlist
\item
  \textbf{Select the destination paper.} Keep its Paper Files inspector visible.
\item
  \textbf{Begin dragging the standalone item.} The main paper remains the active destination while you drag.
\item
  \textbf{Move the pointer into the Paper Files drop zone.} Wait for the drop feedback that identifies the target paper.
\item
  \textbf{Release the pointer.} Gewu Insight validates the move, adopts the file into the paper bundle, and removes its independent Library position.
\item
  \textbf{Classify the attachment.} Choose \textbf{Supplementary material} for an actual paper supplement, or \textbf{Other} for a general image, dataset, code file, slide deck, or note. Expand the size footer to set Annotation and Source separately. See \hyperref[versions-current]{the complete current file-type list}.
\end{enumerate}

\begin{GewuResult}
\textbf{Result.} The physical file and database link now belong to the destination paper; the former standalone entry is gone.
\end{GewuResult}

\subsection{If the conversion cannot finish}\label{if-the-conversion-cannot-finish}

The file moves and database changes are coordinated as one recoverable operation. If any step fails, Gewu Insight restores the original standalone files and rolls back the paper, Library item, resource relationships, and fingerprint fields. The source should remain usable rather than becoming a half-converted attachment.

\begin{quote}
\textbf{Important}

If the source item is shared by another Library reference or protected relationship, Gewu Insight may refuse a destructive conversion. Add a copy with \textbf{Add File} instead.
\end{quote}

\begin{GewuRelated}
\textbf{See also:} \hyperref[attachments]{Add and Classify Attachments} · \hyperref[organize]{Folders, Shortcuts, and Spaces}
\end{GewuRelated}

\section{Delete and Recover Items}\label{recycle}

\begin{GewuLead}
Deleting a paper reference either removes only that reference or moves the last Library reference to Recycle. Final deletion inside Recycle is reference-aware and uses the macOS Trash for managed files.
\end{GewuLead}

\subsection{Delete a Library reference}\label{delete-a-library-reference}

\begin{enumerate}
\tightlist
\item
  \textbf{Control-click the item.}
\item
  \textbf{Choose \textbf{Delete}.}
\item
  \textbf{Check the result.} If another Library location or My Publications still refers to the same work, only the selected reference is removed. Otherwise the item moves to \textbf{Recycle} and remains recoverable.
\end{enumerate}

\subsection{Delete from My Publications}\label{delete-from-my-publications}

Select one or more publications and choose \textbf{Delete}. If a Library, Space, Shortcut or Recycle location still refers to the work, its shared article record, Gewu Key and files remain. Otherwise, deleting this last placement also removes the article record and its identifier tasks from SQLite; private managed files move to macOS Trash. External originals and bytes shared with other articles are preserved. Import history does not keep an otherwise unreferenced article alive.

Removing every currently saved article of a running profile import pauses that import. Continue Import does not restore previously removed source entries. An explicit Add All Publications recheck after completion can add those articles again.

\subsection{Recover an item}\label{recover-an-item}

\begin{enumerate}
\tightlist
\item
  \textbf{Open \textbf{Recycle}.}
\item
  \textbf{Control-click the item and choose \textbf{Recover}.}
\item
  \textbf{Confirm the restored location.}
\end{enumerate}

\subsection{Delete permanently}\label{delete-permanently}

Final deletion of Library items takes place inside Recycle. Choosing \textbf{Delete} on a folder moves that folder intact, including its subfolders and unregistered contents, to macOS Trash. The confirmation lists the folders and files to move. Original imported source files outside Gewu Insight are not touched. If another Library, Space, Recycle or My Publications location still refers to the paper, its shared record, Gewu Key, Paper Files and Default remain; any required local files are retained outside the folder being removed. Only papers with no remaining references are removed from SQLite.

\begin{quote}
\textbf{Empty Recycle}

Control-click the Recycle root folder or blank space directly inside it to choose \textbf{Empty Recycle}. This clears every Recycle entry and moves all contents of its local disk directory to macOS Trash, including unregistered files, empty folders and broken shortcuts. The Recycle root remains. This command is not offered on files or folders inside Recycle.
\end{quote}

\begin{GewuRelated}
\textbf{See also:} \hyperref[backup-migrate]{Back Up or Move to Another Mac} · \hyperref[storage-model]{How Gewu Insight Stores Files}
\end{GewuRelated}

\chapter{Import and Transfer}\label{chapter-13}

\section{Use Folder Loader}\label{folder-loader}

\begin{GewuLead}
Folder Loader imports an existing directory tree into managed Library storage while preserving relative paths.
\end{GewuLead}

\par\noindent\begin{minipage}{\linewidth}

\begin{center}\includegraphics[width=0.96\linewidth,height=0.72\textheight,keepaspectratio]{/Users/lingyuan/PaperMonitor-20260516/assets/user-guide/screenshots/folder-loader.png}\end{center}

\begin{GewuCaption}
Folder Loader. Choose a source directory and metadata mode before Import. The live log and Result Guide fill the remaining window height and scroll independently.
\end{GewuCaption}

\end{minipage}\par\medskip

\subsection{Import a folder}\label{import-a-folder}

\begin{enumerate}
\tightlist
\item
  \textbf{Choose \textbf{Loader \textgreater{} Open Loader}.} Outside Browser, you can also press \texttt{{\symbolfont ⌘}}\texttt{L}; inside Browser that shortcut focuses its address field.
\item
  \textbf{Choose the source folder.} Review the displayed path before starting.
\item
  \textbf{Select a metadata mode.}

  \begin{itemize}
  \tightlist
  \item
    \textbf{Skip metadata} --- fastest; import files and structure only.
  \item
    \textbf{Quick refresh} --- recommended for most migrations.
  \item
    \textbf{Full refresh} --- more network work and a longer run.
  \end{itemize}
\item
  \textbf{Choose \textbf{Import}.} Follow progress in the live log. Use \textbf{Stop Import} if necessary.
\item
  \textbf{Review the result.} Open previous logs or the result guide when items need attention.
\end{enumerate}

\subsection{What Folder Loader preserves}\label{what-folder-loader-preserves}

\begin{itemize}
\tightlist
\item
  The source folder's relative hierarchy.
\item
  Existing managed targets and files on repeated imports.
\item
  An existing primary PDF rather than overwriting it.
\item
  A verified exported Paper Files relationship when a Gewu Insight manifest is present.
\end{itemize}

Folder Loader copies into Gewu Insight's managed storage. It does not move, rename, or delete the external source files and folders you selected.

\subsection{Identity mismatches and partial imports}\label{identity-mismatches-and-partial-imports}

Loader and Finder folder drops share the same matching and file-preservation policy. When a DOI, title or existing identifier disagrees with the matched paper, the incoming file is still saved in its Paper Files and the paper is placed in the target folder. Existing metadata and Default are preserved. Different PDFs with the same DOI are kept as separate files; the import continues without a mandatory review prompt.

The live log and \textbf{Result Guide} identify each mismatch by filename, detected DOI, existing DOI and Gewu Key. Correct a wrong assignment later from the file card\textquotesingle s \textbf{Move to Another Paper\ldots{}} or \textbf{Create Paper from File\ldots{}} command. This preserves the PDF; it does not certify that the identified DOI is correct.

A disk or other import error can still stop a batch. Completed items remain available, and the result shows the failing file and actual progress. For a background drop, use the bottom bar\textquotesingle s \textbf{Details} to reopen its current result and \textbf{Dismiss} when finished. Previously saved results remain available in Loader logs.

\subsection{Delete an import batch}\label{delete-an-import-batch}

\textbf{Delete Import} is a rollback for the selected Loader batch, not a way to forget a log. Stop an active import first. The command removes database records created by that batch, eligible empty folders, its Loader reports, and managed local copies associated with the batch; eligible paths are sent to the macOS Trash. It never deletes the original external source directory.

\begin{quote}
\textbf{Important}

Review the selected batch and make a backup before using Delete Import, especially if you have already reorganized or reused imported items.
\end{quote}

\begin{quote}
\textbf{Note}

For an ordinary folder, Loader does not guess which files are ``primary'' versus ``associated.'' It imports the collection as supplied. Special reconstruction occurs only for a verified Gewu Insight export manifest.
\end{quote}

\begin{GewuRelated}
\textbf{See also:} \hyperref[desktop-copy]{Make a Portable Desktop Copy} · \hyperref[loader-maintenance]{Refresh an Imported Library} · \hyperref[backup-migrate]{Back Up or Move to Another Mac}
\end{GewuRelated}

\section{Refresh an Imported Library}\label{loader-maintenance}

\begin{GewuLead}
Loader maintenance commands fill metadata in batches after the files and folder structure are safely imported.
\end{GewuLead}

\par\noindent\begin{minipage}{\linewidth}

\begin{center}\includegraphics[width=0.96\linewidth,height=0.72\textheight,keepaspectratio]{/Users/lingyuan/PaperMonitor-20260516/assets/user-guide/screenshots/loader-metadata.png}\end{center}

\begin{GewuCaption}
Refresh All Metadata. The maintenance window shows its Library scope, metadata mode and skip rule before Start Refresh.
\end{GewuCaption}

\end{minipage}\par\medskip

\par\noindent\begin{minipage}{\linewidth}

\begin{center}\includegraphics[width=0.96\linewidth,height=0.72\textheight,keepaspectratio]{/Users/lingyuan/PaperMonitor-20260516/assets/user-guide/screenshots/loader-identifiers.png}\end{center}

\begin{GewuCaption}
Find All Missing Identifiers. Start Search checks eligible Library PDFs for missing identifiers. The right pane explains the source policy and retry cooldowns.
\end{GewuCaption}

\end{minipage}\par\medskip

\subsection{Refresh known works}\label{refresh-known-works}

\begin{enumerate}
\tightlist
\item
  \textbf{Choose \textbf{Loader \textgreater{} Refresh All Metadata}.}
\item
  \textbf{Choose \textbf{Start Refresh}.} Eligible Library papers with DOI or arXiv identifiers are refreshed sequentially.
\item
  \textbf{Review the log.} Recently successful records may be skipped to avoid unnecessary requests.
\end{enumerate}

\subsection{Find missing identifiers}\label{find-missing-identifiers}

\begin{enumerate}
\tightlist
\item
  \textbf{Choose \textbf{Loader \textgreater{} Find All Missing Identifiers}.}
\item
  \textbf{Choose \textbf{Start Search}.} The app checks active Library and Space PDFs using Crossref, OpenAlex, and Semantic Scholar where available.
\item
  \textbf{Review uncertain matches.} Refresh or recover DOI evidence as needed, and edit an incorrect arXiv identifier inline before downloading versions.
\end{enumerate}

Verified ``not found'' checks and transient failures use cooldowns so a repeated run does not immediately send the same requests. Bulk missing-identifier work does not use Google Scholar.

\begin{GewuRelated}
\textbf{See also:} \hyperref[citation-settings]{Open Source Settings} · \hyperref[network-access]{Publisher and Network Access} · \hyperref[identifiers]{Edit DOI and arXiv Identifiers}
\end{GewuRelated}

\section{Make a Portable Desktop Copy}\label{desktop-copy}

\begin{GewuLead}
Make Copy Desktop creates a self-contained copy for transfer, inspection, or re-import.
\end{GewuLead}

\subsection{Create a copy}\label{create-a-copy}

\begin{enumerate}
\tightlist
\item
  \textbf{Control-click the non-root Library folder you want to export.}
\item
  \textbf{Choose \textbf{Make Copy Desktop}.}
\item
  \textbf{Wait for the completion message.} Open the resulting Desktop folder in Finder.
\item
  \textbf{Confirm representative files.} Check a paper with versions or attachments as well as a normal folder.
\end{enumerate}

\subsection{What is included}\label{what-is-included}

The export replaces managed filesystem shortcuts with real files, includes relevant Associated Files, and writes a verified manifest describing paper-file relationships. This lets Folder Loader restore the bundle on another installation without guessing from filenames.

\begin{quote}
\textbf{Important}

Do not edit or discard the export manifest before migration. Without it, Loader correctly treats the files as an ordinary folder and does not infer hidden associations.
\end{quote}

\begin{GewuRelated}
\textbf{See also:} \hyperref[folder-loader]{Use Folder Loader} · \hyperref[backup-migrate]{Back Up or Move to Another Mac} · \hyperref[storage-model]{How Gewu Insight Stores Files}
\end{GewuRelated}

\section{Back Up or Move to Another Mac}\label{backup-migrate}

\begin{GewuLead}
A reliable transfer preserves both the SQLite database and the managed file roots. The app bundle alone is not your library.
\end{GewuLead}

\subsection{Create a full backup}\label{create-a-full-backup}

\begin{enumerate}
\tightlist
\item
  \textbf{Quit Gewu Insight with {{\symbolfont ⌘}Q}.} Wait until background refresh or download work has stopped. Closing a window is not the same as quitting; also stop any separately launched local UI or CLI process before copying its database.
\item
  \textbf{Copy the complete project directory.} This is the safest recovery unit because it keeps \texttt{data}, \texttt{reports}, \texttt{config}, the app source, and managed paths together.
\item
  \textbf{Preserve directory names and permissions.} Do not separate SQLite from the managed file roots it describes.
\item
  \textbf{Store the backup separately.} Use a second disk or a versioned backup location.
\end{enumerate}

\subsection{Back Up and Upgrade}\label{back-up-and-upgrade}

\begin{enumerate}
\tightlist
\item
  \textbf{Quit the current version and make a full backup.} Preserve SQLite and every managed file root together.
\item
  \textbf{Install or build the newer application.} A normal software upgrade does not overwrite the existing \texttt{data} directory.
\item
  \textbf{Start the new version once.} Required database schema changes run automatically against the existing SQLite database.
\item
  \textbf{Verify before deleting the backup.} Open the default PDF, another Paper Files version, an attachment, and a My Publications entry.
\end{enumerate}

\begin{quote}
\textbf{Note}

Schema migration updates database structure; it is not a reason to rebuild the database or re-import the Library. Keep the pre-upgrade backup until the new version has been checked.
\end{quote}

\subsection{Move to another Mac}\label{move-to-another-mac}

\begin{enumerate}
\tightlist
\item
  \textbf{Install or update Gewu Insight on the destination Mac.}
\item
  \textbf{Keep the destination app closed.}
\item
  \textbf{Restore the matching database and managed file roots together.} Do not mix a new database with an unrelated old Associated Files tree.
\item
  \textbf{Start the app and test representative works.} Open a paper, an older version, and an attachment; verify My Publications and one monitor.
\end{enumerate}

\subsection{Transfer only part of a Library}\label{transfer-only-part-of-a-library}

Use \textbf{Make Copy Desktop} on the source Mac, move the complete exported folder, then use Folder Loader on the destination Mac. Keep its verified manifest intact.

\subsection{Browser and credential backup boundaries}\label{browser-and-credential-backup-boundaries}

The project database includes Browser bookmarks, folders, history, and saved sidebar preferences. A bookmark HTML export transfers only bookmarks, not history or research papers. Open webpage tabs and unsubmitted forms are memory-only and are not restored after quitting.

WebKit website cookies, sign-ins, and cache belong to a macOS-managed Browser profile outside this project. Copying the project or a bookmark export is not a supported transfer of that profile; expect to sign in again on another Mac. Apple Passwords and its passkeys remain under macOS/iCloud management and are never part of a Gewu export. Optional metadata API keys saved in SQLite do travel with the database: keep full backups private.

\begin{quote}
\textbf{Important}

Never migrate by copying only \texttt{paper\_monitor.sqlite} or only the visible Library folder. Either choice leaves half of the managed state behind.
\end{quote}

\begin{GewuRelated}
\textbf{See also:} \hyperref[desktop-copy]{Make a Portable Desktop Copy} · \hyperref[storage-reference]{Storage Reference} · \hyperref[troubleshooting]{Troubleshooting}
\end{GewuRelated}

\chapter{My Publications}\label{chapter-14}

\section{My Publications}\label{my-publications}

\begin{GewuLead}
My Publications keeps a managed publication list, citation sources, author-role statistics, and export tools for your own work.
\end{GewuLead}

\par\noindent\begin{minipage}{\linewidth}

\begin{center}\includegraphics[width=0.96\linewidth,height=0.72\textheight,keepaspectratio]{/Users/lingyuan/PaperMonitor-20260516/assets/user-guide/screenshots/my-publications.png}\end{center}

\begin{GewuCaption}
My Publications. The active profile, DOI-resolved publication list, citation evidence, article figure, and authorship summary share one view.
\end{GewuCaption}

\end{minipage}\par\medskip

\subsection{Choose your profile}\label{choose-your-profile}

\begin{enumerate}
\tightlist
\item
  \textbf{Enter a name, optional affiliations, or an exact Google Scholar profile URL.} For a name search, use \texttt{Name;\ affiliation\ A;\ affiliation\ B}.
\item
  \textbf{Choose \textbf{Find Profile}.} Review candidate profiles carefully.
\item
  \textbf{Choose \textbf{Use Profile} to review candidates individually, or \textbf{Add All Publications} to import the complete profile.} The saved strip shows the active name and profile photo. Select the photo to open the profile using \hyperref[website-link-settings]{Website Links}; use the remove control to forget it.
\end{enumerate}

Already have the paper in Library? Control-click its row and choose \textbf{Add to My Publications}. Multiple selected papers can be added together. The publication list and Library then refer to the same work, files, and Gewu Key.

\subsection{Add all publications or choose individual papers}\label{add-all-publications-or-choose-individual-papers}

\textbf{Use Profile} opens a candidate list where each article retains its own \textbf{Add} button. To import the author\textquotesingle s entire profile, choose \textbf{Add All Publications} beside the correct search profile or at the top of its article candidate list. These controls appear only inside the search/add area; the normal publication table keeps its original layout. The complete import reads successive source pages until the profile ends; it is independent of the 200-title candidate preview and has no 300- or 5,000-paper cutoff. Existing canonical articles and My Publications entries are reused.

The bottom status bar continuously shows active import counts and DOI/arXiv progress, including running, processed, paused, and retry-needed states. Completed imports move into \textbf{Import history}, which distinguishes original import counts from articles currently saved in My Publications. When only history remains, the footer entry reads \textbf{History}. Old imports do not appear in a new author search; its own status or error stays visible in the search panel. Choose \textbf{Details} in that bar to view full progress and import controls without searching again. The search/add panel opens when searching or choosing a profile and stays open until you choose its \textbf{Close} button; clicking a publication or completing a task does not dismiss it. Closing the panel keeps background work running and preserves the table position. \textbf{Pause Import} preserves completed work; \textbf{Continue Import} resumes it. An interrupted source page is retried without adding its successful articles twice. A blocked, unreadable, repeated, or interrupted source page is reported as incomplete with a short reason and guidance to try again later. Technical connection details stay in the diagnostic log. Individual identity conflicts do not stop later articles and remain reported for review and retry. Adding all again after completion explicitly checks the profile from its first page for additional articles.

Both individual Add and Add All automatically queue missing DOI/arXiv identifiers. Articles are saved first, and identifier completion runs separately in the background. Its progress distinguishes matched, unmatched, failed, and pending work. Switching tabs does not stop either task; unfinished work resumes after app restart. Identifier updates preserve the article\textquotesingle s Gewu Key and file relationships. The existing per-candidate \textbf{Retry ID} remains available. When saved articles have unmatched or failed identifiers, \textbf{Retry Unresolved IDs} appears with progress inside the search/add area. Pausing the profile import does not pause identifier work already queued. Articles without a confident match remain saved. This operation does not automatically download PDFs.

After Scholar reports an access restriction, Gewu Insight applies a shared 24-hour pause to Scholar requests. The prompt shows the local retry time; this is an application pause, not a Google-provided IP-ban deadline. Changing authors or restarting does not clear it.

\subsection{Add and update publications}\label{add-and-update-publications}

\begin{itemize}
\tightlist
\item
  Use the identifier/title field and \textbf{Add} to add a DOI, arXiv record, supported identifier, or title.
\item
  Choose \textbf{Refresh Metrics} to update citation services and PubMed data.
\item
  \textbf{First as listed} is calculated automatically from the active profile name and the resolved author order, including the profile\textquotesingle s English and Chinese name forms.
\item
  Use \textbf{Co-1st}, \textbf{Co-corr.}, and \textbf{Rev.} in Paper Details to mark co-first authorship, corresponding authorship, and review/perspective papers. Leave Co-1st unchecked for a solo first-author paper. These manual roles are preserved during metric refresh.
\item
  Use the star to highlight a paper. Star state is independent from citation refresh and export inclusion.
\item
  Drag rows into a manual order when a curated sequence matters more than date or citation sorting.
\item
  Choose \textbf{Comments} to build or open a citation-comments report from citing-paper metadata. This is not a personal note field.
\item
  Use \textbf{Hide} or \textbf{Show} in the context menu to control export inclusion. Use \textbf{Open in Finder} or \textbf{Delete} for the corresponding local action.
\end{itemize}

\subsection{Edit manual records and author marks}\label{edit-manual-records-and-author-marks}

A publication without either a DOI or arXiv identifier exposes \textbf{Manual Metadata}. Edit Title, Authors, Journal, Date, Volume, Issue, and Pages, then choose \textbf{Save Metadata}. Separate author names with semicolons or new lines. Once an identifier is present, use identifier-based metadata controls instead.

Select an author name in the publication\textquotesingle s author list to inspect affiliations and set \textbf{First author}, \textbf{Co-first author}, or \textbf{Corresponding author} for that named author. Checkbox changes save immediately. These individual author marks are separate from the publication\textquotesingle s Co-1st, Co-corr., and Rev. summary controls.

\subsection{Build a citation-comments report}\label{build-a-citation-comments-report}

\textbf{Comments} gathers citing-paper metadata and available Semantic Scholar citation-context sentences, then writes a local report that groups strong praise, positive impact statements, cautions, and unmatched citing records. It helps you review how later literature describes a publication; it does not modify the publication metadata or replace your own notes.

The report is a generated local reading artifact. Network services may return fewer contexts than citation records, and a temporary rate limit can leave some entries without quoted context. Reopen or rebuild the report later to retry the available sources; an existing report remains on disk until you deliberately replace or remove it.

\subsection{Work with files}\label{work-with-files}

My Publications uses the same Paper Details and flat Paper Files model as Library. Expand \textbf{Files · n} within the details to inspect its resources; Library instead has an adjacent Paper Files column. The publication Files panel shows File type, Annotation, and Source controls directly; the size-footer expansion belongs to Library cards. A publication can keep a default PDF, earlier versions, supplements, notes, presentations, data, and other attachments under the same Gewu Key.

\subsection{Read the summary}\label{read-the-summary}

The lower summary includes publication counts, identifier coverage, h-index and i10 information, journal distribution, merged citation maxima, and major-paper role counts. Select a role statistic to filter the table where supported.

\subsection{Export a list}\label{export-a-list}

Choose \textbf{Export List}, then select LaTeX PDF, LaTeX source, Overleaf project ZIP, or Word document. Papers excluded from export remain visible in My Publications but are omitted from the generated list. The export follows the current visible table, including its search, role/journal filters, and row order, and excludes rows marked Hide. Clear filters and search first if you want the full included list. Exporting never replaces the SQLite source.

The Overleaf option creates a project ZIP, copies it to the Desktop, and opens Overleaf using \hyperref[website-link-settings]{Website Links}. The status message reports the saved path. It does not upload your files or create an Overleaf project automatically; import the ZIP yourself on the website.

\begin{GewuRelated}
\textbf{See also:} \hyperref[citation-settings]{Open Source Settings} · \hyperref[paper-details]{Use Paper Details} · \hyperref[privacy]{Local Data and Privacy}
\end{GewuRelated}

\chapter{Journal Monitor}\label{chapter-15}

\section{Journal Monitor}\label{journal-monitor}

\begin{GewuLead}
Journal Monitor collects papers for a selected publication date and lets you narrow the result by journal, preprint source, and field.
\end{GewuLead}

\par\noindent\begin{minipage}{\linewidth}

\begin{center}\includegraphics[width=0.96\linewidth,height=0.72\textheight,keepaspectratio]{/Users/lingyuan/PaperMonitor-20260516/assets/user-guide/screenshots/journal-monitor.png}\end{center}

\begin{GewuCaption}
Journal Monitor. Nature is selected for August 12, 2026, with scientific figures visible beside the paper cards. The independent sidebar retains calendar counts, the monitor summary, and journal and field filters.
\end{GewuCaption}

\end{minipage}\par\medskip

\subsection{Refresh a date}\label{refresh-a-date}

\begin{enumerate}
\tightlist
\item
  \textbf{Select a day in the calendar.} Small numbers below a day show stored result counts.
\item
  \textbf{Choose \textbf{Refresh}.} The app requests all configured journal sources for that date.
\item
  \textbf{Control-click the journal in Journal Filter and use \textbf{Refresh This Journal} when only that journal needs another attempt.} This avoids re-requesting journals that are already complete.
\item
  \textbf{Choose \textbf{arXiv} separately if needed.} This keeps journal and preprint work explicit.
\item
  \textbf{Watch the bottom status area.} It reports the current source and completion without changing the page height.
\end{enumerate}

\subsection{Manage monitored journals}\label{manage-monitored-journals}

Open \textbf{Journal Settings} to add or manage a journal by name, ISSN, or known publication history. When an ISSN is missing, the app can look up a matching journal record before saving it. Review the match rather than relying on a similar title alone.

Drag journals vertically in \textbf{Journal Filter} or \textbf{Journal Settings} to set their order. The saved order controls the filter list and the journal groups in the report. Removing a journal disables the subscription without erasing its previously collected papers or its saved history.

Your enabled list, disabled history, ISSNs, source identifier, section policy, logo source, and display order are saved in SQLite. Product updates can add better built-in source definitions, but those defaults do not replace this personal state. Adding a previously removed journal restores it at the end of the active order.

\subsection{Understand counts and filters}\label{understand-counts-and-filters}

The number beside each journal is the total stored count for that journal on the selected day. It is deliberately independent of \textbf{Field Filter}. The field checkboxes narrow the paper cards you can see; they do not rewrite the journal totals. A journal can therefore show a nonzero count while the current field selection hides every one of its papers.

Calendar counts summarize stored day results. A zero count is meaningful only after the relevant source completed successfully. Read the source status before interpreting zero as ``the journal published nothing.''

\subsection{Daily scheduled updates}\label{daily-scheduled-updates}

Journal Monitor is designed to keep a daily cached record through its scheduled update workflow. A scheduled run and a manual refresh use the same stored day results, so opening a date does not require downloading it again. Use a manual refresh when a scheduled source was missed, partial, or needs verification.

The scheduler belongs to the running local UI service; it does not wake a sleeping Mac or run while that service is stopped. See \hyperref[startup]{When automatic updates run} for the current timing and restart boundaries.

\subsection{Journal logos and mathematical text}\label{journal-logos-and-mathematical-text}

Configured journal logos are downloaded into a local cache and reused from disk. The first successful request may require network access; changing dates and later launches should not repeatedly download the same logo. The cache is presentation data and can be rebuilt without changing subscriptions or papers.

Supported inline LaTeX in titles and abstracts is typeset after the math renderer becomes ready. On a cold launch, raw source may be visible briefly, then the current day is rendered again automatically; changing to another date and back should not be necessary. Title normalization also preserves recognized proper names and source capitalization such as \emph{Su--Schrieffer--Heeger} and \emph{Monte Carlo} while still cleaning all-uppercase metadata and formula spacing.

\subsection{How discovery provenance is preserved}\label{how-discovery-provenance-is-preserved}

Each stored Journal result is related to the canonical paper, the fetch run, and the source that reported it. The run preserves the journal, requested date window, status, and timing. Seeing the same paper from Crossref and OpenAlex records both source observations against the same paper identity; seeing it again within one run does not create a duplicate relationship. Later runs remain separately traceable.

The same canonical paper can also acquire Subject Monitor and Author Monitor relationships. These discovery paths explain why the work is visible; they do not automatically add a Library placement or create a local Paper Files resource.

\subsection{Filter and inspect}\label{filter-and-inspect}

Use \textbf{Journal Filter}, \textbf{Preprint Filter}, and \textbf{Field Filter}. Journal and field choices combine to determine the visible cards, while the journal counts continue to describe the full selected day. Enable \textbf{Figures} to show a representative image when available. Selecting a result opens its report and actions.

\subsection{Act on a result}\label{act-on-a-result}

\begin{itemize}
\tightlist
\item
  \textbf{Open Website} opens the publisher or best known record using \hyperref[website-link-settings]{Website Links}: a new built-in Browser tab by default, or the macOS default browser when selected.
\item
  \textbf{Add to Library} begins the placement workflow described below; after that, \textbf{Show in Library} returns to the managed paper.
\item
  \textbf{Download} retrieves the paper file. \textbf{PDF} opens an existing local PDF in Preview, and \textbf{Delete PDF} removes that managed download and resets its status.
\item
  When a reading note exists, \textbf{Note} opens it in the right-side drawer. The list compresses but keeps its own scroll position. \textbf{Delete Note} removes the note and resets analysis status.
\end{itemize}

\subsection{Read an existing Note}\label{read-an-existing-note}

\textbf{Note} appears when a reading note has already been generated. It opens the saved note on demand in a right-side drawer; the report area compresses and each side scrolls independently. Drag the divider to resize the note. \textbf{Hide list} gives the note more space; \textbf{Show list} restores the report. Closing the drawer keeps the note. \textbf{Delete Note} is a separate action that removes it and resets its analysis status. Links in the note follow \hyperref[website-link-settings]{Website Links}.

\subsection{Add a Journal result to Library}\label{add-a-journal-result-to-library}

\begin{enumerate}
\tightlist
\item
  \textbf{Choose \textbf{Add to Library} on the result.} The Library destination chooser opens without leaving the selected date.
\item
  \textbf{Choose the target Library folder.} This is the visible folder placement for the managed paper.
\item
  \textbf{Complete the add.} The card changes to \textbf{Show in Library}; use it to reveal the existing placement later.
\end{enumerate}

The operation links the Journal discovery to the canonical paper, preserves its source and run provenance, and adds the chosen Library placement. It does not automatically download a PDF. Choose \textbf{Download} separately when you also want a local file. If the paper is already managed, \textbf{Show in Library} reveals it instead of creating a duplicate identity.

Per-paper status text and the bottom progress area make completed downloads and background work visible, which helps avoid repeating the same action.

\begin{quote}
\textbf{Note}

``No matching papers'' and ``source fetch failed'' are different outcomes. Check the source status before concluding that a date is empty.
\end{quote}

\begin{GewuRelated}
\textbf{See also:} \hyperref[subject-monitor]{Subject Monitor} · \hyperref[network-access]{Publisher and Network Access} · \hyperref[import-files]{Add Papers and Files}
\end{GewuRelated}

\chapter{Subject Monitor}\label{chapter-16}

\section{Subject Monitor}\label{subject-monitor}

\begin{GewuLead}
Subject Monitor discovers papers for reusable topics. Each subject stores a name and literal discovery terms; an optional copied prompt can help you draft those terms outside the app.
\end{GewuLead}

\par\noindent\begin{minipage}{\linewidth}

\begin{center}\includegraphics[width=0.96\linewidth,height=0.72\textheight,keepaspectratio]{/Users/lingyuan/PaperMonitor-20260516/assets/user-guide/screenshots/subject-monitor.png}\end{center}

\begin{GewuCaption}
Subject Monitor. Superconductivity is selected, with matching papers and numeric calendar counts for that subject. Recorded counts include zero; dates without stored results remain blank.
\end{GewuCaption}

\end{minipage}\par\medskip

\subsection{Create a subject}\label{create-a-subject}

\begin{enumerate}
\tightlist
\item
  \textbf{Choose \textbf{Add Subject}.}
\item
  \textbf{Enter a clear name and discovery keywords.} Describe the field broadly enough to find variants, but avoid unrelated general terms.
\item
  \textbf{Use \textbf{Copy Suggested Prompt} if you want Codex-assisted wording.} Review the prompt before using it.
\item
  \textbf{Save and arrange the subject.} Subjects can be edited, deleted, and reordered.
\end{enumerate}

Keywords use case-insensitive literal matching. Enter the words or phrases you expect to appear; they are not regular expressions and are not silently expanded into unrelated concepts.

\par\noindent\begin{minipage}{\linewidth}

\begin{center}\includegraphics[width=0.96\linewidth,height=0.72\textheight,keepaspectratio]{/Users/lingyuan/PaperMonitor-20260516/assets/user-guide/screenshots/subject-editor.png}\end{center}

\begin{GewuCaption}
Add Subject. Enter a name and one keyword or phrase per line. The same form is used when editing an existing subject; Cancel leaves its saved keywords unchanged.
\end{GewuCaption}

\end{minipage}\par\medskip

\subsection{Edit keywords without losing history}\label{edit-keywords-without-losing-history}

Control-click a subject and choose \textbf{Edit}. Saving a new name or keywords keeps its existing paper list, calendar counts, selection, and reading position. New keywords apply to future retrieval. Existing date windows, including empty and incomplete results, remain preserved after a keyword change; background retries and later local arXiv matching do not silently reclassify those windows.

Choose \textbf{Refresh This Subject} explicitly to rebuild that subject\textquotesingle s selected date with the current keywords. A successful refresh replaces only that result window; other dates and subjects remain unchanged. An incomplete refresh preserves previous results. Editing during an active retrieval discards the obsolete result instead of replacing history. Deleting a subject still clears its Subject Monitor matches, without deleting the underlying shared papers or Library files. Previously cleared matches from an older software version are not automatically reconstructed by this fix.

\subsection{Keep one subject selected across dates}\label{keep-one-subject-selected-across-dates}

Selecting a subject makes it the active topic until you choose another one. Moving to a date with no result does not switch the tab back to the first subject, and a subject with zero papers remains selectable. The calendar\textquotesingle s day numbers describe the active subject, so you can see which dates contain papers for that topic rather than a total across every subject. The count beside each subject tab describes that subject on the selected date.

The calendar shows the stored count as a number, including \textbf{0}, for the selected subject and date. Counts have no completeness marker. A date with no stored results and no fetch record remains blank. While a fetch is running, the progress indicator remains visible.

The \textbf{Discovery} row describes the recorded retrieval status for the selected subject. A completed retrieval is not a guarantee that every relevant paper has been found. Switching subjects changes the counts and recorded status; reopening the app preserves the status stored in SQLite.

\subsection{Refresh subjects}\label{refresh-subjects}

\begin{enumerate}
\tightlist
\item
  \textbf{Select a calendar date.} Selecting a date only loads the result already stored for that day; it does not start a network request.
\item
  \textbf{Choose \textbf{Refresh}.} The command refreshes every configured subject for the currently selected date. The arXiv display option is enabled by default; clear it when you want to hide arXiv cards without changing what is stored.
\item
  \textbf{Read source status.} A source can be complete, partial, empty, or failed. Source-level retries are managed by the background scheduler.
\item
  \textbf{Cancel if necessary.} Completed results remain distinguishable from unfinished work.
\end{enumerate}

\subsection{Refresh one subject}\label{refresh-one-subject}

Control-click a subject and choose \textbf{Refresh This Subject} when only that topic needs another pass. Use the main \textbf{Refresh} command when every configured subject should run for the selected date.

The current candidate layer can draw from local arXiv data, local APS RSS data, OpenAlex, Crossref, Semantic Scholar, and PubMed. It does not depend on a fixed number of feeds. Results still pass topic filtering and deduplication before display.

\subsection{Inspect and keep a result}\label{inspect-and-keep-a-result}

Select a paper to read its title, authors, abstract, source, date, identifier, topic score, and matched subject terms. Use \textbf{Open Website}, \textbf{Download}, \textbf{PDF}, and the note controls in the same way as Journal Monitor. Choose \textbf{Add to Library} to create a managed reference, or \textbf{Show in Library} when one already exists.

\subsection{Add a Subject result to Library}\label{add-a-subject-result-to-library}

\begin{enumerate}
\tightlist
\item
  \textbf{Choose \textbf{Add to Library} on the matching paper.} The subject and selected date remain in place while the destination chooser opens.
\item
  \textbf{Select the target Library folder.} The same canonical paper can be placed in a durable topic hierarchy without losing the Subject relationship.
\item
  \textbf{Complete the add.} Use \textbf{Show in Library} on later visits to reveal the managed placement.
\end{enumerate}

Adding the result preserves the matched subject, score, terms, and discovery provenance. The paper remains visible in Subject Monitor immediately after adding it; its action becomes \textbf{Show in Library}, while the selected subject, date, and reading position stay in place. It does not automatically download a PDF; use \textbf{Download} separately. Repeated discovery of the same identity updates relationships rather than creating unrelated paper records.

\begin{quote}
\textbf{Important}

Subject Monitor rejects Zenodo records during ingestion and when cached results are displayed. There is no Zenodo show/hide option.
\end{quote}

\begin{GewuRelated}
\textbf{See also:} \hyperref[journal-monitor]{Journal Monitor} · \hyperref[network-access]{Publisher and Network Access}
\end{GewuRelated}

\chapter{Author Monitor}\label{chapter-17}

\section{Author Monitor}\label{author-monitor}

\begin{GewuLead}
Author Monitor follows researchers and separates a first-import baseline from newly discovered work.
\end{GewuLead}

\par\noindent\begin{minipage}{\linewidth}

\begin{center}\includegraphics[width=0.96\linewidth,height=0.72\textheight,keepaspectratio]{/Users/lingyuan/PaperMonitor-20260516/assets/user-guide/screenshots/author-monitor.png}\end{center}

\begin{GewuCaption}
Author Monitor. Hong Ding\textquotesingle s profile is selected. The paper title leads the details pane, followed by compact actions and a visible Gewu Key; local portraits keep the author list stable.
\end{GewuCaption}

\end{minipage}\par\medskip

\subsection{Add an author}\label{add-an-author}

\begin{enumerate}
\tightlist
\item
  \textbf{Choose the add-author control.}
\item
  \textbf{Enter a name and optional affiliations, or an exact profile URL.} Google Scholar and ResearchGate profile URLs are supported where the profile can be resolved.
\item
  \textbf{Confirm the correct person.} The initial result set becomes the baseline rather than a batch of unread ``new'' papers.
\end{enumerate}

\subsection{Understand profile and publication sources}\label{understand-profile-and-publication-sources}

Google Scholar profiles refresh through Scholar with shared pacing and cooldown. A ResearchGate URL identifies the public profile; its publications are matched separately through OpenAlex when configured. If unavailable, Gewu Insight automatically falls back to Semantic Scholar. Confirm the proposed person, especially for similar names.

Without a verified publication-data match, the ResearchGate profile and homepage can be saved, but automatic publication refresh remains unavailable. Profile, publication, and citation sources remain distinct evidence; a fallback does not change the author identity.

\par\noindent\begin{minipage}{\linewidth}

\begin{center}\includegraphics[width=0.96\linewidth,height=0.72\textheight,keepaspectratio]{/Users/lingyuan/PaperMonitor-20260516/assets/user-guide/screenshots/author-add.png}\end{center}

\begin{GewuCaption}
Add Monitored Author. Enter an author name or a profile URL. This real capture also shows the local Google Scholar pause and retry time; ResearchGate metadata sources remain available.
\end{GewuCaption}

\end{minipage}\par\medskip

\subsection{Review updates}\label{review-updates}

Later refreshes can mark newly discovered publications as unread. Use the \textbf{New} filter and other author/result controls to focus the list. Select \textbf{Date} or \textbf{Citations} to sort; selecting the active sort reverses its direction. \textbf{Conferences · n} opens conference publications in a separate, independently scrolling popover, so they are not missing from the ordinary article list. Close the popover with its close control or Escape.

Control-click \textbf{New} and choose \textbf{Read All}. For one selected author, it marks every unread publication as read, including offscreen papers; other authors keep their unread state, even for a shared paper. With \textbf{All Authors} selected, it marks all authors\textquotesingle{} new publications as read. Papers and Library placements are preserved. The New filter remains active; click New again to see read papers. Read state persists after reopening.

\subsection{Refresh one author or find missing identifiers}\label{refresh-one-author-or-find-missing-identifiers}

Control-click the author and choose \textbf{Refresh This Author} to update that person\textquotesingle s publications. \textbf{Find Missing Identifiers} checks eligible papers for that author. To check only one unidentified paper, Control-click the paper and choose \textbf{Find Missing Identifier}; it is disabled when an identifier already exists or a lookup is running. These online lookups differ from \textbf{Get Identifier from Files}, which uses an available local PDF.

\subsection{Update citation profiles}\label{update-citation-profiles}

For a Google Scholar author, Control-click the author, choose \textbf{Update All Scholar Citations}, review the request estimate, then confirm \textbf{Update All Citations}. This updates that author's full profile paper by paper and can take a long time. Follow the bottom progress status until it finishes.

\subsection{Manage the profile image}\label{manage-the-profile-image}

Control-click an author and choose \textbf{Edit Homepage} to edit the website, or \textbf{Change Photo} / \textbf{Add Photo} to choose a local portrait. Downloaded Scholar avatars are saved locally. Opening the author list does not need to reopen SQLite or contact the image host for every row; a manual photo takes precedence. Without a valid local photo, the list can show initials until you refresh or replace the image.

\subsection{Discovery entries and managed papers}\label{discovery-entries-and-managed-papers}

An Author Monitor result is a cataloged shared record, with a Gewu Key shown in its details even before it has Library placements or local files. A missing key displays \texttt{(undefined)} until record data is available. Add to Library preserves the existing key and updates the view; downloading or attaching a file creates a resource relationship.

\subsection{Add an author result to Library}\label{add-an-author-result-to-library}

\begin{enumerate}
\tightlist
\item
  \textbf{Select the publication and choose \textbf{Add to Library}.} This opens the Library destination chooser without removing the author from view.
\item
  \textbf{Choose the target Library folder.} The selected placement becomes the durable organizational path for the shared paper record.
\item
  \textbf{Complete the add.} The action becomes \textbf{Show in Library}, which reveals the managed paper later.
\end{enumerate}

The author-publication relationship and canonical identity remain intact. Adding a placement does not automatically download a PDF, and a paper already managed elsewhere is reused rather than duplicated.

\subsection{Remove a monitored author}\label{remove-a-monitored-author}

Control-click the author and choose \textbf{Remove from Monitor}. Removing an author deletes the monitoring relationship and its monitor state. It does not delete papers already added to Library, their managed PDFs, or their Paper Files bundles.

\begin{quote}
\textbf{Note}

Google Scholar requests are paced and protected by a cooldown. Repeated manual refreshes during a block do not improve the result.
\end{quote}

\begin{GewuRelated}
\textbf{See also:} \hyperref[my-publications]{My Publications} · \hyperref[network-access]{Publisher and Network Access} · \hyperref[privacy]{Local Data and Privacy}
\end{GewuRelated}

\chapter{Browser}\label{chapter-18}

\section{Browse, Bookmark, and Revisit Websites}\label{browser}

\begin{GewuLead}
Browser is the sixth main tab, after Author Monitor. By default, website links across Library, My Publications and all three Monitors switch here, opening a new webpage tab without replacing existing pages. Choose Setting \textgreater{} Website Links to use the system default browser instead.
\end{GewuLead}

\par\noindent\begin{minipage}{\linewidth}

\begin{center}\includegraphics[width=0.96\linewidth,height=0.72\textheight,keepaspectratio]{/Users/lingyuan/PaperMonitor-20260516/assets/user-guide/screenshots/browser.png}\end{center}

\begin{GewuCaption}
Browser in use. Captured from the original App displaying a fully loaded Nature Physics PDF. The real bookmark collection, webpage tab and Save to Gewu button are visible. This is a live interface capture, not a generated demonstration.
\end{GewuCaption}

\end{minipage}\par\medskip

\subsection{Search and navigate}\label{search-and-navigate}

\begin{enumerate}
\tightlist
\item
  \textbf{Open Browser.} On first use, Home displays your local Favorites folder. Later switches return to your selected webpage; use Home to return to Favorites. A paper link opens its destination directly in the unused first tab. No Google webpage is loaded for Home.
\item
  \textbf{Use the address field.} Enter a website address or search terms. Plain search terms go to Google; an address opens directly. Use {{\symbolfont ⌘}L} to focus the field.
\item
  \textbf{Navigate normally.} Back, Forward, Reload/Stop, and Home act on the selected tab. Home uses the content area for Favorites; websites, search results, sign-in pages and PDFs use it for browsing.
\item
  \textbf{Switch between workspaces.} Returning from Library or a Monitor preserves the current page and browsing position without reloading merely because the main tab changed.
\end{enumerate}

\textbf{Return to the source workspace: {{\symbolfont ⌘}{[}}.} After a website link in Library, My Publications, Journal Monitor, Subject Monitor, or Author Monitor brings you here, press Command-left bracket to return to that workspace\textquotesingle s retained selection and scroll position. The physical left-bracket key also works with a Chinese input method selected. It works from the webpage or address field and leaves Browser tabs open; it is not webpage Back. The source lasts for this Browser detour only: manually leaving Browser clears it, and entering Browser directly does not create one. Close any open dialog first.

\textbf{Local Home, online search.} Favorites is available without loading Google, even after restarting the App. Home has no Google logo, search box or account controls; use the toolbar address field to search. Reload on Home refreshes local bookmarks, not a website. Site icons may need a connection; missing icons show initials. Webpages and sign-in remain online and share Browser\textquotesingle s existing profile.

When you search or follow a link from Home, Home remains visible while the destination begins loading, then gives way to that destination. It should not flash the page you visited before Home. A cancelled or failed pending navigation leaves Home available rather than presenting that older page as the result.

{{\symbolfont ⌘}F} opens Find on Page. Enter moves to the next match; Shift-Enter moves to the previous match. Browser is available in the native \textbf{Gewu Insight.app}; opening the local UI in a separate web browser does not embed another browser engine.

\textbf{Favorites on Home.} Favorites/Favourites (including localized names) supplies the tiles; use \textbf{Create Favorites Folder} if missing. Open nested folders and return through breadcrumbs; use \textbf{Add Favorite} or right-click to edit. The grid uses the whole main area, without a Google section or divider. Excess tiles use Previous/Next, not vertical scrolling. Bottom bars align; hiding the sidebar keeps Favorites visible. Opening google.com explicitly shows an ordinary webpage, without Home-specific styling or zoom. Bookmarks are not sent to Google.

\par\noindent\begin{minipage}{\linewidth}

\begin{center}\includegraphics[width=0.96\linewidth,height=0.72\textheight,keepaspectratio]{/Users/lingyuan/PaperMonitor-20260516/assets/user-guide/screenshots/browser-home.png}\end{center}

\begin{GewuCaption}
Browser Home. The original App displays the user\textquotesingle s real Favorites and bookmarks. Home is local; use the address field or a bookmark to open a website.
\end{GewuCaption}

\end{minipage}\par\medskip

\subsection{Use multiple webpage tabs}\label{use-multiple-webpage-tabs}

Below the address toolbar, adjoining rectangular tabs show icons, titles and \textbf{×} above the webpage column only. The sidebar begins directly below the toolbar; it is not pushed down by the webpage tabs. The top-right \textbf{+} (\textbf{New Tab}, {{\symbolfont ⌘}T}) opens Home. Switching preserves each page, position and Back/Forward history without reloading or stopping background loads.

Right-click links for \textbf{Open Link in New Tab}, or Command-click. New-window links also create tabs, including inline PDFs. Bookmarks/Favorites/History offer \textbf{Open in New Tab}. New tabs are selected immediately; ordinary clicks use the current tab.

\textbf{×} or {{\symbolfont ⌘}W} closes a tab; closing the last opens Home. Left/Right switches focused tabs. Up to 32 tabs stay in memory, not restored after App quit. Save drafts before closing; bookmarks, history, downloads and sign-ins remain. AutoFill only triggers in the visible page.

\subsection{Save and organize bookmarks}\label{save-and-organize-bookmarks}

The sidebar\textquotesingle s top tabs switch between \textbf{Bookmarks} and \textbf{History}, without a repeated heading. The compact \textbf{B} (Bookmarks menu), \textbf{+} (New Folder), and Search sit along the bottom, even at the minimum sidebar width. History shows Clear History and Search there instead. Drag the sidebar\textquotesingle s right edge to resize it or use the toolbar to collapse it. Width, tab, and expanded folders are saved locally.

\begin{enumerate}
\tightlist
\item
  \textbf{Add a page.} Use the bookmark toolbar button, \textbf{Share \textgreater{} Add Bookmark}, or {{\symbolfont ⌘}D}. The native sheet starts with the page title and address.
\item
  \textbf{Choose a location.} Edit the name or address and choose a folder under \textbf{Add to}. The previous destination is remembered. Saving the same address in the same folder reuses the existing bookmark.
\item
  \textbf{Create nested folders.} Use \textbf{New Folder}, or right-click a folder to create a child folder. The editor can also rename a folder or choose another parent.
\item
  \textbf{Arrange the tree.} Drag an item into the middle of a folder row to move it inside, or to a row edge to reorder it. Drop on the top \textbf{Bookmarks} tab to move to the root. A folder cannot move inside itself or its descendants.
\item
  \textbf{Save by dragging.} Drag the website icon beside the address field, or a webpage link, into a bookmark folder. Confirm the name and destination in the same native sheet.
\end{enumerate}

Search filters bookmark names and addresses. Right-click a bookmark for \textbf{Open}, \textbf{Edit Bookmark}, or \textbf{Delete}. Deleting a folder requires confirmation and removes its contained bookmarks, not websites, browsing history, downloaded files, or papers. A bookmark is a saved link, not an offline webpage copy.

\subsection{Import and export bookmarks}\label{import-and-export-bookmarks}

Choose \textbf{B \textgreater{} Import} at the bottom of Bookmarks and select a Safari or Chrome bookmark HTML file. Safari versions with \textbf{File \textgreater{} Export Bookmarks} produce HTML; newer Safari versions may use \textbf{File \textgreater{} Export Browsing Data to File} and produce ZIP. Select bookmarks only when exporting. Gewu accepts a ZIP containing exactly one bookmark HTML file and ignores all other members. In Chrome, use \textbf{Bookmark Manager \textgreater{} menu \textgreater{} Export Bookmarks}.

Imports go into \textbf{Imported Bookmarks}, retaining folders, order, names, and supported links. Repeating an import merges matching folder names and reuses the same URL in the same folder; existing bookmarks are not overwritten. The result reports added bookmarks, new folders, duplicates, and skipped links. Local addresses, application links, and recognized authorization links are skipped. Limits are 10 MB, 20,000 source entries, and 64 folder levels; invalid files do not partially import.

\textbf{B \textgreater{} Export} saves the full tree as standard bookmark HTML; right-click a folder for \textbf{Export Folder}. Safari and Chrome can import this format. Exports contain no history, cookies, or passwords, but their bookmark names and URLs are readable by anyone with the file. Keep personal exports private.

\subsection{Sync one source folder into Gewu}\label{sync-one-source-folder-into-gewu}

\begin{enumerate}
\tightlist
\item
  \textbf{Choose \textbf{B \textgreater{} Sync from Safari} or \textbf{Sync from Chrome}.} This is a manual, one-way operation, not a background watcher or a write-back to either browser.
\item
  \textbf{Select the source bookmark file.} For Safari, choose \texttt{\textasciitilde{}/Library/Safari/Bookmarks.plist}. For Chrome, choose \texttt{Bookmarks} in the intended profile under \texttt{\textasciitilde{}/Library/Application\ Support/Google/Chrome/}, such as Default or Profile 1. The picker remembers the last chosen location; select the same file for later refreshes. Only this explicitly selected bookmark file is read.
\item
  \textbf{Choose one folder.} The next sheet lists source folder paths and bookmark counts. \textbf{Sync Selected Folder} includes that folder and its descendants, not its siblings or the whole profile. Cancel leaves Gewu bookmarks unchanged.
\item
  \textbf{Review the result.} New links and folders are added. Stable source identifiers allow later source title/address changes to update entries created by sync. Pre-existing matches and items edited or moved in Gewu are kept. Source deletions do not delete local bookmarks. The result reports additions, updates, local changes kept, and skipped links.
\end{enumerate}

The selected folder is placed at the Bookmarks root; matching folder names and same-folder URLs are reused. Sync tracking lives in SQLite, independently of papers and website sessions. Legacy source entries without stable identifiers merge by folder/name/URL rather than overwriting local edits. Switching source files or profiles is a separate source. Limits match import (10 MB, 20,000 entries, 64 levels), with at most 2,000 folders in the chooser. Invalid input or a failed transaction leaves no partial sync.

If macOS denies Safari file access, use a Safari bookmark HTML export with \textbf{B \textgreater{} Import}. No privacy settings are changed, and no password, cookie, history, or cloud-account databases are read. Only bookmarks already saved in the chosen local source file are available.

\subsection{Navigate Browser with arrows}\label{navigate-browser-with-arrows}

From the top \textbf{Browser} workspace tab, press Down to enter the toolbar. Left/Right moves among available toolbar controls; Enter activates one. At the address field, Enter or typing starts editing; its arrows then edit text or browse address suggestions. Escape returns to the toolbar. Up from the toolbar returns to the top Browser tab, where Left/Right changes workspaces. The Library exception is unchanged: Up within Library content does not ascend to its top workspace tab.

Down from the toolbar selects the current webpage tab. Left/Right selects adjacent webpage tabs; Down enters Favorites or the native webpage. Favorites uses directional tile selection and Enter to open. At the left edge, Left enters the visible sidebar on its currently active Bookmarks/History tab; Down traverses its list and footer controls. Right returns to the main area. A hidden sidebar is never selected. In the sidebar header, Left focuses the other collection tab and Enter activates it.

Inside an ordinary webpage, arrows remain WebKit scrolling controls. Once Up reaches the document top, it can hand focus back to webpage tabs; two more Up presses reach the top Browser tab. Left at a document\textquotesingle s left boundary can enter the visible sidebar. Inputs, editable text, framed/PDF content, and recognized or uncertain custom widgets keep their own keys; no form values are read or filled. A boundary check runs only after an eligible arrow press, without background polling. Use {{\symbolfont ⌘}L}, then Escape, to return to the toolbar from a document that retains its own arrows.

\subsection{Website icons}\label{website-icons}

Bookmarks, Favorites and webpage tabs use a disposable local icon cache, separate from SQLite. Icons come directly from sites (at most two concurrent requests), without cookies, credentials, whole-article fetching or third-party icon services. Visits may supply a nonstandard icon URL. Missing icons use globes or Home initials; bookmark lists are never sent to an icon service.

\subsection{Revisit and clear history}\label{revisit-and-clear-history}

History groups loaded main pages by Today, Yesterday, and date, with names, sites, visit times, search, and pagination. Subresources, main-tab switches, and cached Home returns add no visits; Reload updates the existing record. Click to revisit; right-click to bookmark or delete a visit.

\textbf{Clear History} confirms Last hour, Last 24 hours, Last 7 days, or All history. It preserves bookmarks, downloads, cookies, and sign-ins. History remains until deleted; authorization callbacks and recognized signed-download URLs are excluded.

\subsection{Address suggestions from your history}\label{address-suggestions-from-your-history}

As you type in the toolbar address field, Gewu searches saved local visit titles and URLs. Address prefixes rank first, followed by matching titles; visit frequency and recency help order the results. A matching address can complete the typed prefix with its remaining text selected. These are local history suggestions, not live Google suggestions, and typing does not send a suggestion request to a search engine.

Use {{\symbolfont ↑}}/{{\symbolfont ↓}}, then Enter to open a result, or click it to open directly. Escape dismisses suggestions and restores your typed query. Without a chosen completion, Enter opens an address or searches Google normally. Clearing history also removes those visits from suggestions; bookmark search remains a separate sidebar control.

\subsection{Read PDFs and save downloads}\label{read-pdfs-and-save-downloads}

Viewable PDFs stay inline, including attachment/download links, with article/PDF Back/Forward navigation. Save via Download or \textbf{Share \textgreater{} Save Download}; unsupported attachments use a save sheet. This uses Browser\textquotesingle s session, not Python\textquotesingle s downloader. For HTML viewers, use the site\textquotesingle s PDF link. Save Download writes an ordinary downloaded file. To connect an article and its PDF to your research library, use \textbf{Save to Gewu} below or \hyperref[import-files]{import the saved file}. Library PDFs still open in macOS Preview.

\subsection{System networking and Clash rules}\label{system-networking-and-clash-rules}

WebKit follows macOS networking without changing proxies/Clash. When system proxy or TUN routes traffic to Clash, its rules choose DIRECT/proxy. \textbf{System network} is policy, not proof of routing. DIRECT needs a school-recognized exit or institutional login. PDF/login/CDN domains match separately; browser-specific extensions/process rules are not inherited. Verify Clash connections; see \hyperref[network-access]{Publisher and Network Access}.

\subsection{Save an article and PDF to Library}\label{browser-save-gewu}

\textbf{Save to Gewu} opens a compact dialog with \textbf{File type}, \textbf{Add to Library}, Cancel and Save to Library. Click the double-arrow \textbf{Details} to expand article identity and PDF controls. Set its default under \textbf{Setting \textgreater{} Website Links}.

\begin{enumerate}
\tightlist
\item
  \textbf{Open the article or PDF.} Follow \textbf{Open Website} from a research workspace, or find the article directly in Browser. Wait for the page to finish loading and click \textbf{Save to Gewu}.
\item
  \textbf{Review identity in Details.} Check the title, DOI and Gewu Key. An existing DOI or known arXiv identity reuses its paper record, including a record discovered only in a Monitor. A new article receives its own Key when saved. If needed, enter a DOI/arXiv ID and click \textbf{Identify}. Without an identifier, enter the article title and choose \textbf{Use this title}; metadata can be completed later.
\item
  \textbf{Choose the PDF in Details.} A directly opened PDF is retrieved automatically into temporary staging. On an article page, select a detected PDF link and click \textbf{Get PDF}, or choose a local PDF. Each file must be a readable PDF up to 64 MB. Check the filename and leave \textbf{Include this PDF in Paper Files} selected to attach it.
\item
  \textbf{Choose the destination.} Click the folder button under \textbf{Add to Library}. The same column picker used by Journal and Subject Monitor shows \textbf{My Library} and Spaces on the left, then subfolders in successive columns. Check the full path and click \textbf{Confirm} to select the destination. Choose the file type, then click \textbf{Save to Library} to import. Recognized publisher article PDFs suggest \textbf{Published version}. Recognized arXiv, bioRxiv, medRxiv, ChemRxiv and Research Square preprints suggest \textbf{Preprint}, regardless of the article's journal metadata. Review the suggestion; uncertain files remain \textbf{Unknown}. The first article PDF becomes Default in the main folder, even if Unknown. Supplements stay separate, and manual Default is preserved. Source follows the PDF's actual website: Publisher, Preprint server, Institutional repository or Other. Local files use User import without a webpage source.
\end{enumerate}

Opening Details, identifying or cancelling imports nothing. Save reuses an existing active Library or Space placement, preferring the selected occurrence of that same paper. Attaching a PDF does not add a root-level entry. To add another location, explicitly choose and confirm its folder. Papers with no active placement default to unsorted unless you confirm another folder. Paper Files refreshes immediately. Exact duplicates retain their existing source and settings; distinct versions stay separate, and manual Default is preserved. \textbf{Show in Library} selects the saved occurrence. Browser keeps the webpage.

\textbf{Source context is a hint.} An Open Website link may lead to another work. Gewu checks the current article and PDF identity before attachment. If identification fails, you can explicitly choose the source article when offered. Conflicting identifiers block saving; open the correct PDF separately or select its matching local file.

\textbf{If PDF retrieval fails:} retry after opening the site\textquotesingle s PDF, choose a downloaded local PDF, or save the article alone. Errors stay visible; missing or conflicting identity opens Details. If storage fails after saving the article, retry reuses its record and placement. Browser capture uses its current WebKit session without exporting sign-in cookies. It cannot bypass a site\textquotesingle s login, access rights or embedded-browser restrictions. The collector currently belongs to Gewu\textquotesingle s built-in Browser; a Safari extension is not included.

\par\noindent\begin{minipage}{\linewidth}

\begin{center}\includegraphics[width=0.96\linewidth,height=0.72\textheight,keepaspectratio]{/Users/lingyuan/PaperMonitor-20260516/assets/user-guide/screenshots/save-gewu-compact.png}\end{center}

\begin{GewuCaption}
Save to Gewu: compact review. After the article and PDF have been identified, the compact panel keeps File type, Add to Library and the final save action visible.
\end{GewuCaption}

\end{minipage}\par\medskip

\par\noindent\begin{minipage}{\linewidth}

\begin{center}\includegraphics[width=0.96\linewidth,height=0.72\textheight,keepaspectratio]{/Users/lingyuan/PaperMonitor-20260516/assets/user-guide/screenshots/save-gewu-details.png}\end{center}

\begin{GewuCaption}
Save to Gewu: Details. Expanded Details identifies this existing Nature Physics article by its Gewu Key and DOI. The PDF is Ready and classified as Published version before saving.
\end{GewuCaption}

\end{minipage}\par\medskip

\par\noindent\begin{minipage}{\linewidth}

\begin{center}\includegraphics[width=0.96\linewidth,height=0.72\textheight,keepaspectratio]{/Users/lingyuan/PaperMonitor-20260516/assets/user-guide/screenshots/library-destination-picker.png}\end{center}

\begin{GewuCaption}
Choose a Library destination. Read the columns from workspace to folder to subfolder. The bottom path follows the selection. Confirm selects a destination; Save to Library performs the import.
\end{GewuCaption}

\end{minipage}\par\medskip

\subsection{Privacy and compatibility boundaries}\label{privacy-and-compatibility-boundaries}

Bookmarks/history use \hyperref[privacy]{local SQLite}; cookies/cache use Browser\textquotesingle s persistent profile, not Safari synchronization. Browsing/search reaches websites, but external pages have no Library/credential bridge. Local/file/application addresses are blocked. Standard internal child documents use WebKit\textquotesingle s origin and sandbox rules; they cannot become address-bar destinations. HTTP recovery retries once as same-host HTTPS GET, without replaying forms or bypassing certificates; errors include reason/code. The local server rejects cross-site API requests and unrecognized hosts.

\subsection{Use macOS Passwords}\label{use-macos-passwords}

\textbf{Open Passwords} opens Apple\textquotesingle s app (or use Passwords in System Settings on older Macs). Unlock and choose credentials there yourself. Gewu has no password vault and never imports, exports, lists or relays credentials.

\textbf{Conservative AutoFill trigger.} A single, unmodified click on an explicitly marked HTTPS login username or current-password input can open \textbf{AutoFill \textgreater{} Passwords}. Search, ordinary email, registration, codes, embedded frames and uncertain fields are skipped, as are keyboard focus, double-clicks and dragging. Each field kind is offered once per site per App session, including cancellation; right-click AutoFill retries manually. Navigation or intervening input cancels a pending trigger. macOS handles authentication, selection and filling; Gewu never reads credentials or submits login. Site matching, Safari-style suggestions and successful filling are not guaranteed.

A macOS prompt for \textbf{Gewu Insight WebCrypto Master Key} concerns WebKit\textquotesingle s protection of locally stored webpage cryptographic keys, not a metadata API key. Local development rebuilds can require renewed keychain permission. Handle this system prompt yourself; do not share your Mac password or delete the key to suppress the dialog.

\textbf{Password versus passkey.} A Password entry contains a text password; a Passkey entry is a different sign-in credential and may show only a creation date. A passkey cannot be inserted as text into a password field. Use the website\textquotesingle s passkey sign-in flow where supported, or select a genuine Password entry yourself. See \hyperref[troubleshooting]{AutoFill troubleshooting}; the trigger is not a guarantee that macOS or a particular website can fill the field.

\begin{quote}
\textbf{Website links}

Publisher, Google Scholar, author/profile, preprint, citation-source, Note, remote-PDF and Overleaf links from research views follow \hyperref[website-link-settings]{Website Links}. They open here by default; System Default Browser sends them to the browser chosen in macOS. Navigation already inside Browser stays in Browser. Local PDF/Preview and Finder actions stay native. For incompatible sites, explicitly use \textbf{Open in Safari}; cookies are not transferred. The browser-hosted UI opens ordinary tabs. Never bypass access controls.
\end{quote}

\chapter{Settings and Safety}\label{chapter-19}

\section{Open Source Settings}\label{citation-settings}

\begin{GewuLead}
Built-in sources work without setup. Optional credentials improve access or coverage for supported services.
\end{GewuLead}

\par\noindent\begin{minipage}{\linewidth}

\begin{center}\includegraphics[width=0.96\linewidth,height=0.72\textheight,keepaspectratio]{/Users/lingyuan/PaperMonitor-20260516/assets/user-guide/screenshots/open-source.png}\end{center}

\begin{GewuCaption}
Open Source Settings. Configure discovery and metadata services in one panel. Saved credentials remain masked; each service has its own Save, Test and Clear controls.
\end{GewuCaption}

\end{minipage}\par\medskip

\subsection{Open settings}\label{open-settings}

Choose \textbf{Setting \textgreater{} Open Source Settings}, or press \texttt{{\symbolfont ⌘}}\texttt{,}.

\subsection{Available services}\label{available-services}

{\def\LTcaptype{none} % do not increment counter
\GewuTableStart
\begin{longtable}[]{@{}>{\RaggedRight\arraybackslash}p{0.27\textwidth}>{\RaggedRight\arraybackslash}p{0.63\textwidth}@{}}
\toprule\noalign{}
\rowcolor{GewuTableHeader}
Service & Configuration \\
\midrule\noalign{}
\endhead
\bottomrule\noalign{}
\endlastfoot
Crossref, arXiv, DataCite & Built in; no key field is required. \\
OpenAlex & Optional OpenAlex API key. \\
Semantic Scholar & Optional API key. Anonymous access remains available. \\
NASA ADS & Optional API token for astronomy and physics coverage. \\
\end{longtable}
\GewuTableEnd
}

\subsection{Save a credential}\label{save-a-credential}

\begin{enumerate}
\tightlist
\item
  \textbf{Paste the credential into its service field.}
\item
  \textbf{Choose \textbf{Save} beside that service.} Each service has its own Save, Test, and Clear controls.
\item
  \textbf{Choose \textbf{Test} and read the status underneath.} This checks service access; it does not refresh every paper. Then refresh one paper to inspect the available citation evidence.
\end{enumerate}

A configured field shows fixed round dots by default, while the detailed status below continues to report saving, connection, or quota information. The dots do not expose the key or its length. Click \textbf{Show} inside the field to reveal the saved key, or inspect a newly typed key; click \textbf{Hide} to mask it again. A saved key is retrieved only on Show, and its displayed value is cleared on Hide or when settings closes. Type a replacement and choose \textbf{Save} to change the key; \textbf{Clear} removes the locally saved key. A configured environment key can remain available. Avoid screenshots or screen sharing while a key is visible.

Semantic Scholar sends its key in the official \texttt{x-api-key} header. Requests share a low-rate queue and use exponential backoff after rate limits or temporary failures.

OpenAlex status reports the remaining daily free allowance and its reset time when available. Gewu uses only that free allowance and falls back to other metadata sources when it is exhausted; it does not spend paid OpenAlex credits. A saved key is not proof that every remote request will succeed.

\begin{quote}
\textbf{Important}

Credentials are local-only: store them in local SQLite or approved environment variables. Do not place keys in source files, logs, exported reports, tests, or Git.
\end{quote}

\begin{GewuRelated}
\textbf{See also:} \hyperref[arxiv-scope-settings]{arXiv Scope Settings} · \hyperref[field-settings]{Field Settings} · \hyperref[privacy]{Local Data and Privacy} · \hyperref[network-access]{Publisher and Network Access}
\end{GewuRelated}

\section{arXiv Scope Settings}\label{arxiv-scope-settings}

\begin{GewuLead}
Choose one or more arXiv subject archives for discovery. Selections can span several research fields or combine several groups within one field.
\end{GewuLead}

\par\noindent\begin{minipage}{\linewidth}

\begin{center}\includegraphics[width=0.96\linewidth,height=0.72\textheight,keepaspectratio]{/Users/lingyuan/PaperMonitor-20260516/assets/user-guide/screenshots/arxiv-scope.png}\end{center}

\begin{GewuCaption}
arXiv Scope Settings. Choose subject archives for future Journal Monitor discovery. Expand a group to inspect its categories; the displayed Condensed Matter selection is an example of a saved preference.
\end{GewuCaption}

\end{minipage}\par\medskip

\subsection{Open arXiv Scope Settings}\label{open-arxiv-scope-settings}

Choose \textbf{Setting \textgreater{} arXiv Scope Settings}. The window lists every supported primary subject and its archive groups. Select at least one group, then choose \textbf{Save}.

\subsection{Understand the hierarchy}\label{understand-the-hierarchy}

The first level identifies a broad arXiv subject such as Physics, Computer Science, Mathematics, Economics, Electrical Engineering and Systems Science, Quantitative Biology, Quantitative Finance, or Statistics. Each checkbox represents an archive group beneath that subject. For example, \textbf{Physics \textgreater{} Condensed Matter} includes all \texttt{cond-mat.*} leaf categories.

You may select several groups under Physics, or combine groups across primary subjects. Categories shared by more than one scope are fetched only once.

\subsection{Default and saved behavior}\label{default-and-saved-behavior}

\begin{itemize}
\tightlist
\item
  The default selection is only \textbf{Physics \textgreater{} Condensed Matter}.
\item
  \textbf{Restore Defaults} changes the selection in the window; choose \textbf{Save} to apply it.
\item
  Saving changes future arXiv refresh coverage and rebuilds the available \textbf{Preprint Filter} categories.
\item
  Previously saved papers remain in SQLite. Changing scopes never deletes earlier discoveries.
\end{itemize}

\begin{quote}
\textbf{Note}

The category checkboxes in Journal Monitor are display filters. They do not narrow the discovery scope selected here.
\end{quote}

\begin{GewuRelated}
\textbf{See also:} \hyperref[journal-monitor]{Journal Monitor} · \hyperref[field-settings]{Field Settings} · \hyperref[citation-settings]{Open Source Settings} · \hyperref[network-access]{Publisher and Network Access}
\end{GewuRelated}

\section{Field Settings}\label{field-settings}

\begin{GewuLead}
Choose which research fields appear as separate filters across the workspaces.
\end{GewuLead}

\par\noindent\begin{minipage}{\linewidth}

\begin{center}\includegraphics[width=0.96\linewidth,height=0.72\textheight,keepaspectratio]{/Users/lingyuan/PaperMonitor-20260516/assets/user-guide/screenshots/field-settings.png}\end{center}

\begin{GewuCaption}
Field Settings. Choose which research fields appear as separate filters. Unchecked fields remain searchable under Other Fields.
\end{GewuCaption}

\end{minipage}\par\medskip

\subsection{Open Field Settings}\label{open-field-settings}

Choose \textbf{Setting \textgreater{} Field Settings}. Select a field to show it as a separate filter and label. Clear it to collect that field under \textbf{Other Fields}, then choose \textbf{Save}. Closing without saving leaves the earlier presentation in effect.

\subsection{Default organization}\label{default-organization}

{\def\LTcaptype{none} % do not increment counter
\GewuTableStart
\begin{longtable}[]{@{}
  >{\raggedright\arraybackslash}p{(\linewidth - 2\tabcolsep) * \real{0.5000}}
  >{\raggedright\arraybackslash}p{(\linewidth - 2\tabcolsep) * \real{0.5000}}@{}}
\toprule\noalign{}
\rowcolor{GewuTableHeader}
\begin{minipage}[b]{\linewidth}\raggedright
Shown separately by default
\end{minipage} & \begin{minipage}[b]{\linewidth}\raggedright
Grouped under Other Fields by default
\end{minipage} \\
\midrule\noalign{}
\endhead
\bottomrule\noalign{}
\endlastfoot
\begin{minipage}[t]{\linewidth}\raggedright
Physics\\
Quantum Information\\
Materials Science\\
Chemistry\strut
\end{minipage} & \begin{minipage}[t]{\linewidth}\raggedright
Astronomy \& Astrophysics\\
Particle \& Nuclear Physics\\
Mathematics\\
Computer Science\\
Biology \& Medicine\\
Statistics\\
Economics \& Finance\\
Electrical Engineering\\
Earth \& Environmental Science\\
Social Sciences\\
Unclassified records\strut
\end{minipage} \\
\end{longtable}
\GewuTableEnd
}

\subsection{How the visible filters behave}\label{how-the-visible-filters-behave}

\begin{itemize}
\tightlist
\item
  \textbf{Other Fields} appears immediately above \textbf{All Fields}.
\item
  \textbf{All Fields} is always the final choice and includes every classified paper.
\item
  Changing this setting updates field filters, counts, badges, and Paper Details across Journal Monitor, Subject Monitor, Library, My Publications, and Author Monitor.
\item
  Canonical field assignments stay in SQLite and reappear when a field is selected again.
\end{itemize}

\begin{quote}
\textbf{Note}

Use \textbf{Restore Defaults} to save the four-field default immediately. Field Settings do not change arXiv discovery coverage; configure that separately in \textbf{arXiv Scope Settings}.
\end{quote}

\begin{GewuRelated}
\textbf{See also:} \hyperref[arxiv-scope-settings]{arXiv Scope Settings} · \hyperref[citation-settings]{Open Source Settings} · \hyperref[journal-monitor]{Journal Monitor} · \hyperref[subject-monitor]{Subject Monitor}
\end{GewuRelated}

\section{Quick Look Settings}\label{quick-look-settings}

\begin{GewuLead}
Choose whether Space-bar previews stay inside Gewu Insight, use macOS Quick Look, or divide the work by file format.
\end{GewuLead}

\par\noindent\begin{minipage}{\linewidth}

\begin{center}\includegraphics[width=0.96\linewidth,height=0.72\textheight,keepaspectratio]{/Users/lingyuan/PaperMonitor-20260516/assets/user-guide/screenshots/quick-look-settings.png}\end{center}

\begin{GewuCaption}
Quick Look Settings. Choose built-in previews for all files, a mixed PDF/native workflow, or macOS Quick Look for all supported files.
\end{GewuCaption}

\end{minipage}\par\medskip

\subsection{Open Quick Look Settings}\label{open-quick-look-settings}

Choose \textbf{Setting \textgreater{} Quick Look Settings}. Select a mode and choose \textbf{Save}. The setting applies to Space-bar previews in Library, Spaces, Paper Files, and the other research views when a local file is available. Closing without saving leaves the previous mode in effect. \textbf{Restore Default} immediately saves Gewu Insight for all files. The saved mode persists after restarting the App; it does not change Open in Preview. If a preview is already open, saving a different mode refreshes that preview with the newly selected behavior.

{\def\LTcaptype{none} % do not increment counter
\GewuTableStart
\begin{longtable}[]{@{}>{\RaggedRight\arraybackslash}p{0.27\textwidth}>{\RaggedRight\arraybackslash}p{0.63\textwidth}@{}}
\toprule\noalign{}
\rowcolor{GewuTableHeader}
Mode & Behavior \\
\midrule\noalign{}
\endhead
\bottomrule\noalign{}
\endlastfoot
\textbf{Gewu Insight for all files} & The default. Every preview stays inside the workspace. PDFs use the built-in reader; unsupported formats keep the built-in file card and \textbf{Open File} action. \\
\textbf{Gewu Insight for PDFs, macOS for other files} & PDFs use the built-in reader. Presentations, media, images, and other formats use the macOS Quick Look panel. \\
\textbf{macOS Quick Look for all files} & PDFs and every other format supported by macOS use the system Quick Look panel. \\
\end{longtable}
\GewuTableEnd
}

\subsection{Choose a mode for your workflow}\label{choose-a-mode-for-your-workflow}

\begin{itemize}
\tightlist
\item
  Use the default when you want one consistent, contained workspace with an embedded PDF reader.
\item
  Use the mixed mode when PDFs are your main reading material but presentations, video, audio, and images should use native macOS preview support.
\item
  Use the all-native mode for the macOS Quick Look panel, including when you prefer its behavior during rapid PDF switching. In the native App, this is separate from opening the full Preview application.
\end{itemize}

\begin{quote}
\textbf{Note}

When the macOS panel is used, keyboard focus remains in Gewu Insight so \texttt{{\symbolfont ↑}} and \texttt{{\symbolfont ↓}} continue changing the selected Library or Paper Files item instead of turning into page-navigation commands in the system panel.
\end{quote}

\begin{GewuRelated}
\textbf{See also:} \hyperref[quick-look]{Preview Files with Quick Look} · \hyperref[paper-files]{Understand Paper Files} · \hyperref[navigation]{Navigate with Mouse and Keyboard}
\end{GewuRelated}

\section{Website Links}\label{website-link-settings}

\begin{GewuLead}
Choose one browser preference for website actions across your research workspaces.
\end{GewuLead}

\par\noindent\begin{minipage}{\linewidth}

\begin{center}\includegraphics[width=0.96\linewidth,height=0.72\textheight,keepaspectratio]{/Users/lingyuan/PaperMonitor-20260516/assets/user-guide/screenshots/website-links.png}\end{center}

\begin{GewuCaption}
Website Links. Choose the browser used by Open Website and the initial Details state of Save to Gewu. The checkbox shown here is a saved user preference; the shipped Details default is collapsed.
\end{GewuCaption}

\end{minipage}\par\medskip

Open \textbf{Setting \textgreater{} Website Links}, select an option, and click \textbf{Save}. The choice applies to the next website action and is remembered after restarting Gewu Insight. Closing the dialog without saving leaves the previous preference in effect.

{\def\LTcaptype{none} % do not increment counter
\GewuTableStart
\begin{longtable}[]{@{}>{\RaggedRight\arraybackslash}p{0.27\textwidth}>{\RaggedRight\arraybackslash}p{0.63\textwidth}@{}}
\toprule\noalign{}
\rowcolor{GewuTableHeader}
Option & Behavior \\
\midrule\noalign{}
\endhead
\bottomrule\noalign{}
\endlastfoot
\textbf{Gewu Insight Browser} (default) & Opens a new built-in Browser tab, preserving existing pages. Press \texttt{{\symbolfont ⌘}}\texttt{{[}} to return to the source workspace. \\
\textbf{System Default Browser} & Opens the website in the default web browser selected in macOS. Your current Gewu workspace stays in place. \\
\end{longtable}
\GewuTableEnd
}

This applies to Library (including Spaces), My Publications, Journal Monitor, Subject Monitor, and Author Monitor: \textbf{Open Website}, publisher arrows, DOI/arXiv links, Google Scholar, author profiles and homepages, citation sources, remote PDF links, Overleaf, and links inside Note. \textbf{Restore Default} saves Gewu Insight Browser and collapses Save to Gewu Details by default.

\textbf{Show Details by default} controls how Save to Gewu opens. It is off initially; Save remembers it after restart. Toggling Details inside one dialog preserves its fields and does not change this preference. Unsaved settings changes are discarded.

Local PDFs still open in Preview; Finder, Quick Look, and file downloads retain their own actions. Browsing directly inside the Browser tab stays there. Its explicit \textbf{Open in Safari} command still opens Safari. Selecting a different browser does not transfer website sign-ins, cookies, or existing tabs. When using the local UI in a separate web browser, the built-in option opens an ordinary tab in that browser.

\begin{GewuRelated}
\textbf{See also:} \hyperref[browser]{Browse, Bookmark, and Revisit Websites} · \hyperref[paper-details]{Use Paper Details} · \hyperref[quick-look-settings]{Quick Look Settings}
\end{GewuRelated}

\section{AI Access}\label{ai-access}

\begin{GewuLead}
Give an AI client access to selected papers through their stable Gewu Keys, without navigating the interface.
\end{GewuLead}

\par\noindent\begin{minipage}{\linewidth}

\begin{center}\includegraphics[width=0.96\linewidth,height=0.72\textheight,keepaspectratio]{/Users/lingyuan/PaperMonitor-20260516/assets/user-guide/screenshots/ai-access.png}\end{center}

\begin{GewuCaption}
AI Access. The enable switch, existing connections and Recent Activity share a compact settings panel. Expand New Connection to choose a scope and permissions.
\end{GewuCaption}

\end{minipage}\par\medskip

\textbf{Setting \textgreater{} AI Access} is off by default. Normal App use needs no AI token. Enabling it does not start a model, analyze papers or upload your library. Requests happen only when a connected AI client calls the local service.

\subsection{API, MCP and your token}\label{api-mcp-and-your-token}

{\def\LTcaptype{none} % do not increment counter
\GewuTableStart
\begin{longtable}[]{@{}>{\RaggedRight\arraybackslash}p{0.27\textwidth}>{\RaggedRight\arraybackslash}p{0.63\textwidth}@{}}
\toprule\noalign{}
\rowcolor{GewuTableHeader}
Name & Meaning in this connection \\
\midrule\noalign{}
\endhead
\bottomrule\noalign{}
\endlastfoot
API & \textbf{Application Programming Interface}: the operations another program can request from Gewu Insight. \\
MCP & \textbf{Model Context Protocol}: a shared way for an AI client to discover and call tools. Gewu\textquotesingle s local MCP adapter translates those calls into its API requests. \\
Access token & A private credential issued by Gewu Insight. It authorizes a connection\textquotesingle s paper collection and permissions; it is separate from your AI account login and from text tokens counted by models. \\
Gewu Key & The stable identifier of a paper. It selects the article; the access token authorizes access to it. \\
\end{longtable}
\GewuTableEnd
}

\subsection{Methodology: a stateful information infrastructure}\label{methodology-a-stateful-information-infrastructure}

Gewu Insight gives people and AI clients a shared, persistent research workspace. \textbf{Stateful} means that cataloged papers, their identities, registered files and saved analyses remain available between tasks. SQLite is the source of truth for records and relationships; the documents remain local files. The Gewu Key lets a person or an authorized AI return to the same paper across conversations and clients.

The API exposes supported operations with defined inputs and structured results. An AI can retrieve a paper by key without navigating the interface or depending on the internal database schema. MCP makes those operations discoverable as tools; a script can use the HTTP API directly. Both use the same underlying records and authorization rules. Gewu supplies the material and stores approved results; the connected AI performs the analysis.

{\def\LTcaptype{none} % do not increment counter
\GewuTableStart
\begin{longtable}[]{@{}>{\RaggedRight\arraybackslash}p{0.27\textwidth}>{\RaggedRight\arraybackslash}p{0.63\textwidth}@{}}
\toprule\noalign{}
\rowcolor{GewuTableHeader}
Design choice & How it helps research work \\
\midrule\noalign{}
\endhead
\bottomrule\noalign{}
\endlastfoot
Stable paper identity & Use one Gewu Key to connect the article, its files and analyses across tasks. A title or filename is not the identity, and knowing a key does not grant access. \\
Explicit, bounded operations & Search with pagination and read selected pages or text chunks. Programs can compose these steps and limit unnecessary retrieval without relying on screen coordinates. \\
Persistent, traceable results & Save a Markdown analysis under the same key with source file references. People can inspect it in Paper Files, and another authorized session can retrieve it later. \\
\end{longtable}
\GewuTableEnd
}

This provides a foundation for continuing literature analysis and reusing saved knowledge. It does not automatically preserve an AI\textquotesingle s conversation or reasoning, build a semantic memory index, or synchronize model memory. A later session must explicitly retrieve the records and notes it needs; access still depends on its current grant and the continued availability of those records and files.

\subsection{Two API routes, one authorization boundary}\label{two-api-routes-one-authorization-boundary}

{\def\LTcaptype{none} % do not increment counter
\GewuTableStart
\begin{longtable}[]{@{}>{\RaggedRight\arraybackslash}p{0.27\textwidth}>{\RaggedRight\arraybackslash}p{0.63\textwidth}@{}}
\toprule\noalign{}
\rowcolor{GewuTableHeader}
Route & Where permission is checked \\
\midrule\noalign{}
\endhead
\bottomrule\noalign{}
\endlastfoot
AI client {\symbolfont →} MCP adapter {\symbolfont →} Gewu API & The adapter sends the connection token to the API. MCP does not grant additional access. \\
AI client or script {\symbolfont →} Gewu API & The caller supplies the same kind of token directly. Skipping MCP does not skip API authorization. \\
\end{longtable}
\GewuTableEnd
}

For either route, the API checks that AI access is enabled, the token is valid, unexpired and unrevoked, the requested paper is in the connection\textquotesingle s collection, and the operation is permitted. Metadata access alone does not authorize file reading or analysis writing. The native App uses a separate private workspace session for normal human use; an AI token cannot call those internal UI endpoints.

\begin{quote}
\textbf{Interface authorization and local file permissions}

AI Access controls requests through Gewu\textquotesingle s interfaces. It does not encrypt or isolate SQLite and document files from software that already has operating-system permission to read them. Such software can read the disk without contacting MCP or the API; administrator privileges are not necessarily required. Turning off AI Access or revoking a token stops access through that grant, but does not revoke independent filesystem permissions. This broader local access is outside the AI Access permission boundary.
\end{quote}

\subsection{Create a connection}\label{create-a-connection}

\begin{enumerate}
\tightlist
\item
  Keep Gewu Insight running and enable \textbf{AI access}.
\item
  Expand \textbf{New Connection} and enter a name you will recognize.
\item
  Choose Library, Spaces and My Publications; one folder and its subfolders using the usual folder picker; or specific Gewu Keys. A folder connection follows current membership.
\item
  Article metadata and file names are readable by default. Separately allow file contents and/or adding new Markdown analyses when your task needs them.
\item
  Choose an expiry, then \textbf{Create Connection}. The default is 30 days.
\item
  Use \textbf{Copy MCP Configuration} to obtain the local adapter path, Python executable, port and token. Enter these in your AI client\textquotesingle s MCP settings using its required format. See the Codex example below.
\end{enumerate}

\begin{quote}
\textbf{Keep the connection private}

The token is shown only after creation. Store it in client configuration, not in a chat message or shared document. If it is lost, revoke that connection and create another. An AI client may send authorized content to its model provider.
\end{quote}

\par\noindent\begin{minipage}{\linewidth}

\begin{center}\includegraphics[width=0.96\linewidth,height=0.72\textheight,keepaspectratio]{/Users/lingyuan/PaperMonitor-20260516/assets/user-guide/screenshots/ai-new-connection.png}\end{center}

\begin{GewuCaption}
New AI connection. Set the connection name, accessible papers, file-reading permission, analysis-writing permission and expiry before Create Connection. This example leaves optional permissions unchecked.
\end{GewuCaption}

\end{minipage}\par\medskip

\subsection{What an authorized AI can do}\label{what-an-authorized-ai-can-do}

{\def\LTcaptype{none} % do not increment counter
\GewuTableStart
\begin{longtable}[]{@{}>{\RaggedRight\arraybackslash}p{0.27\textwidth}>{\RaggedRight\arraybackslash}p{0.63\textwidth}@{}}
\toprule\noalign{}
\rowcolor{GewuTableHeader}
Tool & Result \\
\midrule\noalign{}
\endhead
\bottomrule\noalign{}
\endlastfoot
\texttt{search\_papers} & Search authorized titles, DOI, Gewu Keys and abstracts, with pagination. \\
\texttt{get\_paper} & Get the article metadata, file inventory and Default file ID for a Gewu Key. \\
\texttt{get\_paper\_file} & Get an original PDF or registered attachment as a read-only temporary copy on the local AI client\textquotesingle s machine, with paper metadata, file ID and actual SHA-256. Requires file-reading permission. \\
\texttt{read\_paper} & Extract PDF text by page or read text/Markdown chunks. Use \texttt{get\_paper\_file} for the original file when text extraction is insufficient. \\
\texttt{save\_analysis} & Add a new Markdown analysis to the same paper\textquotesingle s Paper Files when writing analyses is allowed. \\
\end{longtable}
\GewuTableEnd
}

For example, give the AI a Gewu Key, ask it to read the Default PDF, and save a Chinese analysis with source page references. The saved note records its connection, time, key, request ID and supplied source file hash/pages. It remains an ordinary associated Markdown file, available in built-in Quick Look. Returning to Library refreshes the selected paper\textquotesingle s file panel.

Saving an analysis preserves the article identity, existing files, metadata and Default choice. A retry using the same request ID and identical content reuses the saved file. A new analysis needs a new request ID. This API cannot edit, overwrite or delete any existing Paper Files, including earlier AI analyses, delete papers, change folders or alter the Default.

New analyses are registered as Markdown notes: \textbf{File type: Other}, \textbf{Annotation: Unannotated}, and \textbf{Source: Other} with \textbf{AI analysis} and the connection name as source detail. File size, content hash, creation time and analysis provenance are recorded. These categories are assigned by Gewu; the AI cannot relabel its note as a published paper.

\subsection{Give the original PDF to the AI client}\label{give-the-original-pdf-to-the-ai-client}

Use \texttt{get\_paper\_file} with the Gewu Key and a file ID from \texttt{get\_paper}, or \texttt{default}. When no Default exists, choose an explicit file ID. Gewu checks the connection and transfers the original bytes. The client receives a read-only temporary copy, title, authors, DOI, file classification and actual SHA-256. Its own tools can read the PDF, inspect pages and figures, or run OCR. Gewu does not render pages, run OCR, upload the file to a model or fetch a replacement from the DOI.

The delivered copy is separate from Paper Files and never syncs changes back. Write any rendered pages or other intermediate results separately. Copies last until the MCP process exits normally; finish reading before restarting the connection. A file is limited to 64 MiB, with up to 256 MiB or 64 distinct copies per MCP session. Repeat delivery of the same file version reuses its temporary destination. A client needs access to the same local filesystem as its STDIO adapter; a returned path is not automatically a cloud PDF attachment. If the client cannot access it, report the limitation rather than reading the source Library directly.

Revocation, expiry or disabling AI Access prevents new deliveries. Copies already delivered to a client cannot be recalled. File transfer uses local disk and memory; PDF processing done by a client on the same computer can still compete for resources. Gewu adds no automatic whole-document OCR workload.

\subsection{A repeatable research workflow}\label{a-repeatable-research-workflow}

\begin{enumerate}
\tightlist
\item
  \textbf{Identify.} Give the AI an authorized Gewu Key, or use \texttt{search\_papers} to find candidates and confirm the intended paper with \texttt{get\_paper}.
\item
  \textbf{Read.} Inspect the file inventory, then use \texttt{get\_paper\_file} to obtain the original PDF and read its temporary copy with the AI client\textquotesingle s tools. For text-only requests, \texttt{read\_paper} remains available. Retain the returned file ID, SHA-256 and page references.
\item
  \textbf{Analyze and save.} Ask the AI to produce a Chinese Markdown analysis. With analysis-writing permission, call \texttt{save\_analysis} under the same key and include the source references. Reuse the request ID only for an identical retry; changed source bytes require a fresh read before saving.
\item
  \textbf{Review and continue.} Inspect the note in Paper Files. In another task, retrieve that same paper and explicitly read the saved note before continuing. Each session works with the current authorized state.
\end{enumerate}

These steps can be combined into a scripted workflow or repeated over a selected collection. Pagination and bounded reads help control how much material is retrieved; they do not guarantee model accuracy or a particular speed. Review the analysis against its cited sources.

\subsection{Example: connect Codex on macOS}\label{example-connect-codex-on-macos}

This example connects a local Codex client to Gewu Insight on the same Mac. Keep Gewu running. Create a connection named \textbf{Codex} using the steps above. Enable \textbf{Read PDF, attachment and note contents} to read full text; metadata-only access cannot read PDFs. Enable \textbf{Add new Markdown analyses to Paper Files} only when you also want Codex to save analyses.

\begin{enumerate}
\tightlist
\item
  Open the desktop client\textquotesingle s \textbf{Settings \textgreater{} MCP servers \textgreater{} Add server}.
\item
  Name the server \texttt{gewu-insight} and choose \textbf{STDIO}. This starts our adapter as a local process.
\item
  Fill the command, one argument and two environment variables shown below. Use the values from your own \textbf{Copy MCP Configuration}.
\item
  Save, then choose \textbf{Restart} to reload the MCP connection.
\end{enumerate}

These fields follow \href{https://learn.chatgpt.com/docs/extend/mcp?surface=cli}{OpenAI\textquotesingle s MCP configuration documentation} (checked September 9, 2026). The same local Codex host shares MCP configuration across the desktop app, CLI and IDE extension.

{\def\LTcaptype{none} % do not increment counter
\GewuTableStart
\begin{longtable}[]{@{}>{\RaggedRight\arraybackslash}p{0.27\textwidth}>{\RaggedRight\arraybackslash}p{0.63\textwidth}@{}}
\toprule\noalign{}
\rowcolor{GewuTableHeader}
Codex field & Example value \\
\midrule\noalign{}
\endhead
\bottomrule\noalign{}
\endlastfoot
Name & \texttt{gewu-insight} \\
Type & \texttt{STDIO} \\
Command & \texttt{\textless{}GEWU\_ROOT\textgreater{}/.venv/bin/python} \\
Arguments: one entry & \texttt{\textless{}GEWU\_ROOT\textgreater{}/scripts/gewu\_mcp.py} \\
\end{longtable}
\GewuTableEnd
}

\begin{quote}
\textbf{What \textless GEWU\_ROOT\textgreater{} means}

\texttt{\textless{}GEWU\_ROOT\textgreater{}} is a placeholder for the absolute path to your Gewu program folder: the folder containing \texttt{scripts/gewu\_mcp.py}. It is not the Library folder or the \texttt{.app} bundle. This setup does not require one fixed installation location. Replace the entire placeholder, including the angle brackets, before saving; Codex will not expand it automatically. The command example assumes a project-local \texttt{.venv}. If your Python executable is elsewhere, use the complete \texttt{command} from \textbf{Copy MCP Configuration} instead. Its first \texttt{args} entry supplies the actual script path. Keep a path containing spaces as one field. If you move the installation later, update both paths from a newly generated configuration.
\end{quote}

\subsubsection{Fill in Environment variables}\label{fill-in-environment-variables}

Under \textbf{Environment variables}, enter the first row below: the left \textbf{Key} box holds the variable name and the right \textbf{Value} box holds its value. Click \textbf{Add environment variable} to add the second row. Enter each name and value separately, without quotation marks or an equals sign.

{\def\LTcaptype{none} % do not increment counter
\GewuTableStart
\begin{longtable}[]{@{}>{\RaggedRight\arraybackslash}p{0.27\textwidth}>{\RaggedRight\arraybackslash}p{0.63\textwidth}@{}}
\toprule\noalign{}
\rowcolor{GewuTableHeader}
Key (left box) & Value (right box) \\
\midrule\noalign{}
\endhead
\bottomrule\noalign{}
\endlastfoot
\texttt{GEWU\_API\_URL} & \texttt{http://127.0.0.1:8765}, or the exact address under \texttt{env} in your generated configuration if the port differs. \\
\texttt{GEWU\_API\_TOKEN} & Paste the token just created in Gewu Insight using \textbf{Copy Token}. \\
\end{longtable}
\GewuTableEnd
}

Leave \textbf{Environment variable passthrough} empty. It is for forwarding variables already present in the client\textquotesingle s environment; this example supplies both values explicitly above. The API address belongs in these variables while the server type remains \textbf{STDIO}. Copy the token into settings, not the chat.

\subsubsection{Alternative: edit Codex\textquotesingle s configuration file}\label{alternative-edit-codexs-configuration-file}

\textbf{Copy MCP Configuration} currently produces generic JSON. Codex\textquotesingle s file uses TOML, so do not paste that JSON directly into it. If you prefer editing the file, merge this equivalent block into \texttt{\textasciitilde{}/.codex/config.toml}. Replace \texttt{\textless{}GEWU\_ROOT\textgreater{}} as explained above and \texttt{\textless{}GEWU\_TOKEN\textgreater{}} with your copied token; use the generated command and port if they differ. Preserve your other settings; update an existing \texttt{gewu-insight} block instead of adding a duplicate. Keep the quotation marks in this TOML version and reload the connection after saving.

\begin{verbatim}
[mcp_servers.gewu-insight]
command = "<GEWU_ROOT>/.venv/bin/python"
args = ["<GEWU_ROOT>/scripts/gewu_mcp.py"]

[mcp_servers.gewu-insight.env]
GEWU_API_URL = "http://127.0.0.1:8765"
GEWU_API_TOKEN = "<GEWU_TOKEN>"
\end{verbatim}

\subsubsection{Verify a real paper, then request an analysis}\label{verify-a-real-paper-then-request-an-analysis}

Type \texttt{/mcp} in Codex to inspect the connected server. Replace \texttt{YOUR\_GEWU\_KEY} below with the exact key of an authorized paper that has a local PDF, then send:

\begin{quote}
Use the gewu-insight MCP tools to get the paper with Gewu Key YOUR\_GEWU\_KEY. Show its title and Default filename. Use get\_paper\_file to obtain the original PDF, then read the first three pages of its delivered copy with your PDF tools and summarize them in Chinese. Preserve the file ID and SHA-256; do not access the source Library directly.
\end{quote}

This checks both metadata access and PDF reading. When analysis-writing permission is enabled, you can follow with:

\begin{quote}
Read the relevant remaining PDF pages, then save a Chinese Markdown analysis to this same paper through save\_analysis. Include the source file ID, SHA-256 and page references, and use a unique request ID.
\end{quote}

The saved analysis appears as an associated Markdown file under that paper\textquotesingle s Paper Files; the Default PDF stays unchanged. If a tool is missing, check that connection\textquotesingle s permissions and reload MCP. If the service is unavailable, start Gewu; if access is rejected, check the AI switch, expiry, revocation and selected paper collection. A connection that you create later supplies a new token, which must also be updated in Codex.

\subsection{Pause, revoke and review}\label{pause-revoke-and-review}

Turn off \textbf{Enable AI access} to pause all connections. \textbf{Revoke} stops one connection; expired connections also stop working. Re-enabling access resumes connections that are still valid. Expand \textbf{Recent Activity} to review requests and outcomes, and choose \textbf{Refresh} to update the list. Tokens, prompts and full file text are not included in that activity log.

\subsection{Limits and troubleshooting}\label{limits-and-troubleshooting}

The AI client must run on the same computer, and Gewu Insight\textquotesingle s local service must be running. The local MCP adapter supports protocol 2025-11-25 and earlier negotiated versions. There is no remote MCP endpoint in this version. Developers can retrieve the authenticated API v1 specification at \texttt{/api/v1/openapi.json}.

Original file delivery has a 64 MiB file limit and includes every page. PDF text extraction through \texttt{read\_paper} is limited to ten pages per request and a bounded extraction time; text attachments are returned in chunks. Scanned PDFs may need OCR in the AI client; no OCR or model runs automatically in Gewu. Analysis notes are limited to 256 KiB. A busy response means retry later. An unavailable paper or file may be outside the connection\textquotesingle s collection, missing, or no longer accessible.

Source editing, building and temporary tests require no AI token; calling this API against your collection does require one. Configure each client with the collection and operations its task needs. The interface and filesystem permission boundary is explained above.

\begin{GewuRelated}
\textbf{See also:} \hyperref[paper-files]{Understand Paper Files} · \hyperref[quick-look]{Preview Files with Quick Look} · \hyperref[privacy]{Local Data and Privacy}
\end{GewuRelated}

\section{Naming and PDF Rename Rules}\label{naming}

\begin{GewuLead}
Display naming and physical PDF renaming are separate choices. Review metadata before renaming files.
\end{GewuLead}

\par\noindent\begin{minipage}{\linewidth}

\begin{center}\includegraphics[width=0.96\linewidth,height=0.72\textheight,keepaspectratio]{/Users/lingyuan/PaperMonitor-20260516/assets/user-guide/screenshots/pdf-rename-settings.png}\end{center}

\begin{GewuCaption}
PDF Rename Settings. Edit the filename template and inspect its preview before saving. Available tokens appear beside the template.
\end{GewuCaption}

\end{minipage}\par\medskip

\subsection{Choose the Library naming rule}\label{choose-the-library-naming-rule}

Choose \textbf{Setting \textgreater{} Naming Rule \textgreater{} Use Title} or \textbf{Use Citation}. The choice is saved immediately and supplies the default name for a newly added Library placement that has no name yet. It does not rename the existing collection. The default rule is Use Citation; physical PDF renaming uses the separate template below.

\subsection{Configure PDF renaming}\label{configure-pdf-renaming}

Open \textbf{Setting \textgreater{} PDF Rename Settings}. A template can use tokens such as:

\begin{itemize}
\tightlist
\item
  \texttt{\{author\}} or \texttt{\{authors\}}
\item
  \texttt{\{year\}}
\item
  \texttt{\{title\}}
\item
  \texttt{\{journal\}}
\item
  \texttt{\{doi\}}
\end{itemize}

Edit \textbf{Filename Template} and inspect the sample \textbf{Preview}, then choose \textbf{Save} or press Return. \textbf{Default} immediately saves the standard \texttt{\{author\}\_\{year\}\_\{title\}} template. Saving a template does not rename existing PDFs; apply it through the Library rename action below.

\subsection{Rename one item from metadata}\label{rename-one-item-from-metadata}

\begin{enumerate}
\tightlist
\item
  \textbf{Confirm its DOI, title, author, journal, and year.}
\item
  \textbf{Control-click the item.}
\item
  \textbf{Choose \textbf{Auto rename from Metadata}.}
\item
  \textbf{Verify the result in the Library and Finder.} To rename several PDFs, select only local PDF paper rows, Control-click, and choose \textbf{Auto rename n selected}. Each file uses its own metadata; the progress message reports renamed, already named, and skipped items.
\end{enumerate}

\begin{quote}
\textbf{Note}

Choosing a Default PDF can move it to the main file folder and archive the previous main PDF in Associated Files. This does not apply the filename template or rewrite file contents; name conflicts use a unique filename. Use a naming or rename action to apply your preferred naming rule.
\end{quote}

\begin{GewuRelated}
\textbf{See also:} \hyperref[versions-current]{Choose File Types and the Default File} · \hyperref[organize]{Folders, Shortcuts, and Spaces}
\end{GewuRelated}

\section{How Gewu Insight Stores Files}\label{storage-model}

\begin{GewuLead}
The visible Library stays simple: its main tree contains a compatibility primary path. Other resources can live in a mirrored Associated Files area, while the UI presents every resource as a peer in one flat Paper Files list.
\end{GewuLead}

\subsection{One work, several representations}\label{one-work-several-representations}

Each managed paper has one database identity and a stable \textbf{Gewu Key}. DOI and arXiv are external identifiers attached to that identity; neither is used as the sole storage identity.

\subsection{Where the files go}\label{where-the-files-go}

\begin{itemize}
\tightlist
\item
  A compatibility primary PDF stays in the visible Library, Space, or My Publications tree.
\item
  Older PDF versions go to a mirrored \texttt{Versions} area.
\item
  Supporting information goes to \texttt{Supplements}.
\item
  Images, notes, presentations, data, code, and other attachments go to \texttt{Files}.
\end{itemize}

The associated bundle mirrors the paper's managed path under \nolinkurl{data/Associated Files}. When you move the paper in the UI, Gewu Insight moves the associated bundle and refreshes live paths.

\subsection{Why the UI and Finder look different}\label{why-the-ui-and-finder-look-different}

The UI presents a paper as one item with a flat Paper Files list inside it. Finder shows a compatibility primary file in the normal tree and the related bundle in Associated Files. The UI Default marker is a reading preference and does not require physical file promotion or renaming.

\subsection{Copies and filesystem shortcuts}\label{copies-and-filesystem-shortcuts}

When the same work appears in more than one Library folder, those locations can share one database record and one canonical managed file. Symbolic links can represent secondary filesystem locations. \textbf{Make Copy Desktop} resolves those shortcuts into real files for a portable export.

\begin{quote}
\textbf{Important}

Do not manually move an Associated Files bundle or replace an alias in Finder. Use Gewu Insight so the database and physical paths change together.
\end{quote}

\begin{GewuRelated}
\textbf{See also:} \hyperref[storage-reference]{Storage Reference} · \hyperref[desktop-copy]{Make a Portable Desktop Copy} · \hyperref[paper-files]{Understand Paper Files}
\end{GewuRelated}

\section{Local Data and Privacy}\label{privacy}

\begin{GewuLead}
Gewu Insight runs a local interface and keeps its authoritative library database on your Mac. Network requests occur only for the features that need external metadata, pages, profiles, or files.
\end{GewuLead}

\subsection{Local components}\label{local-components}

\begin{itemize}
\tightlist
\item
  The UI server listens on the local loopback address, normally \texttt{127.0.0.1:8765}.
\item
  SQLite stores papers, identities, relationships, settings, and operational state.
\item
  Browser bookmarks, folders, history, and sidebar preferences also live in SQLite. Browser cookies and cache use a separate persistent WebKit profile; they are not shared with Safari or exported with the database.
\item
  Managed PDFs, attachments, photos, and generated reports stay in local project directories.
\item
  JSON files are exports or diagnostics; they are not the authoritative database.
\end{itemize}

\subsection{External requests}\label{external-requests}

Metadata refresh, citation counts, journal/subject/author discovery, profile images, website opening, and downloads can contact their named services. The app paces services that require strict request limits.

\subsection{Protect credentials}\label{protect-credentials}

Optional service credentials belong only in Open Source Settings or approved environment variables. They must not appear in a screenshot, support message, Git commit, or exported diagnostic file.

\begin{GewuRelated}
\textbf{See also:} \hyperref[citation-settings]{Open Source Settings} · \hyperref[backup-migrate]{Back Up or Move to Another Mac} · \hyperref[network-access]{Publisher and Network Access}
\end{GewuRelated}

\part{Product Identity}\label{part-4}

\GewuIdentityChapter{Visual Identity and Cultural Elements}\label{visual-identity}

\begin{GewuLead}
Gewu Insight uses three related marks---the calligraphic Chinese name, the app icon, and a seal-script watermark---to connect a contemporary research workspace with the longer Chinese tradition of examining things in order to form insight.
\end{GewuLead}

The marks do not decorate every surface. Each has a different job: the icon identifies the application, the calligraphic name gives the product a human signature, and the quiet seal-script pattern belongs to transitional or otherwise empty space.

\par\noindent\begin{minipage}{\linewidth}

\begin{center}\includegraphics[width=0.48\linewidth,height=0.72\textheight,keepaspectratio]{/Users/lingyuan/PaperMonitor-20260516/assets/user-guide/screenshots/about-window.png}\end{center}

\begin{GewuCaption}
About Gewu Insight. Open the application menu and choose About Gewu Insight to see the product name, version and calligraphic signature.
\end{GewuCaption}

\end{minipage}\par\medskip

\GewuIdentityChapter{The calligraphic name}\label{the-calligraphic-name}

\begin{center}\includegraphics[width=0.84\linewidth,keepaspectratio]{/Users/lingyuan/PaperMonitor-20260516/assets/branding/GewuAboutCalligraphy.png}\end{center}

\begin{GewuCaption}
The transparent production artwork used in the About panel. From left to right the picture contains 見、洞、物、格; read in the traditional direction from right to left, it says 格物洞見.
\end{GewuCaption}

This is a \textbf{祝枝山行草集字} composition: four separately documented historical character forms were assembled into one wordmark. It is therefore accurate to call it ``a collection of characters in the calligraphic style of Zhu Zhishan,'' but not ``a title written by Zhu Zhishan.'' No surviving single inscription titled {格物洞見} is being claimed.

\section{How the four characters were assembled}\label{how-the-four-characters-were-assembled}

{\def\LTcaptype{none} % do not increment counter
\GewuTableStart
\begin{longtable}[]{@{}>{\RaggedRight\arraybackslash}p{0.14\textwidth}>{\RaggedRight\arraybackslash}p{0.20\textwidth}>{\RaggedRight\arraybackslash}p{0.22\textwidth}>{\RaggedRight\arraybackslash}p{0.24\textwidth}@{}}
\toprule\noalign{}
\rowcolor{GewuTableHeader}
Reading order & Reference character & Recorded hand and script & Role in the product artwork \\
\midrule\noalign{}
\endhead
\bottomrule\noalign{}
\endlastfoot
1 & \href{https://www.shufa5.com/images/shufa/zi/5/35045.gif}{格} & 祝枝山 · 行书 & Begins the product name when read from the right. \\
2 & \href{https://www.shufa5.com/images/shufa/zi/6/165906.gif}{物} & 祝枝山 · 草书 & Completes the classical term 格物. \\
3 & \href{https://www.shufa5.com/images/shufa/zi/2/162202.gif}{洞} & 祝枝山 · 草书 & Introduces 洞見, insight gained by seeing through a subject. \\
4 & \href{https://www.shufa5.com/images/shufa/zi/2/179662.gif}{見} & 祝枝山 · 草书 & Closes the traditional form of the name. \\
\end{longtable}
\GewuTableEnd
}

The selected glyphs were optically balanced into a common line: scale, spacing, ink tone, edge cleanup, and transparent alpha were normalized while the characteristic strokes were retained. The production preparation used a Codex-assisted image workflow. The interface and accessibility text use the contemporary simplified product name \textbf{格物洞见}; the artwork deliberately preserves the traditional character \textbf{見} and traditional right-to-left composition.

\begin{quote}
\textbf{Provenance and reuse}

The character links above are retained as provenance for the design process. The source thumbnails do not state a clear open-reuse licence. They should not be treated as public-domain assets merely because they are viewable online. The assembled Gewu Insight product artwork is distributed under the repository\textquotesingle s all-rights-reserved notice.
\end{quote}

\GewuIdentityChapter{The app icon and the seal-script watermark}\label{the-app-icon-and-the-seal-script-watermark}

\section{The app icon}\label{the-app-icon}

\begin{center}\includegraphics[width=0.18\linewidth,keepaspectratio]{/Users/lingyuan/PaperMonitor-20260516/assets/gewu-insight-app-icon.png}\end{center}

The icon\textquotesingle s two blue-violet fields read simultaneously as an open book and two sides of a field of view. The white eye and central point suggest observation and focus; the tapered center also recalls a pen nib or the spine of a book.

\begin{GewuCaption}
This description is the product\textquotesingle s visual reading of the finished mark, not a claim that every interpretation was specified in an original design brief.
\end{GewuCaption}

\section{The seal-script watermark}\label{the-seal-script-watermark}

\begin{center}\begin{tcolorbox}[enhanced,width=0.30\linewidth,colback=GewuSealPanel,colframe=GewuSealPanel,boxrule=0pt,arc=3mm,boxsep=0pt,left=4mm,right=4mm,top=4mm,bottom=4mm]\includegraphics[width=\linewidth,keepaspectratio]{/Users/lingyuan/PaperMonitor-20260516/assets/branding/GewuSealGlyphs.png}\end{tcolorbox}\end{center}

The 2×2 motif is a digital rendering in \href{https://www.hanyi.com.cn/productdetail?id=450&type=0}{汉仪篆书繁}. Read down the right column and then down the left, its four characters are 格、物、洞、见. It uses simplified 见 because that is the supported glyph in the selected digital typeface.

\begin{GewuCaption}
The mark is a light-line screen motif inspired by seal composition, not a photograph or facsimile of a historical carved seal. Typeface licensing remains separate from the product artwork.
\end{GewuCaption}

\GewuIdentityChapter{Why the three marks are used differently}\label{why-the-three-marks-are-used-differently}

{\def\LTcaptype{none} % do not increment counter
\GewuTableStart
\begin{longtable}[]{@{}>{\RaggedRight\arraybackslash}p{0.27\textwidth}>{\RaggedRight\arraybackslash}p{0.63\textwidth}@{}}
\toprule\noalign{}
\rowcolor{GewuTableHeader}
Mark & Role \\
\midrule\noalign{}
\endhead
\bottomrule\noalign{}
\endlastfoot
\textbf{Icon · recognition} & {Use at application scale: Dock, startup identity, About, and platform packaging.} \\
\textbf{Calligraphy · signature} & {Use where the Chinese name can be read deliberately, especially the About panel and identity materials.} \\
\textbf{Seal script · atmosphere} & {Use at very low opacity in startup transitions or quiet lower-right space, never where it competes with paper details.} \\
\end{longtable}
\GewuTableEnd
}

Together the three registers establish the product\textquotesingle s visual identity: the icon is contemporary and operational; the calligraphy is personal and gestural; the seal-script watermark is historical in reference but restrained in use. Their common purpose is to make the name visible without turning the research workspace into a themed illustration.

\begin{GewuRelated}
\textbf{See also:} \hyperref[welcome]{Welcome} · \hyperref[personal-research-workspace]{The Personal Research Workspace} · \hyperref[startup]{Startup, Monitor Cache, and Text Size}
\end{GewuRelated}

\part{Technical and Operational Reference}\label{part-5}

\chapter{Keyboard Shortcuts}\label{keyboard}

\begin{GewuLead}
Use standard macOS editing commands throughout the app. The shortcuts below are specific to navigation and Gewu Insight commands.
\end{GewuLead}

{\def\LTcaptype{none} % do not increment counter
\GewuTableStart
\begin{longtable}[]{@{}>{\RaggedRight\arraybackslash}p{0.27\textwidth}>{\RaggedRight\arraybackslash}p{0.63\textwidth}@{}}
\toprule\noalign{}
\rowcolor{GewuTableHeader}
Shortcut & Action \\
\midrule\noalign{}
\endhead
\bottomrule\noalign{}
\endlastfoot
\texttt{{\symbolfont ↑}} / \texttt{{\symbolfont ↓}} / \texttt{{\symbolfont ←}} / \texttt{{\symbolfont →}}, Enter & In Browser chrome, move between the top workspace tab, toolbar, webpage tabs, Favorites, and visible sidebar; Enter activates the focused control. Native webpages retain scrolling and editing keys. Up at an ordinary document\textquotesingle s top can return to Browser chrome; {\symbolfont ⌘}L then Escape also returns to the toolbar. \\
\texttt{{\symbolfont ↑}} / \texttt{{\symbolfont ↓}} & Move through the focused Library or Paper Files list; on the source rail, move among Library, shortcuts, and Spaces and open the highlighted source. \\
\texttt{{\symbolfont →}} & Enter the highlighted source from the vertical rail, enter a folder, or move from a selected paper into Paper Files. \\
\texttt{{\symbolfont ←}} & Leave Paper Files for the selected paper, return to a parent folder, or move from the leftmost Library column to the source rail. \\
\texttt{Space} & Open or close Quick Look for the selected local paper/Paper Files item, including available monitor PDFs; not while editing text. \\
\texttt{Esc} & Dismiss an active transient control or return focus where supported. \\
\texttt{{\symbolfont ⇧}}\texttt{{\symbolfont ⌘}}\texttt{V} & Toggle List View and Table View. \\
\texttt{{\symbolfont ⌘}}\texttt{L} & Focus the address/search field in Browser; otherwise open Folder Loader. \\
\texttt{{\symbolfont ⌘}}\texttt{D} & In Browser, bookmark the current page. \\
\texttt{{\symbolfont ⌘}}\texttt{R} & In Browser, reload the current page. \\
\texttt{{\symbolfont ⌘}}\texttt{F} & In Browser, Find on Page. In Library, My Publications, and the three Monitors, focus the current workspace search. \\
\texttt{{\symbolfont ⌘}}\texttt{T} & In Browser, open a new webpage tab at Home. \\
\texttt{{\symbolfont ⌘}}\texttt{{[}} & Return from Browser to the source workspace; keep its position and webpage tabs. \\
\texttt{{\symbolfont ⌘}}\texttt{,} & Show Open Source Settings. \\
\texttt{{\symbolfont ⌘}}\texttt{W} & In Browser, close the selected webpage tab; closing the last creates Home. Outside Browser, dismiss active Quick Look first, otherwise close the current window. \\
\texttt{{\symbolfont ⌘}}\texttt{Q} & Quit Gewu Insight. Save webpage drafts first; open Browser tabs are not restored after quitting. \\
\end{longtable}
\GewuTableEnd
}

\section{Monitor keyboard navigation}\label{monitor-keyboard-navigation}

In a Journal or Subject calendar, Left/Right moves one day and Up/Down moves one week. Enter selects the date and enters the papers. In the paper area, Up/Down selects papers, Left/Right changes the journal or subject, and Escape returns to the calendar. Space opens an available local PDF in Quick Look.

In Author Monitor, Left selects the author rail and Right selects its papers; Up/Down moves through the active pane. Up from the first item enters the main-tab navigation level. At that level, Left/Right chooses a main tab and Down or Enter returns to its content. These commands do not apply while typing in an input.

\section{Browser keyboard context}\label{browser-keyboard-context}

Find on Page uses Enter for the next match, Shift-Enter for the previous match, and Escape to close. In address suggestions, Up/Down selects a result and Escape restores your query. A focused webpage tab uses Left/Right to change tabs. These contexts are separate from Library navigation; {{\symbolfont ⌘}L} in Browser does not open Folder Loader.

\textbf{Paper Files mouse behavior:} single-click selects, double-click opens, and Control-click opens the context menu.

\begin{GewuRelated}
\textbf{See also:} \hyperref[navigation]{Navigate with Mouse and Keyboard} · \hyperref[paper-files]{Understand Paper Files}
\end{GewuRelated}

\chapter{Publisher and Network Access}\label{network-access}

\begin{GewuLead}
Publishers can rate-limit shared VPN addresses or automated-looking traffic. Gewu Insight reduces overlap and request rate, but it cannot remove a publisher block.
\end{GewuLead}

\section{If a publisher page is blocked}\label{if-a-publisher-page-is-blocked}

\begin{enumerate}
\tightlist
\item
  \textbf{Stop repeated refreshes.} Additional requests can extend a temporary block.
\item
  \textbf{Try your direct connection.} If the page works with VPN off, the VPN exit IP is likely the blocked client.
\item
  \textbf{Use \textbf{Open Website} or an external browser.} Complete any normal institutional login or publisher interaction yourself.
\item
  \textbf{Download only when authorized.} Add a legitimately obtained PDF with \textbf{Add File}.
\end{enumerate}

\section{Why figures may arrive later}\label{why-figures-may-arrive-later}

Figure and article-page retrieval is queued behind discovery work and uses caching. This reduces simultaneous RSS, page, and image requests to the same publisher. A paper can appear before its representative image.

\section{Browser and your system network}\label{browser-and-your-system-network}

The Browser tab uses native WebKit requests and does not proxy pages through the Python server. System proxy or TUN capture can send its traffic into Clash; only then do the active rules determine DIRECT or proxy routing. A DIRECT rule uses your current physical network, which must itself provide institutional access. Gewu Insight neither changes these rules nor guarantees a particular route; use Clash connection records to verify the actual route. See \hyperref[browser]{Browse, Bookmark, and Revisit Websites} for downloads, sign-in limits, and the separate browser profile.

\section{Service rate limits}\label{service-rate-limits}

Semantic Scholar, Google Scholar, and some publisher sources use shared pacing, cooldowns, or exponential backoff. Let the bottom status area settle before starting the same action again.

\begin{quote}
\textbf{Important}

Do not attempt to bypass a paywall, CAPTCHA, browser verification, or access control. Use your authorized browser session or attach a file you obtained legitimately.
\end{quote}

\begin{GewuRelated}
\textbf{See also:} \hyperref[journal-monitor]{Journal Monitor} · \hyperref[download-sync]{Download and Sync Versions} · \hyperref[troubleshooting]{Troubleshooting}
\end{GewuRelated}

\chapter{Startup, Monitor Cache, and Text Size}\label{startup}

\begin{GewuLead}
The macOS app shows its branded startup view immediately, starts or connects to the local UI service, and reveals the workspace when Library is ready.
\end{GewuLead}

\par\noindent\begin{minipage}{\linewidth}

\begin{center}\includegraphics[width=0.96\linewidth,height=0.72\textheight,keepaspectratio]{/Users/lingyuan/PaperMonitor-20260516/assets/user-guide/screenshots/font-size.png}\end{center}

\begin{GewuCaption}
Font Size. Preview one of five reading sizes. Save keeps the selected size; Cancel restores the previous preference.
\end{GewuCaption}

\end{minipage}\par\medskip

\section{Normal startup}\label{normal-startup}

The first window frame contains the warm background, packaged icon, and Starting Gewu Insight panel. The native startup view is installed before the window is shown, so an empty white web background does not precede the branding. The panel stays visible while the local service and first Library view prepare, then fades into the workspace.

The first launch after an additive database upgrade can take longer because the existing SQLite records are migrated and older provenance or file metadata is backfilled. The native launcher allows up to 120 seconds for this preparation. Keep the single startup panel open rather than launching a second copy; the migration preserves the existing database and managed files.

The Dock tile, native startup panel, and in-window loading state use the icon packaged with the same app build. Startup has a native launch phase and a local-web bootstrap phase, but the visual handoff should not introduce a second icon frame or fall back to an older product mark.

\section{Automatic local-server recovery}\label{automatic-local-server-recovery}

If the local UI server exits unexpectedly while the macOS window remains open, the launcher starts it again. The window rechecks local health and reloads the workspace when the service becomes ready. A startup-failure sheet is nonblocking; if the server recovers after that sheet appears, the stale sheet is dismissed automatically instead of continuing to interrupt a healthy workspace.

A brief reload during recovery does not mean the SQLite database was rebuilt or that another copy of Gewu Insight should be opened. If repeated recovery never reaches a healthy Library view, use the launcher log steps below.

Recovery reloads the workspace, so Browser\textquotesingle s memory-only tabs and unsubmitted forms can be lost. Bookmarks and history already saved in SQLite remain. Save work on websites regularly; recovery is not a web-session or draft backup.

\section{When automatic updates run}\label{when-automatic-updates-run}

Automatic monitoring runs inside the local UI service, with startup delays, periodic checks, and saved completion/retry state. Keep the Mac awake and the service running. Closing a window does not necessarily stop that service. Quitting the App stops a server it started; a separately launched server must be stopped separately. A late launch checks work that is due rather than guaranteeing a run at an exact minute.

{\def\LTcaptype{none} % do not increment counter
\GewuTableStart
\begin{longtable}[]{@{}>{\RaggedRight\arraybackslash}p{0.27\textwidth}>{\RaggedRight\arraybackslash}p{0.63\textwidth}@{}}
\toprule\noalign{}
\rowcolor{GewuTableHeader}
Monitor & Current automatic schedule \\
\midrule\noalign{}
\endhead
\bottomrule\noalign{}
\endlastfoot
Journal Monitor & Daily work becomes due from 09:00 in the Mac\textquotesingle s local time zone; stored completion and pending retry state determine what still needs work. \\
Subject Monitor & Daily work becomes due from 09:00 local time; incomplete sources have their own bounded retry state. \\
Author Monitor & Checks due authors from 10:00 local time; already checked authors are not blindly fetched again on every scheduler tick. \\
arXiv & Uses its separate announcement-ready schedule and buffer, not the journal\textquotesingle s 09:00 rule. Read the displayed update/countdown state. \\
\end{longtable}
\GewuTableEnd
}

These are local scheduler rules, not guarantees of remote source completion or a hosted monitoring service. Manual Refresh remains available for the selected scope; repeated broad refreshes can worsen a source\textquotesingle s rate limit.

\section{If the app cannot start}\label{if-the-app-cannot-start}

\begin{enumerate}
\tightlist
\item
  \textbf{Quit any old Gewu Insight window.} An earlier local server may still own port 8765.
\item
  \textbf{Open the launcher log.} Check \nolinkurl{logs/paper_monitor_launcher.log} for the first error.
\item
  \textbf{Allow the first startup to finish.} After an upgrade, wait through the migration window before treating the preparation panel as a failure.
\item
  \textbf{Start the app again once if recovery does not complete.} Avoid opening several copies while the first is still preparing.
\item
  \textbf{Preserve the database.} Do not delete SQLite as a startup fix.
\end{enumerate}

\section{Clear a stale interface}\label{clear-a-stale-interface}

Choose \textbf{Setting \textgreater{} Clear Monitor Cache} if a monitor still shows stale downloaded images or rendered web assets after an app update. This clears disposable monitor presentation cache, not the SQLite database, saved paper records, or managed PDFs.

This is not Browser\textquotesingle s Clear History command and does not clear Browser sign-in cookies or Apple Passwords. Use \textbf{Browser \textgreater{} History \textgreater{} Clear History} for visit records only.

Monitor cache and macOS app identity are separate. Clearing Monitor Cache cannot repair a stale Dock tile. After replacing the app bundle, quit every running copy and open the rebuilt \textbf{Gewu Insight.app}. If a permanently pinned Dock tile still points to an older bundle at another path, remove that tile and add the current app again.

\section{Choose a comfortable text size}\label{choose-a-comfortable-text-size}

Choose \textbf{Setting \textgreater{} Font Size}, then select \textbf{Compact}, \textbf{Standard}, \textbf{Large}, \textbf{Larger}, or \textbf{Extra Large}. The workspace previews the choice immediately while keeping panel dimensions, icons, and window controls unchanged.

\textbf{Cancel}, the Escape key, or clicking outside the sheet restores the previously saved size. \textbf{Restore Default} previews Standard and still requires Save. Choose \textbf{Save} to keep the preference in the local SQLite settings. New app sessions apply the saved size automatically.

\begin{GewuRelated}
\textbf{See also:} \hyperref[troubleshooting]{Troubleshooting} · \hyperref[privacy]{Local Data and Privacy}
\end{GewuRelated}

\chapter{Troubleshooting}\label{troubleshooting}

\begin{GewuLead}
Start with the visible status and the selected work. Avoid broad repair actions until you know whether the problem is network, metadata, database state, or a local path.
\end{GewuLead}

\section{A PDF cannot be found after a move}\label{a-pdf-cannot-be-found-after-a-move}

\begin{enumerate}
\tightlist
\item
  Reselect the paper and try \textbf{Show in Finder}.
\item
  If the move happened in Gewu Insight, close and reopen the item to refresh the view.
\item
  If the file was moved manually in Finder, return it if possible, then organize it through the app.
\item
  Check Recycle before importing another copy.
\end{enumerate}

\section{Paper Files shows Missing File}\label{paper-files-shows-missing-file}

This label is a path-state warning, not a failed metadata lookup: SQLite still knows the resource, but the recorded file path does not currently exist. Check \textbf{Show in Finder}, the corresponding \nolinkurl{data/Associated Files} bundle, and Recycle. If the file was moved outside Gewu Insight, restore it before reorganizing it through the app. Do not delete the row merely because its filename resembles the default PDF; explicitly typed attachments and the default resource remain visible for recovery.

\section{Quick Look is slow on the first open}\label{quick-look-is-slow-on-the-first-open}

A large or previously unseen PDF may need an initial structural read and thumbnail render. The main document should appear before the full thumbnail strip. Leave the preview open briefly instead of pressing Space repeatedly; later opens can reuse page information and cached thumbnails. If the viewer never produces a page, use \textbf{Open in Preview} to check whether the file itself is valid.

\section{PDF switching flashes in built-in Quick Look}\label{pdf-switching-flashes-in-built-in-quick-look}

The built-in preview can flash as it loads another PDF. Open \textbf{Setting \textgreater{} Quick Look Settings}, choose \textbf{macOS Quick Look for all files}, and Save to use the system panel for rapid comparison. Up and Down still follow the selected files. This changes the preview method; it does not alter your PDF files or default-file choices.

\section{The wrong file is the default}\label{the-wrong-file-is-the-default}

\begin{enumerate}
\tightlist
\item
  Open Paper Files and verify every file type.
\item
  Correct a mislabeled Corrected published version, Published version, Accepted manuscript, or Preprint; expand the size footer to set annotation separately.
\item
  Choose \textbf{Set as Default} for the desired valid file, or choose \textbf{Use Automatic Selection} to resume automatic selection.
\end{enumerate}

\section{A duplicate appears after download}\label{a-duplicate-appears-after-download}

Compare file size and classification. New downloads use exact SHA-256 and PDF content fingerprints. A copy with annotations may be retained intentionally. Exact byte duplicates are reused, while clean and annotated peers can both remain. Keep the intended default before deleting any redundant non-default row.

\section{A monitor shows no papers}\label{a-monitor-shows-no-papers}

Read the source status. An empty successful source is not the same as a failed fetch. In Subject Monitor, allow scheduled source-level retries to finish, or use the selected subject\textquotesingle s refresh command when a manual retry is needed. In Journal Monitor, confirm the date and filters. If a VPN is involved, test a direct connection without repeatedly refreshing.

\section{A journal count and the visible cards disagree}\label{a-journal-count-and-the-visible-cards-disagree}

This is usually expected when \textbf{Field Filter} is active. Journal counts describe all stored papers for the selected journal and day; field choices narrow the visible cards. Clear or broaden the field selection before diagnosing a missing fetch.

\section{A journal disappeared from the configured list}\label{a-journal-disappeared-from-the-configured-list}

Open \textbf{Journal Settings} and inspect the disabled history before editing configuration files. Removing a journal disables its subscription but preserves previously collected papers and history. Add it again to restore monitoring. Because the subscription list is SQLite state, replacing source files should not be used as a journal-list repair.

\section{LaTeX remains visible as source text}\label{latex-remains-visible-as-source-text}

On a cold launch, allow the math renderer a moment to become ready. The active day should re-render automatically. If raw \texttt{\$...\$} text remains, check network availability for the renderer, then reload the same view once; switching calendar dates should not be required.

\section{Folder Loader looks incomplete}\label{folder-loader-looks-incomplete}

Open the live or previous Loader log and the result guide. Confirm the metadata mode, source path, and any stop action. Re-running the same source should reuse existing files rather than overwrite the primary PDF.

\section{A default file cannot be deleted}\label{a-default-file-cannot-be-deleted}

A Default PDF is kept as the protected main file. Choose another PDF with \textbf{Set as Default} to move the previous primary PDF into Associated Files, where it remains available as a separate file. A non-PDF default changes the reading preference and leaves the main PDF in place. Remove ordinary non-primary resources through Paper Files. To remove the managed paper itself, use the normal Library reference/Recycle workflow, which preserves shared references. Do not bypass a primary-file warning in Finder. An \textbf{Invalid file} remains an audit record and is never eligible to become default.

\section{Change where website links open}\label{change-where-website-links-open}

Open \hyperref[website-link-settings]{Setting \textgreater{} Website Links} and save \textbf{System Default Browser} to use the browser selected in macOS. The default, \textbf{Gewu Insight Browser}, opens a new built-in webpage tab or reuses its unused first tab. This covers website actions in all five research views, including publisher arrows, DOI/arXiv links, scholarly profiles, citation-source links, and links inside Note. Existing web pages stay open. Local PDF, Preview, Finder, download, and file-import actions are unchanged.

\section{Home works offline but a website will not load}\label{home-works-offline-but-a-website-will-not-load}

Home displays local Favorites without a Google request. Missing site icons fall back to initials; they do not prevent opening Home. Search results and account actions require a connection. Read the load error, check the destination and Clash route, and try Reload once. A website may reject embedded browsers or have its own session policy; an unsupported-browser banner alone does not prove why a sign-in ended. Open in Safari is a manual compatibility fallback, with a separate sign-in session.

\section{A publisher stays on security verification}\label{a-publisher-stays-on-security-verification}

A verification loop does not establish that you were identified as AI. Browser supports standard internal child documents used by embedded widgets, but a publisher\textquotesingle s verification can still reject an embedded browser. Try Reload once; if the loop persists, use \textbf{Open in Safari}. Safari has a separate session. For publisher support, note the public URL, time, and displayed Ray ID, without sharing cookies or passwords. Do not repeatedly retry or bypass the challenge.

\section{AutoFill does not appear or does not fill}\label{autofill-does-not-appear-or-does-not-fill}

Automatic triggering is deliberately limited to an explicit click in a clearly marked login field. Ordinary email, search, registration, verification-code, and uncertain fields do not qualify. Cancellation uses up the automatic offer for that field kind/site during the App session. Right-click the actual input and choose \textbf{AutoFill \textgreater{} Passwords} to retry manually.

Unlock and select the appropriate entry yourself. If the item says \textbf{Passkey} and shows a creation date, it is not a text password. If a real Password entry closes without filling, macOS or that website may not support the return path in this embedded view. Gewu does not read your vault, relay credentials, or submit forms to work around it. Use the website\textquotesingle s supported login flow or Safari; do not send passwords or unmasked credential screenshots in a support report.

\begin{quote}
\textbf{Important}

Before any manual filesystem repair, quit the app and make a full backup of the database and managed files together.
\end{quote}

\begin{GewuRelated}
\textbf{See also:} \hyperref[backup-migrate]{Back Up or Move to Another Mac} · \hyperref[startup]{Startup, Monitor Cache, and Text Size} · \hyperref[network-access]{Publisher and Network Access}
\end{GewuRelated}

\chapter{Storage Reference}\label{storage-reference}

\begin{GewuLead}
These paths describe the managed project layout. They are useful for backup and diagnosis, not for routine manual reorganization.
\end{GewuLead}

{\def\LTcaptype{none} % do not increment counter
\GewuTableStart
\begin{longtable}[]{@{}>{\RaggedRight\arraybackslash}p{0.27\textwidth}>{\RaggedRight\arraybackslash}p{0.63\textwidth}@{}}
\toprule\noalign{}
\rowcolor{GewuTableHeader}
Location & Purpose \\
\midrule\noalign{}
\endhead
\bottomrule\noalign{}
\endlastfoot
\nolinkurl{data/paper_monitor.sqlite} & {Authoritative database for papers, identities, relationships, settings, operational state, Browser bookmarks, and visit history. May contain saved metadata API keys; treat backups as private.} \\
\nolinkurl{data/cache/journal_logos} & {Rebuildable local image cache for configured Journal Monitor logos.} \\
\nolinkurl{data/Library} & {Visible main Library root and compatibility primary files.} \\
\nolinkurl{data/Spaces} & {Independent Space roots.} \\
\nolinkurl{data/My Publications} & {Managed local files for My Publications.} \\
\nolinkurl{data/Associated Files} & {Mirrored bundles for versions, supplements, and attachments.} \\
\nolinkurl{data/pdfs} & {Legacy managed PDF location retained for compatibility where needed.} \\
\texttt{reports} & {Generated reading notes and reports. These are artifacts, not the database.} \\
\nolinkurl{logs/paper_monitor_launcher.log} & {macOS launcher diagnostics when the local UI cannot start.} \\
\end{longtable}
\GewuTableEnd
}

Quick Look page thumbnails use a separate, automatically trimmed temporary cache. It is rebuildable preview data, not part of the Library, Associated Files, or the SQLite record, so it does not belong in a research backup.

Browser website data and its disposable favicon cache are managed by macOS outside the project. They are separate from SQLite bookmarks/history and from the monitor caches listed above. See \hyperref[backup-migrate]{Browser and credential backup boundaries} before moving to another Mac.

\section{Associated Files layout}\label{associated-files-layout}

An associated bundle follows the managed scope and path, then a work-specific bundle. Its children separate \texttt{Versions}, \texttt{Supplements}, and \texttt{Files}. The exact physical name can change as the paper moves; the Gewu Key remains stable.

\section{Database and exports}\label{database-and-exports}

SQLite is the only source of truth. JSON, HTML, Markdown, Word, LaTeX, and PDF outputs are exports, diagnostics, or reading artifacts. Editing one of them does not edit the Library record.

\begin{quote}
\textbf{Important}

Do not delete, rebuild, or ``clean'' the database because a generated JSON file looks stale. Diagnose the selected record and local path first.
\end{quote}

\begin{GewuRelated}
\textbf{See also:} \hyperref[storage-model]{How Gewu Insight Stores Files} · \hyperref[backup-migrate]{Back Up or Move to Another Mac} · \hyperref[privacy]{Local Data and Privacy}
\end{GewuRelated}

\Needspace{7\baselineskip}
\section*{Copyright and product notice}
\addcontentsline{toc}{section}{Copyright and product notice}
Gewu Insight is proprietary software. Copyright \textcopyright\ 2026 Lingyuan Kong and Zhengran Wu. All Rights Reserved.

This manual describes the product as shipped at the version printed on the title page. Online services, publisher access, and third-party metadata may change independently of the application.
\end{document}
